\documentclass[11pt]{article}

    \usepackage[breakable]{tcolorbox}
    \usepackage{parskip} % Stop auto-indenting (to mimic markdown behaviour)
    

    % Basic figure setup, for now with no caption control since it's done
    % automatically by Pandoc (which extracts ![](path) syntax from Markdown).
    \usepackage{graphicx}
    % Keep aspect ratio if custom image width or height is specified
    \setkeys{Gin}{keepaspectratio}
    % Maintain compatibility with old templates. Remove in nbconvert 6.0
    \let\Oldincludegraphics\includegraphics
    % Ensure that by default, figures have no caption (until we provide a
    % proper Figure object with a Caption API and a way to capture that
    % in the conversion process - todo).
    \usepackage{caption}
    \DeclareCaptionFormat{nocaption}{}
    \captionsetup{format=nocaption,aboveskip=0pt,belowskip=0pt}

    \usepackage{float}
    \floatplacement{figure}{H} % forces figures to be placed at the correct location
    \usepackage{xcolor} % Allow colors to be defined
    \usepackage{enumerate} % Needed for markdown enumerations to work
    \usepackage{geometry} % Used to adjust the document margins
    \usepackage{amsmath} % Equations
    \usepackage{amssymb} % Equations
    \usepackage{textcomp} % defines textquotesingle
    % Hack from http://tex.stackexchange.com/a/47451/13684:
    \AtBeginDocument{%
        \def\PYZsq{\textquotesingle}% Upright quotes in Pygmentized code
    }
    \usepackage{upquote} % Upright quotes for verbatim code
    \usepackage{eurosym} % defines \euro

    \usepackage{iftex}
    \ifPDFTeX
        \usepackage[T1]{fontenc}
        \IfFileExists{alphabeta.sty}{
              \usepackage{alphabeta}
          }{
              \usepackage[mathletters]{ucs}
              \usepackage[utf8x]{inputenc}
          }
    \else
        \usepackage{fontspec}
        \usepackage{unicode-math}
    \fi

    \usepackage{fancyvrb} % verbatim replacement that allows latex
    \usepackage{grffile} % extends the file name processing of package graphics
                         % to support a larger range
    \makeatletter % fix for old versions of grffile with XeLaTeX
    \@ifpackagelater{grffile}{2019/11/01}
    {
      % Do nothing on new versions
    }
    {
      \def\Gread@@xetex#1{%
        \IfFileExists{"\Gin@base".bb}%
        {\Gread@eps{\Gin@base.bb}}%
        {\Gread@@xetex@aux#1}%
      }
    }
    \makeatother
    \usepackage[Export]{adjustbox} % Used to constrain images to a maximum size
    \adjustboxset{max size={0.9\linewidth}{0.9\paperheight}}

    % The hyperref package gives us a pdf with properly built
    % internal navigation ('pdf bookmarks' for the table of contents,
    % internal cross-reference links, web links for URLs, etc.)
    \usepackage{hyperref}
    % The default LaTeX title has an obnoxious amount of whitespace. By default,
    % titling removes some of it. It also provides customization options.
    \usepackage{titling}
    \usepackage{longtable} % longtable support required by pandoc >1.10
    \usepackage{booktabs}  % table support for pandoc > 1.12.2
    \usepackage{array}     % table support for pandoc >= 2.11.3
    \usepackage{calc}      % table minipage width calculation for pandoc >= 2.11.1
    \usepackage[inline]{enumitem} % IRkernel/repr support (it uses the enumerate* environment)
    \usepackage[normalem]{ulem} % ulem is needed to support strikethroughs (\sout)
                                % normalem makes italics be italics, not underlines
    \usepackage{soul}      % strikethrough (\st) support for pandoc >= 3.0.0
    \usepackage{mathrsfs}
    

    
    % Colors for the hyperref package
    \definecolor{urlcolor}{rgb}{0,.145,.698}
    \definecolor{linkcolor}{rgb}{.71,0.21,0.01}
    \definecolor{citecolor}{rgb}{.12,.54,.11}

    % ANSI colors
    \definecolor{ansi-black}{HTML}{3E424D}
    \definecolor{ansi-black-intense}{HTML}{282C36}
    \definecolor{ansi-red}{HTML}{E75C58}
    \definecolor{ansi-red-intense}{HTML}{B22B31}
    \definecolor{ansi-green}{HTML}{00A250}
    \definecolor{ansi-green-intense}{HTML}{007427}
    \definecolor{ansi-yellow}{HTML}{DDB62B}
    \definecolor{ansi-yellow-intense}{HTML}{B27D12}
    \definecolor{ansi-blue}{HTML}{208FFB}
    \definecolor{ansi-blue-intense}{HTML}{0065CA}
    \definecolor{ansi-magenta}{HTML}{D160C4}
    \definecolor{ansi-magenta-intense}{HTML}{A03196}
    \definecolor{ansi-cyan}{HTML}{60C6C8}
    \definecolor{ansi-cyan-intense}{HTML}{258F8F}
    \definecolor{ansi-white}{HTML}{C5C1B4}
    \definecolor{ansi-white-intense}{HTML}{A1A6B2}
    \definecolor{ansi-default-inverse-fg}{HTML}{FFFFFF}
    \definecolor{ansi-default-inverse-bg}{HTML}{000000}

    % common color for the border for error outputs.
    \definecolor{outerrorbackground}{HTML}{FFDFDF}

    % commands and environments needed by pandoc snippets
    % extracted from the output of `pandoc -s`
    \providecommand{\tightlist}{%
      \setlength{\itemsep}{0pt}\setlength{\parskip}{0pt}}
    \DefineVerbatimEnvironment{Highlighting}{Verbatim}{commandchars=\\\{\}}
    % Add ',fontsize=\small' for more characters per line
    \newenvironment{Shaded}{}{}
    \newcommand{\KeywordTok}[1]{\textcolor[rgb]{0.00,0.44,0.13}{\textbf{{#1}}}}
    \newcommand{\DataTypeTok}[1]{\textcolor[rgb]{0.56,0.13,0.00}{{#1}}}
    \newcommand{\DecValTok}[1]{\textcolor[rgb]{0.25,0.63,0.44}{{#1}}}
    \newcommand{\BaseNTok}[1]{\textcolor[rgb]{0.25,0.63,0.44}{{#1}}}
    \newcommand{\FloatTok}[1]{\textcolor[rgb]{0.25,0.63,0.44}{{#1}}}
    \newcommand{\CharTok}[1]{\textcolor[rgb]{0.25,0.44,0.63}{{#1}}}
    \newcommand{\StringTok}[1]{\textcolor[rgb]{0.25,0.44,0.63}{{#1}}}
    \newcommand{\CommentTok}[1]{\textcolor[rgb]{0.38,0.63,0.69}{\textit{{#1}}}}
    \newcommand{\OtherTok}[1]{\textcolor[rgb]{0.00,0.44,0.13}{{#1}}}
    \newcommand{\AlertTok}[1]{\textcolor[rgb]{1.00,0.00,0.00}{\textbf{{#1}}}}
    \newcommand{\FunctionTok}[1]{\textcolor[rgb]{0.02,0.16,0.49}{{#1}}}
    \newcommand{\RegionMarkerTok}[1]{{#1}}
    \newcommand{\ErrorTok}[1]{\textcolor[rgb]{1.00,0.00,0.00}{\textbf{{#1}}}}
    \newcommand{\NormalTok}[1]{{#1}}

    % Additional commands for more recent versions of Pandoc
    \newcommand{\ConstantTok}[1]{\textcolor[rgb]{0.53,0.00,0.00}{{#1}}}
    \newcommand{\SpecialCharTok}[1]{\textcolor[rgb]{0.25,0.44,0.63}{{#1}}}
    \newcommand{\VerbatimStringTok}[1]{\textcolor[rgb]{0.25,0.44,0.63}{{#1}}}
    \newcommand{\SpecialStringTok}[1]{\textcolor[rgb]{0.73,0.40,0.53}{{#1}}}
    \newcommand{\ImportTok}[1]{{#1}}
    \newcommand{\DocumentationTok}[1]{\textcolor[rgb]{0.73,0.13,0.13}{\textit{{#1}}}}
    \newcommand{\AnnotationTok}[1]{\textcolor[rgb]{0.38,0.63,0.69}{\textbf{\textit{{#1}}}}}
    \newcommand{\CommentVarTok}[1]{\textcolor[rgb]{0.38,0.63,0.69}{\textbf{\textit{{#1}}}}}
    \newcommand{\VariableTok}[1]{\textcolor[rgb]{0.10,0.09,0.49}{{#1}}}
    \newcommand{\ControlFlowTok}[1]{\textcolor[rgb]{0.00,0.44,0.13}{\textbf{{#1}}}}
    \newcommand{\OperatorTok}[1]{\textcolor[rgb]{0.40,0.40,0.40}{{#1}}}
    \newcommand{\BuiltInTok}[1]{{#1}}
    \newcommand{\ExtensionTok}[1]{{#1}}
    \newcommand{\PreprocessorTok}[1]{\textcolor[rgb]{0.74,0.48,0.00}{{#1}}}
    \newcommand{\AttributeTok}[1]{\textcolor[rgb]{0.49,0.56,0.16}{{#1}}}
    \newcommand{\InformationTok}[1]{\textcolor[rgb]{0.38,0.63,0.69}{\textbf{\textit{{#1}}}}}
    \newcommand{\WarningTok}[1]{\textcolor[rgb]{0.38,0.63,0.69}{\textbf{\textit{{#1}}}}}
    \makeatletter
    \newsavebox\pandoc@box
    \newcommand*\pandocbounded[1]{%
      \sbox\pandoc@box{#1}%
      % scaling factors for width and height
      \Gscale@div\@tempa\textheight{\dimexpr\ht\pandoc@box+\dp\pandoc@box\relax}%
      \Gscale@div\@tempb\linewidth{\wd\pandoc@box}%
      % select the smaller of both
      \ifdim\@tempb\p@<\@tempa\p@
        \let\@tempa\@tempb
      \fi
      % scaling accordingly (\@tempa < 1)
      \ifdim\@tempa\p@<\p@
        \scalebox{\@tempa}{\usebox\pandoc@box}%
      % scaling not needed, use as it is
      \else
        \usebox{\pandoc@box}%
      \fi
    }
    \makeatother

    % Define a nice break command that doesn't care if a line doesn't already
    % exist.
    \def\br{\hspace*{\fill} \\* }
    % Math Jax compatibility definitions
    \def\gt{>}
    \def\lt{<}
    \let\Oldtex\TeX
    \let\Oldlatex\LaTeX
    \renewcommand{\TeX}{\textrm{\Oldtex}}
    \renewcommand{\LaTeX}{\textrm{\Oldlatex}}
    % Document parameters
    % Document title
    \title{Untitled1}
    
    
    
    
    
    
    
% Pygments definitions
\makeatletter
\def\PY@reset{\let\PY@it=\relax \let\PY@bf=\relax%
    \let\PY@ul=\relax \let\PY@tc=\relax%
    \let\PY@bc=\relax \let\PY@ff=\relax}
\def\PY@tok#1{\csname PY@tok@#1\endcsname}
\def\PY@toks#1+{\ifx\relax#1\empty\else%
    \PY@tok{#1}\expandafter\PY@toks\fi}
\def\PY@do#1{\PY@bc{\PY@tc{\PY@ul{%
    \PY@it{\PY@bf{\PY@ff{#1}}}}}}}
\def\PY#1#2{\PY@reset\PY@toks#1+\relax+\PY@do{#2}}

\@namedef{PY@tok@w}{\def\PY@tc##1{\textcolor[rgb]{0.73,0.73,0.73}{##1}}}
\@namedef{PY@tok@c}{\let\PY@it=\textit\def\PY@tc##1{\textcolor[rgb]{0.24,0.48,0.48}{##1}}}
\@namedef{PY@tok@cp}{\def\PY@tc##1{\textcolor[rgb]{0.61,0.40,0.00}{##1}}}
\@namedef{PY@tok@k}{\let\PY@bf=\textbf\def\PY@tc##1{\textcolor[rgb]{0.00,0.50,0.00}{##1}}}
\@namedef{PY@tok@kp}{\def\PY@tc##1{\textcolor[rgb]{0.00,0.50,0.00}{##1}}}
\@namedef{PY@tok@kt}{\def\PY@tc##1{\textcolor[rgb]{0.69,0.00,0.25}{##1}}}
\@namedef{PY@tok@o}{\def\PY@tc##1{\textcolor[rgb]{0.40,0.40,0.40}{##1}}}
\@namedef{PY@tok@ow}{\let\PY@bf=\textbf\def\PY@tc##1{\textcolor[rgb]{0.67,0.13,1.00}{##1}}}
\@namedef{PY@tok@nb}{\def\PY@tc##1{\textcolor[rgb]{0.00,0.50,0.00}{##1}}}
\@namedef{PY@tok@nf}{\def\PY@tc##1{\textcolor[rgb]{0.00,0.00,1.00}{##1}}}
\@namedef{PY@tok@nc}{\let\PY@bf=\textbf\def\PY@tc##1{\textcolor[rgb]{0.00,0.00,1.00}{##1}}}
\@namedef{PY@tok@nn}{\let\PY@bf=\textbf\def\PY@tc##1{\textcolor[rgb]{0.00,0.00,1.00}{##1}}}
\@namedef{PY@tok@ne}{\let\PY@bf=\textbf\def\PY@tc##1{\textcolor[rgb]{0.80,0.25,0.22}{##1}}}
\@namedef{PY@tok@nv}{\def\PY@tc##1{\textcolor[rgb]{0.10,0.09,0.49}{##1}}}
\@namedef{PY@tok@no}{\def\PY@tc##1{\textcolor[rgb]{0.53,0.00,0.00}{##1}}}
\@namedef{PY@tok@nl}{\def\PY@tc##1{\textcolor[rgb]{0.46,0.46,0.00}{##1}}}
\@namedef{PY@tok@ni}{\let\PY@bf=\textbf\def\PY@tc##1{\textcolor[rgb]{0.44,0.44,0.44}{##1}}}
\@namedef{PY@tok@na}{\def\PY@tc##1{\textcolor[rgb]{0.41,0.47,0.13}{##1}}}
\@namedef{PY@tok@nt}{\let\PY@bf=\textbf\def\PY@tc##1{\textcolor[rgb]{0.00,0.50,0.00}{##1}}}
\@namedef{PY@tok@nd}{\def\PY@tc##1{\textcolor[rgb]{0.67,0.13,1.00}{##1}}}
\@namedef{PY@tok@s}{\def\PY@tc##1{\textcolor[rgb]{0.73,0.13,0.13}{##1}}}
\@namedef{PY@tok@sd}{\let\PY@it=\textit\def\PY@tc##1{\textcolor[rgb]{0.73,0.13,0.13}{##1}}}
\@namedef{PY@tok@si}{\let\PY@bf=\textbf\def\PY@tc##1{\textcolor[rgb]{0.64,0.35,0.47}{##1}}}
\@namedef{PY@tok@se}{\let\PY@bf=\textbf\def\PY@tc##1{\textcolor[rgb]{0.67,0.36,0.12}{##1}}}
\@namedef{PY@tok@sr}{\def\PY@tc##1{\textcolor[rgb]{0.64,0.35,0.47}{##1}}}
\@namedef{PY@tok@ss}{\def\PY@tc##1{\textcolor[rgb]{0.10,0.09,0.49}{##1}}}
\@namedef{PY@tok@sx}{\def\PY@tc##1{\textcolor[rgb]{0.00,0.50,0.00}{##1}}}
\@namedef{PY@tok@m}{\def\PY@tc##1{\textcolor[rgb]{0.40,0.40,0.40}{##1}}}
\@namedef{PY@tok@gh}{\let\PY@bf=\textbf\def\PY@tc##1{\textcolor[rgb]{0.00,0.00,0.50}{##1}}}
\@namedef{PY@tok@gu}{\let\PY@bf=\textbf\def\PY@tc##1{\textcolor[rgb]{0.50,0.00,0.50}{##1}}}
\@namedef{PY@tok@gd}{\def\PY@tc##1{\textcolor[rgb]{0.63,0.00,0.00}{##1}}}
\@namedef{PY@tok@gi}{\def\PY@tc##1{\textcolor[rgb]{0.00,0.52,0.00}{##1}}}
\@namedef{PY@tok@gr}{\def\PY@tc##1{\textcolor[rgb]{0.89,0.00,0.00}{##1}}}
\@namedef{PY@tok@ge}{\let\PY@it=\textit}
\@namedef{PY@tok@gs}{\let\PY@bf=\textbf}
\@namedef{PY@tok@ges}{\let\PY@bf=\textbf\let\PY@it=\textit}
\@namedef{PY@tok@gp}{\let\PY@bf=\textbf\def\PY@tc##1{\textcolor[rgb]{0.00,0.00,0.50}{##1}}}
\@namedef{PY@tok@go}{\def\PY@tc##1{\textcolor[rgb]{0.44,0.44,0.44}{##1}}}
\@namedef{PY@tok@gt}{\def\PY@tc##1{\textcolor[rgb]{0.00,0.27,0.87}{##1}}}
\@namedef{PY@tok@err}{\def\PY@bc##1{{\setlength{\fboxsep}{\string -\fboxrule}\fcolorbox[rgb]{1.00,0.00,0.00}{1,1,1}{\strut ##1}}}}
\@namedef{PY@tok@kc}{\let\PY@bf=\textbf\def\PY@tc##1{\textcolor[rgb]{0.00,0.50,0.00}{##1}}}
\@namedef{PY@tok@kd}{\let\PY@bf=\textbf\def\PY@tc##1{\textcolor[rgb]{0.00,0.50,0.00}{##1}}}
\@namedef{PY@tok@kn}{\let\PY@bf=\textbf\def\PY@tc##1{\textcolor[rgb]{0.00,0.50,0.00}{##1}}}
\@namedef{PY@tok@kr}{\let\PY@bf=\textbf\def\PY@tc##1{\textcolor[rgb]{0.00,0.50,0.00}{##1}}}
\@namedef{PY@tok@bp}{\def\PY@tc##1{\textcolor[rgb]{0.00,0.50,0.00}{##1}}}
\@namedef{PY@tok@fm}{\def\PY@tc##1{\textcolor[rgb]{0.00,0.00,1.00}{##1}}}
\@namedef{PY@tok@vc}{\def\PY@tc##1{\textcolor[rgb]{0.10,0.09,0.49}{##1}}}
\@namedef{PY@tok@vg}{\def\PY@tc##1{\textcolor[rgb]{0.10,0.09,0.49}{##1}}}
\@namedef{PY@tok@vi}{\def\PY@tc##1{\textcolor[rgb]{0.10,0.09,0.49}{##1}}}
\@namedef{PY@tok@vm}{\def\PY@tc##1{\textcolor[rgb]{0.10,0.09,0.49}{##1}}}
\@namedef{PY@tok@sa}{\def\PY@tc##1{\textcolor[rgb]{0.73,0.13,0.13}{##1}}}
\@namedef{PY@tok@sb}{\def\PY@tc##1{\textcolor[rgb]{0.73,0.13,0.13}{##1}}}
\@namedef{PY@tok@sc}{\def\PY@tc##1{\textcolor[rgb]{0.73,0.13,0.13}{##1}}}
\@namedef{PY@tok@dl}{\def\PY@tc##1{\textcolor[rgb]{0.73,0.13,0.13}{##1}}}
\@namedef{PY@tok@s2}{\def\PY@tc##1{\textcolor[rgb]{0.73,0.13,0.13}{##1}}}
\@namedef{PY@tok@sh}{\def\PY@tc##1{\textcolor[rgb]{0.73,0.13,0.13}{##1}}}
\@namedef{PY@tok@s1}{\def\PY@tc##1{\textcolor[rgb]{0.73,0.13,0.13}{##1}}}
\@namedef{PY@tok@mb}{\def\PY@tc##1{\textcolor[rgb]{0.40,0.40,0.40}{##1}}}
\@namedef{PY@tok@mf}{\def\PY@tc##1{\textcolor[rgb]{0.40,0.40,0.40}{##1}}}
\@namedef{PY@tok@mh}{\def\PY@tc##1{\textcolor[rgb]{0.40,0.40,0.40}{##1}}}
\@namedef{PY@tok@mi}{\def\PY@tc##1{\textcolor[rgb]{0.40,0.40,0.40}{##1}}}
\@namedef{PY@tok@il}{\def\PY@tc##1{\textcolor[rgb]{0.40,0.40,0.40}{##1}}}
\@namedef{PY@tok@mo}{\def\PY@tc##1{\textcolor[rgb]{0.40,0.40,0.40}{##1}}}
\@namedef{PY@tok@ch}{\let\PY@it=\textit\def\PY@tc##1{\textcolor[rgb]{0.24,0.48,0.48}{##1}}}
\@namedef{PY@tok@cm}{\let\PY@it=\textit\def\PY@tc##1{\textcolor[rgb]{0.24,0.48,0.48}{##1}}}
\@namedef{PY@tok@cpf}{\let\PY@it=\textit\def\PY@tc##1{\textcolor[rgb]{0.24,0.48,0.48}{##1}}}
\@namedef{PY@tok@c1}{\let\PY@it=\textit\def\PY@tc##1{\textcolor[rgb]{0.24,0.48,0.48}{##1}}}
\@namedef{PY@tok@cs}{\let\PY@it=\textit\def\PY@tc##1{\textcolor[rgb]{0.24,0.48,0.48}{##1}}}

\def\PYZbs{\char`\\}
\def\PYZus{\char`\_}
\def\PYZob{\char`\{}
\def\PYZcb{\char`\}}
\def\PYZca{\char`\^}
\def\PYZam{\char`\&}
\def\PYZlt{\char`\<}
\def\PYZgt{\char`\>}
\def\PYZsh{\char`\#}
\def\PYZpc{\char`\%}
\def\PYZdl{\char`\$}
\def\PYZhy{\char`\-}
\def\PYZsq{\char`\'}
\def\PYZdq{\char`\"}
\def\PYZti{\char`\~}
% for compatibility with earlier versions
\def\PYZat{@}
\def\PYZlb{[}
\def\PYZrb{]}
\makeatother


    % For linebreaks inside Verbatim environment from package fancyvrb.
    \makeatletter
        \newbox\Wrappedcontinuationbox
        \newbox\Wrappedvisiblespacebox
        \newcommand*\Wrappedvisiblespace {\textcolor{red}{\textvisiblespace}}
        \newcommand*\Wrappedcontinuationsymbol {\textcolor{red}{\llap{\tiny$\m@th\hookrightarrow$}}}
        \newcommand*\Wrappedcontinuationindent {3ex }
        \newcommand*\Wrappedafterbreak {\kern\Wrappedcontinuationindent\copy\Wrappedcontinuationbox}
        % Take advantage of the already applied Pygments mark-up to insert
        % potential linebreaks for TeX processing.
        %        {, <, #, %, $, ' and ": go to next line.
        %        _, }, ^, &, >, - and ~: stay at end of broken line.
        % Use of \textquotesingle for straight quote.
        \newcommand*\Wrappedbreaksatspecials {%
            \def\PYGZus{\discretionary{\char`\_}{\Wrappedafterbreak}{\char`\_}}%
            \def\PYGZob{\discretionary{}{\Wrappedafterbreak\char`\{}{\char`\{}}%
            \def\PYGZcb{\discretionary{\char`\}}{\Wrappedafterbreak}{\char`\}}}%
            \def\PYGZca{\discretionary{\char`\^}{\Wrappedafterbreak}{\char`\^}}%
            \def\PYGZam{\discretionary{\char`\&}{\Wrappedafterbreak}{\char`\&}}%
            \def\PYGZlt{\discretionary{}{\Wrappedafterbreak\char`\<}{\char`\<}}%
            \def\PYGZgt{\discretionary{\char`\>}{\Wrappedafterbreak}{\char`\>}}%
            \def\PYGZsh{\discretionary{}{\Wrappedafterbreak\char`\#}{\char`\#}}%
            \def\PYGZpc{\discretionary{}{\Wrappedafterbreak\char`\%}{\char`\%}}%
            \def\PYGZdl{\discretionary{}{\Wrappedafterbreak\char`\$}{\char`\$}}%
            \def\PYGZhy{\discretionary{\char`\-}{\Wrappedafterbreak}{\char`\-}}%
            \def\PYGZsq{\discretionary{}{\Wrappedafterbreak\textquotesingle}{\textquotesingle}}%
            \def\PYGZdq{\discretionary{}{\Wrappedafterbreak\char`\"}{\char`\"}}%
            \def\PYGZti{\discretionary{\char`\~}{\Wrappedafterbreak}{\char`\~}}%
        }
        % Some characters . , ; ? ! / are not pygmentized.
        % This macro makes them "active" and they will insert potential linebreaks
        \newcommand*\Wrappedbreaksatpunct {%
            \lccode`\~`\.\lowercase{\def~}{\discretionary{\hbox{\char`\.}}{\Wrappedafterbreak}{\hbox{\char`\.}}}%
            \lccode`\~`\,\lowercase{\def~}{\discretionary{\hbox{\char`\,}}{\Wrappedafterbreak}{\hbox{\char`\,}}}%
            \lccode`\~`\;\lowercase{\def~}{\discretionary{\hbox{\char`\;}}{\Wrappedafterbreak}{\hbox{\char`\;}}}%
            \lccode`\~`\:\lowercase{\def~}{\discretionary{\hbox{\char`\:}}{\Wrappedafterbreak}{\hbox{\char`\:}}}%
            \lccode`\~`\?\lowercase{\def~}{\discretionary{\hbox{\char`\?}}{\Wrappedafterbreak}{\hbox{\char`\?}}}%
            \lccode`\~`\!\lowercase{\def~}{\discretionary{\hbox{\char`\!}}{\Wrappedafterbreak}{\hbox{\char`\!}}}%
            \lccode`\~`\/\lowercase{\def~}{\discretionary{\hbox{\char`\/}}{\Wrappedafterbreak}{\hbox{\char`\/}}}%
            \catcode`\.\active
            \catcode`\,\active
            \catcode`\;\active
            \catcode`\:\active
            \catcode`\?\active
            \catcode`\!\active
            \catcode`\/\active
            \lccode`\~`\~
        }
    \makeatother

    \let\OriginalVerbatim=\Verbatim
    \makeatletter
    \renewcommand{\Verbatim}[1][1]{%
        %\parskip\z@skip
        \sbox\Wrappedcontinuationbox {\Wrappedcontinuationsymbol}%
        \sbox\Wrappedvisiblespacebox {\FV@SetupFont\Wrappedvisiblespace}%
        \def\FancyVerbFormatLine ##1{\hsize\linewidth
            \vtop{\raggedright\hyphenpenalty\z@\exhyphenpenalty\z@
                \doublehyphendemerits\z@\finalhyphendemerits\z@
                \strut ##1\strut}%
        }%
        % If the linebreak is at a space, the latter will be displayed as visible
        % space at end of first line, and a continuation symbol starts next line.
        % Stretch/shrink are however usually zero for typewriter font.
        \def\FV@Space {%
            \nobreak\hskip\z@ plus\fontdimen3\font minus\fontdimen4\font
            \discretionary{\copy\Wrappedvisiblespacebox}{\Wrappedafterbreak}
            {\kern\fontdimen2\font}%
        }%

        % Allow breaks at special characters using \PYG... macros.
        \Wrappedbreaksatspecials
        % Breaks at punctuation characters . , ; ? ! and / need catcode=\active
        \OriginalVerbatim[#1,codes*=\Wrappedbreaksatpunct]%
    }
    \makeatother

    % Exact colors from NB
    \definecolor{incolor}{HTML}{303F9F}
    \definecolor{outcolor}{HTML}{D84315}
    \definecolor{cellborder}{HTML}{CFCFCF}
    \definecolor{cellbackground}{HTML}{F7F7F7}

    % prompt
    \makeatletter
    \newcommand{\boxspacing}{\kern\kvtcb@left@rule\kern\kvtcb@boxsep}
    \makeatother
    \newcommand{\prompt}[4]{
        {\ttfamily\llap{{\color{#2}[#3]:\hspace{3pt}#4}}\vspace{-\baselineskip}}
    }
    

    
    % Prevent overflowing lines due to hard-to-break entities
    \sloppy
    % Setup hyperref package
    \hypersetup{
      breaklinks=true,  % so long urls are correctly broken across lines
      colorlinks=true,
      urlcolor=urlcolor,
      linkcolor=linkcolor,
      citecolor=citecolor,
      }
    % Slightly bigger margins than the latex defaults
    
    \geometry{verbose,tmargin=1in,bmargin=1in,lmargin=1in,rmargin=1in}
    
    

\begin{document}
    
    \maketitle
    
    

    
    \begin{tcolorbox}[breakable, size=fbox, boxrule=1pt, pad at break*=1mm,colback=cellbackground, colframe=cellborder]
\prompt{In}{incolor}{ }{\boxspacing}
\begin{Verbatim}[commandchars=\\\{\}]
                                       \PY{c+c1}{\PYZsh{}PART\PYZhy{}1(Data Import \PYZam{} Inspection )}
\end{Verbatim}
\end{tcolorbox}

    \begin{tcolorbox}[breakable, size=fbox, boxrule=1pt, pad at break*=1mm,colback=cellbackground, colframe=cellborder]
\prompt{In}{incolor}{1}{\boxspacing}
\begin{Verbatim}[commandchars=\\\{\}]
\PY{k+kn}{import}\PY{+w}{ }\PY{n+nn}{pandas}\PY{+w}{ }\PY{k}{as}\PY{+w}{ }\PY{n+nn}{pd}
\end{Verbatim}
\end{tcolorbox}

    \begin{tcolorbox}[breakable, size=fbox, boxrule=1pt, pad at break*=1mm,colback=cellbackground, colframe=cellborder]
\prompt{In}{incolor}{2}{\boxspacing}
\begin{Verbatim}[commandchars=\\\{\}]
\PY{k+kn}{import}\PY{+w}{ }\PY{n+nn}{numpy}\PY{+w}{ }\PY{k}{as}\PY{+w}{ }\PY{n+nn}{np}
\end{Verbatim}
\end{tcolorbox}

    \begin{tcolorbox}[breakable, size=fbox, boxrule=1pt, pad at break*=1mm,colback=cellbackground, colframe=cellborder]
\prompt{In}{incolor}{3}{\boxspacing}
\begin{Verbatim}[commandchars=\\\{\}]
\PY{n}{pip} \PY{n}{install} \PY{n}{matplotlib}
\end{Verbatim}
\end{tcolorbox}

    \begin{Verbatim}[commandchars=\\\{\}]
Requirement already satisfied: matplotlib in
.\textbackslash{}AppData\textbackslash{}Local\textbackslash{}Programs\textbackslash{}Python\textbackslash{}Python314\textbackslash{}Lib\textbackslash{}site-packages (3.11.0)
Requirement already satisfied: contourpy>=1.0.1 in
.\textbackslash{}AppData\textbackslash{}Local\textbackslash{}Programs\textbackslash{}Python\textbackslash{}Python314\textbackslash{}Lib\textbackslash{}site-packages (from matplotlib)
(1.3.3)
Requirement already satisfied: cycler>=0.10 in
.\textbackslash{}AppData\textbackslash{}Local\textbackslash{}Programs\textbackslash{}Python\textbackslash{}Python314\textbackslash{}Lib\textbackslash{}site-packages (from matplotlib)
(0.12.1)
Requirement already satisfied: fonttools>=4.22.0 in
.\textbackslash{}AppData\textbackslash{}Local\textbackslash{}Programs\textbackslash{}Python\textbackslash{}Python314\textbackslash{}Lib\textbackslash{}site-packages (from matplotlib)
(4.63.0)
Requirement already satisfied: kiwisolver>=1.3.1 in
.\textbackslash{}AppData\textbackslash{}Local\textbackslash{}Programs\textbackslash{}Python\textbackslash{}Python314\textbackslash{}Lib\textbackslash{}site-packages (from matplotlib)
(1.5.0)
Requirement already satisfied: numpy>=1.25 in
.\textbackslash{}AppData\textbackslash{}Local\textbackslash{}Programs\textbackslash{}Python\textbackslash{}Python314\textbackslash{}Lib\textbackslash{}site-packages (from matplotlib)
(2.5.1)
Requirement already satisfied: packaging>=20.0 in
.\textbackslash{}AppData\textbackslash{}Local\textbackslash{}Programs\textbackslash{}Python\textbackslash{}Python314\textbackslash{}Lib\textbackslash{}site-packages (from matplotlib)
(26.2)
Requirement already satisfied: pillow>=9 in
.\textbackslash{}AppData\textbackslash{}Local\textbackslash{}Programs\textbackslash{}Python\textbackslash{}Python314\textbackslash{}Lib\textbackslash{}site-packages (from matplotlib)
(12.3.0)
Requirement already satisfied: pyparsing>=3 in
.\textbackslash{}AppData\textbackslash{}Local\textbackslash{}Programs\textbackslash{}Python\textbackslash{}Python314\textbackslash{}Lib\textbackslash{}site-packages (from matplotlib)
(3.3.2)
Requirement already satisfied: python-dateutil>=2.7 in
.\textbackslash{}AppData\textbackslash{}Local\textbackslash{}Programs\textbackslash{}Python\textbackslash{}Python314\textbackslash{}Lib\textbackslash{}site-packages (from matplotlib)
(2.9.0.post0)
Requirement already satisfied: six>=1.5 in
.\textbackslash{}AppData\textbackslash{}Local\textbackslash{}Programs\textbackslash{}Python\textbackslash{}Python314\textbackslash{}Lib\textbackslash{}site-packages (from python-
dateutil>=2.7->matplotlib) (1.17.0)
Note: you may need to restart the kernel to use updated packages.
    \end{Verbatim}

    \begin{tcolorbox}[breakable, size=fbox, boxrule=1pt, pad at break*=1mm,colback=cellbackground, colframe=cellborder]
\prompt{In}{incolor}{4}{\boxspacing}
\begin{Verbatim}[commandchars=\\\{\}]
\PY{k+kn}{import}\PY{+w}{ }\PY{n+nn}{matplotlib}\PY{n+nn}{.}\PY{n+nn}{pyplot}\PY{+w}{ }\PY{k}{as}\PY{+w}{ }\PY{n+nn}{plt}
\end{Verbatim}
\end{tcolorbox}

    \begin{tcolorbox}[breakable, size=fbox, boxrule=1pt, pad at break*=1mm,colback=cellbackground, colframe=cellborder]
\prompt{In}{incolor}{5}{\boxspacing}
\begin{Verbatim}[commandchars=\\\{\}]
\PY{n}{pip} \PY{n}{install} \PY{n}{seaborn}
\end{Verbatim}
\end{tcolorbox}

    \begin{Verbatim}[commandchars=\\\{\}]
Requirement already satisfied: seaborn in
.\textbackslash{}AppData\textbackslash{}Local\textbackslash{}Programs\textbackslash{}Python\textbackslash{}Python314\textbackslash{}Lib\textbackslash{}site-packages (0.13.2)
Requirement already satisfied: numpy!=1.24.0,>=1.20 in
.\textbackslash{}AppData\textbackslash{}Local\textbackslash{}Programs\textbackslash{}Python\textbackslash{}Python314\textbackslash{}Lib\textbackslash{}site-packages (from seaborn)
(2.5.1)
Requirement already satisfied: pandas>=1.2 in
.\textbackslash{}AppData\textbackslash{}Local\textbackslash{}Programs\textbackslash{}Python\textbackslash{}Python314\textbackslash{}Lib\textbackslash{}site-packages (from seaborn)
(3.0.3)
Requirement already satisfied: matplotlib!=3.6.1,>=3.4 in
.\textbackslash{}AppData\textbackslash{}Local\textbackslash{}Programs\textbackslash{}Python\textbackslash{}Python314\textbackslash{}Lib\textbackslash{}site-packages (from seaborn)
(3.11.0)
Requirement already satisfied: contourpy>=1.0.1 in
.\textbackslash{}AppData\textbackslash{}Local\textbackslash{}Programs\textbackslash{}Python\textbackslash{}Python314\textbackslash{}Lib\textbackslash{}site-packages (from
matplotlib!=3.6.1,>=3.4->seaborn) (1.3.3)
Requirement already satisfied: cycler>=0.10 in
.\textbackslash{}AppData\textbackslash{}Local\textbackslash{}Programs\textbackslash{}Python\textbackslash{}Python314\textbackslash{}Lib\textbackslash{}site-packages (from
matplotlib!=3.6.1,>=3.4->seaborn) (0.12.1)
Requirement already satisfied: fonttools>=4.22.0 in
.\textbackslash{}AppData\textbackslash{}Local\textbackslash{}Programs\textbackslash{}Python\textbackslash{}Python314\textbackslash{}Lib\textbackslash{}site-packages (from
matplotlib!=3.6.1,>=3.4->seaborn) (4.63.0)
Requirement already satisfied: kiwisolver>=1.3.1 in
.\textbackslash{}AppData\textbackslash{}Local\textbackslash{}Programs\textbackslash{}Python\textbackslash{}Python314\textbackslash{}Lib\textbackslash{}site-packages (from
matplotlib!=3.6.1,>=3.4->seaborn) (1.5.0)
Requirement already satisfied: packaging>=20.0 in
.\textbackslash{}AppData\textbackslash{}Local\textbackslash{}Programs\textbackslash{}Python\textbackslash{}Python314\textbackslash{}Lib\textbackslash{}site-packages (from
matplotlib!=3.6.1,>=3.4->seaborn) (26.2)
Requirement already satisfied: pillow>=9 in
.\textbackslash{}AppData\textbackslash{}Local\textbackslash{}Programs\textbackslash{}Python\textbackslash{}Python314\textbackslash{}Lib\textbackslash{}site-packages (from
matplotlib!=3.6.1,>=3.4->seaborn) (12.3.0)
Requirement already satisfied: pyparsing>=3 in
.\textbackslash{}AppData\textbackslash{}Local\textbackslash{}Programs\textbackslash{}Python\textbackslash{}Python314\textbackslash{}Lib\textbackslash{}site-packages (from
matplotlib!=3.6.1,>=3.4->seaborn) (3.3.2)
Requirement already satisfied: python-dateutil>=2.7 in
.\textbackslash{}AppData\textbackslash{}Local\textbackslash{}Programs\textbackslash{}Python\textbackslash{}Python314\textbackslash{}Lib\textbackslash{}site-packages (from
matplotlib!=3.6.1,>=3.4->seaborn) (2.9.0.post0)
Requirement already satisfied: tzdata in
.\textbackslash{}AppData\textbackslash{}Local\textbackslash{}Programs\textbackslash{}Python\textbackslash{}Python314\textbackslash{}Lib\textbackslash{}site-packages (from
pandas>=1.2->seaborn) (2026.3)
Requirement already satisfied: six>=1.5 in
.\textbackslash{}AppData\textbackslash{}Local\textbackslash{}Programs\textbackslash{}Python\textbackslash{}Python314\textbackslash{}Lib\textbackslash{}site-packages (from python-
dateutil>=2.7->matplotlib!=3.6.1,>=3.4->seaborn) (1.17.0)
Note: you may need to restart the kernel to use updated packages.
    \end{Verbatim}

    \begin{tcolorbox}[breakable, size=fbox, boxrule=1pt, pad at break*=1mm,colback=cellbackground, colframe=cellborder]
\prompt{In}{incolor}{6}{\boxspacing}
\begin{Verbatim}[commandchars=\\\{\}]
\PY{k+kn}{import}\PY{+w}{ }\PY{n+nn}{seaborn}\PY{+w}{ }\PY{k}{as}\PY{+w}{ }\PY{n+nn}{sns}
\end{Verbatim}
\end{tcolorbox}

    \begin{tcolorbox}[breakable, size=fbox, boxrule=1pt, pad at break*=1mm,colback=cellbackground, colframe=cellborder]
\prompt{In}{incolor}{7}{\boxspacing}
\begin{Verbatim}[commandchars=\\\{\}]
\PY{n}{pip} \PY{n}{install} \PY{n}{scikit}\PY{o}{\PYZhy{}}\PY{n}{learn}
\end{Verbatim}
\end{tcolorbox}

    \begin{Verbatim}[commandchars=\\\{\}]
Requirement already satisfied: scikit-learn in
.\textbackslash{}AppData\textbackslash{}Local\textbackslash{}Programs\textbackslash{}Python\textbackslash{}Python314\textbackslash{}Lib\textbackslash{}site-packages (1.9.0)
Requirement already satisfied: numpy>=1.24.1 in
.\textbackslash{}AppData\textbackslash{}Local\textbackslash{}Programs\textbackslash{}Python\textbackslash{}Python314\textbackslash{}Lib\textbackslash{}site-packages (from scikit-learn)
(2.5.1)
Requirement already satisfied: scipy>=1.10.0 in
.\textbackslash{}AppData\textbackslash{}Local\textbackslash{}Programs\textbackslash{}Python\textbackslash{}Python314\textbackslash{}Lib\textbackslash{}site-packages (from scikit-learn)
(1.18.0)
Requirement already satisfied: joblib>=1.4.0 in
.\textbackslash{}AppData\textbackslash{}Local\textbackslash{}Programs\textbackslash{}Python\textbackslash{}Python314\textbackslash{}Lib\textbackslash{}site-packages (from scikit-learn)
(1.5.3)
Requirement already satisfied: narwhals>=2.0.1 in
.\textbackslash{}AppData\textbackslash{}Local\textbackslash{}Programs\textbackslash{}Python\textbackslash{}Python314\textbackslash{}Lib\textbackslash{}site-packages (from scikit-learn)
(2.24.0)
Requirement already satisfied: threadpoolctl>=3.5.0 in
.\textbackslash{}AppData\textbackslash{}Local\textbackslash{}Programs\textbackslash{}Python\textbackslash{}Python314\textbackslash{}Lib\textbackslash{}site-packages (from scikit-learn)
(3.6.0)
Note: you may need to restart the kernel to use updated packages.
    \end{Verbatim}

    \begin{tcolorbox}[breakable, size=fbox, boxrule=1pt, pad at break*=1mm,colback=cellbackground, colframe=cellborder]
\prompt{In}{incolor}{8}{\boxspacing}
\begin{Verbatim}[commandchars=\\\{\}]
\PY{k+kn}{from}\PY{+w}{ }\PY{n+nn}{sklearn}\PY{+w}{ }\PY{k+kn}{import} \PY{n}{preprocessing}
\end{Verbatim}
\end{tcolorbox}

    \begin{tcolorbox}[breakable, size=fbox, boxrule=1pt, pad at break*=1mm,colback=cellbackground, colframe=cellborder]
\prompt{In}{incolor}{9}{\boxspacing}
\begin{Verbatim}[commandchars=\\\{\}]
\PY{l+s+sd}{\PYZsq{}\PYZsq{}\PYZsq{}The required libraries were successfully imported}

\PY{l+s+sd}{\PYZhy{} Pandas is used for data manipulation}
\PY{l+s+sd}{\PYZhy{} NumPy is used for numerical computations}
\PY{l+s+sd}{\PYZhy{} Matplotlib and Seaborn are used for data visualization}
\PY{l+s+sd}{\PYZhy{} Scikit\PYZhy{}learn provides machine learning algorithms and evaluation metrics\PYZsq{}\PYZsq{}\PYZsq{}}
\end{Verbatim}
\end{tcolorbox}

            \begin{tcolorbox}[breakable, size=fbox, boxrule=.5pt, pad at break*=1mm, opacityfill=0]
\prompt{Out}{outcolor}{9}{\boxspacing}
\begin{Verbatim}[commandchars=\\\{\}]
'The required libraries were successfully imported\textbackslash{}n\textbackslash{}n- Pandas is used for data
manipulation\textbackslash{}n- NumPy is used for numerical computations\textbackslash{}n- Matplotlib and
Seaborn are used for data visualization\textbackslash{}n- Scikit-learn provides machine
learning algorithms and evaluation metrics'
\end{Verbatim}
\end{tcolorbox}
        
    \begin{tcolorbox}[breakable, size=fbox, boxrule=1pt, pad at break*=1mm,colback=cellbackground, colframe=cellborder]
\prompt{In}{incolor}{10}{\boxspacing}
\begin{Verbatim}[commandchars=\\\{\}]
\PY{n}{df}\PY{o}{=}\PY{n}{pd}\PY{o}{.}\PY{n}{read\PYZus{}csv}\PY{p}{(}\PY{l+s+s2}{\PYZdq{}}\PY{l+s+s2}{housing.csv}\PY{l+s+s2}{\PYZdq{}}\PY{p}{)}\PY{c+c1}{\PYZsh{}Dataset loaded successfully!}
\end{Verbatim}
\end{tcolorbox}

    \begin{tcolorbox}[breakable, size=fbox, boxrule=1pt, pad at break*=1mm,colback=cellbackground, colframe=cellborder]
\prompt{In}{incolor}{11}{\boxspacing}
\begin{Verbatim}[commandchars=\\\{\}]
\PY{n+nb}{print}\PY{p}{(}\PY{n}{df}\PY{p}{)}
\end{Verbatim}
\end{tcolorbox}

    \begin{Verbatim}[commandchars=\\\{\}]
       longitude  latitude  housing\_median\_age  total\_rooms  total\_bedrooms  \textbackslash{}
0        -122.23     37.88                41.0        880.0           129.0
1        -122.22     37.86                21.0       7099.0          1106.0
2        -122.24     37.85                52.0       1467.0           190.0
3        -122.25     37.85                52.0       1274.0           235.0
4        -122.25     37.85                52.0       1627.0           280.0
{\ldots}          {\ldots}       {\ldots}                 {\ldots}          {\ldots}             {\ldots}
20635    -121.09     39.48                25.0       1665.0           374.0
20636    -121.21     39.49                18.0        697.0           150.0
20637    -121.22     39.43                17.0       2254.0           485.0
20638    -121.32     39.43                18.0       1860.0           409.0
20639    -121.24     39.37                16.0       2785.0           616.0

       population  households  median\_income  median\_house\_value  \textbackslash{}
0           322.0       126.0         8.3252            452600.0
1          2401.0      1138.0         8.3014            358500.0
2           496.0       177.0         7.2574            352100.0
3           558.0       219.0         5.6431            341300.0
4           565.0       259.0         3.8462            342200.0
{\ldots}           {\ldots}         {\ldots}            {\ldots}                 {\ldots}
20635       845.0       330.0         1.5603             78100.0
20636       356.0       114.0         2.5568             77100.0
20637      1007.0       433.0         1.7000             92300.0
20638       741.0       349.0         1.8672             84700.0
20639      1387.0       530.0         2.3886             89400.0

      ocean\_proximity
0            NEAR BAY
1            NEAR BAY
2            NEAR BAY
3            NEAR BAY
4            NEAR BAY
{\ldots}               {\ldots}
20635          INLAND
20636          INLAND
20637          INLAND
20638          INLAND
20639          INLAND

[20640 rows x 10 columns]
    \end{Verbatim}

    \begin{tcolorbox}[breakable, size=fbox, boxrule=1pt, pad at break*=1mm,colback=cellbackground, colframe=cellborder]
\prompt{In}{incolor}{12}{\boxspacing}
\begin{Verbatim}[commandchars=\\\{\}]
\PY{n}{df}\PY{o}{.}\PY{n}{head}\PY{p}{(}\PY{l+m+mi}{10}\PY{p}{)} \PY{c+c1}{\PYZsh{}Displays the first 10 rows of the dataset.}
\end{Verbatim}
\end{tcolorbox}

            \begin{tcolorbox}[breakable, size=fbox, boxrule=.5pt, pad at break*=1mm, opacityfill=0]
\prompt{Out}{outcolor}{12}{\boxspacing}
\begin{Verbatim}[commandchars=\\\{\}]
   longitude  latitude  housing\_median\_age  total\_rooms  total\_bedrooms  \textbackslash{}
0    -122.23     37.88                41.0        880.0           129.0
1    -122.22     37.86                21.0       7099.0          1106.0
2    -122.24     37.85                52.0       1467.0           190.0
3    -122.25     37.85                52.0       1274.0           235.0
4    -122.25     37.85                52.0       1627.0           280.0
5    -122.25     37.85                52.0        919.0           213.0
6    -122.25     37.84                52.0       2535.0           489.0
7    -122.25     37.84                52.0       3104.0           687.0
8    -122.26     37.84                42.0       2555.0           665.0
9    -122.25     37.84                52.0       3549.0           707.0

   population  households  median\_income  median\_house\_value ocean\_proximity
0       322.0       126.0         8.3252            452600.0        NEAR BAY
1      2401.0      1138.0         8.3014            358500.0        NEAR BAY
2       496.0       177.0         7.2574            352100.0        NEAR BAY
3       558.0       219.0         5.6431            341300.0        NEAR BAY
4       565.0       259.0         3.8462            342200.0        NEAR BAY
5       413.0       193.0         4.0368            269700.0        NEAR BAY
6      1094.0       514.0         3.6591            299200.0        NEAR BAY
7      1157.0       647.0         3.1200            241400.0        NEAR BAY
8      1206.0       595.0         2.0804            226700.0        NEAR BAY
9      1551.0       714.0         3.6912            261100.0        NEAR BAY
\end{Verbatim}
\end{tcolorbox}
        
    \begin{tcolorbox}[breakable, size=fbox, boxrule=1pt, pad at break*=1mm,colback=cellbackground, colframe=cellborder]
\prompt{In}{incolor}{13}{\boxspacing}
\begin{Verbatim}[commandchars=\\\{\}]
\PY{n}{df}\PY{o}{.}\PY{n}{tail}\PY{p}{(}\PY{l+m+mi}{5}\PY{p}{)} \PY{c+c1}{\PYZsh{}Displays the last 5 rows}
\end{Verbatim}
\end{tcolorbox}

            \begin{tcolorbox}[breakable, size=fbox, boxrule=.5pt, pad at break*=1mm, opacityfill=0]
\prompt{Out}{outcolor}{13}{\boxspacing}
\begin{Verbatim}[commandchars=\\\{\}]
       longitude  latitude  housing\_median\_age  total\_rooms  total\_bedrooms  \textbackslash{}
20635    -121.09     39.48                25.0       1665.0           374.0
20636    -121.21     39.49                18.0        697.0           150.0
20637    -121.22     39.43                17.0       2254.0           485.0
20638    -121.32     39.43                18.0       1860.0           409.0
20639    -121.24     39.37                16.0       2785.0           616.0

       population  households  median\_income  median\_house\_value  \textbackslash{}
20635       845.0       330.0         1.5603             78100.0
20636       356.0       114.0         2.5568             77100.0
20637      1007.0       433.0         1.7000             92300.0
20638       741.0       349.0         1.8672             84700.0
20639      1387.0       530.0         2.3886             89400.0

      ocean\_proximity
20635          INLAND
20636          INLAND
20637          INLAND
20638          INLAND
20639          INLAND
\end{Verbatim}
\end{tcolorbox}
        
    \begin{tcolorbox}[breakable, size=fbox, boxrule=1pt, pad at break*=1mm,colback=cellbackground, colframe=cellborder]
\prompt{In}{incolor}{14}{\boxspacing}
\begin{Verbatim}[commandchars=\\\{\}]
\PY{n}{df}\PY{o}{.}\PY{n}{head}\PY{p}{(}\PY{l+m+mi}{5}\PY{p}{)}
\end{Verbatim}
\end{tcolorbox}

            \begin{tcolorbox}[breakable, size=fbox, boxrule=.5pt, pad at break*=1mm, opacityfill=0]
\prompt{Out}{outcolor}{14}{\boxspacing}
\begin{Verbatim}[commandchars=\\\{\}]
   longitude  latitude  housing\_median\_age  total\_rooms  total\_bedrooms  \textbackslash{}
0    -122.23     37.88                41.0        880.0           129.0
1    -122.22     37.86                21.0       7099.0          1106.0
2    -122.24     37.85                52.0       1467.0           190.0
3    -122.25     37.85                52.0       1274.0           235.0
4    -122.25     37.85                52.0       1627.0           280.0

   population  households  median\_income  median\_house\_value ocean\_proximity
0       322.0       126.0         8.3252            452600.0        NEAR BAY
1      2401.0      1138.0         8.3014            358500.0        NEAR BAY
2       496.0       177.0         7.2574            352100.0        NEAR BAY
3       558.0       219.0         5.6431            341300.0        NEAR BAY
4       565.0       259.0         3.8462            342200.0        NEAR BAY
\end{Verbatim}
\end{tcolorbox}
        
    \begin{tcolorbox}[breakable, size=fbox, boxrule=1pt, pad at break*=1mm,colback=cellbackground, colframe=cellborder]
\prompt{In}{incolor}{15}{\boxspacing}
\begin{Verbatim}[commandchars=\\\{\}]
\PY{n}{df}\PY{o}{.}\PY{n}{head}\PY{p}{(}\PY{l+m+mi}{5}\PY{p}{)}
\end{Verbatim}
\end{tcolorbox}

            \begin{tcolorbox}[breakable, size=fbox, boxrule=.5pt, pad at break*=1mm, opacityfill=0]
\prompt{Out}{outcolor}{15}{\boxspacing}
\begin{Verbatim}[commandchars=\\\{\}]
   longitude  latitude  housing\_median\_age  total\_rooms  total\_bedrooms  \textbackslash{}
0    -122.23     37.88                41.0        880.0           129.0
1    -122.22     37.86                21.0       7099.0          1106.0
2    -122.24     37.85                52.0       1467.0           190.0
3    -122.25     37.85                52.0       1274.0           235.0
4    -122.25     37.85                52.0       1627.0           280.0

   population  households  median\_income  median\_house\_value ocean\_proximity
0       322.0       126.0         8.3252            452600.0        NEAR BAY
1      2401.0      1138.0         8.3014            358500.0        NEAR BAY
2       496.0       177.0         7.2574            352100.0        NEAR BAY
3       558.0       219.0         5.6431            341300.0        NEAR BAY
4       565.0       259.0         3.8462            342200.0        NEAR BAY
\end{Verbatim}
\end{tcolorbox}
        
    \begin{tcolorbox}[breakable, size=fbox, boxrule=1pt, pad at break*=1mm,colback=cellbackground, colframe=cellborder]
\prompt{In}{incolor}{16}{\boxspacing}
\begin{Verbatim}[commandchars=\\\{\}]
\PY{n}{df}\PY{o}{.}\PY{n}{dtypes}
\end{Verbatim}
\end{tcolorbox}

            \begin{tcolorbox}[breakable, size=fbox, boxrule=.5pt, pad at break*=1mm, opacityfill=0]
\prompt{Out}{outcolor}{16}{\boxspacing}
\begin{Verbatim}[commandchars=\\\{\}]
longitude             float64
latitude              float64
housing\_median\_age    float64
total\_rooms           float64
total\_bedrooms        float64
population            float64
households            float64
median\_income         float64
median\_house\_value    float64
ocean\_proximity           str
dtype: object
\end{Verbatim}
\end{tcolorbox}
        
    \begin{tcolorbox}[breakable, size=fbox, boxrule=1pt, pad at break*=1mm,colback=cellbackground, colframe=cellborder]
\prompt{In}{incolor}{17}{\boxspacing}
\begin{Verbatim}[commandchars=\\\{\}]
\PY{l+s+sd}{\PYZsq{}\PYZsq{}\PYZsq{}The dataset contains:}
\PY{l+s+sd}{\PYZhy{} Float variables}
\PY{l+s+sd}{\PYZhy{} Integer variables}
\PY{l+s+sd}{\PYZhy{} One categorical variable (object type)\PYZsq{}\PYZsq{}\PYZsq{}}
\end{Verbatim}
\end{tcolorbox}

            \begin{tcolorbox}[breakable, size=fbox, boxrule=.5pt, pad at break*=1mm, opacityfill=0]
\prompt{Out}{outcolor}{17}{\boxspacing}
\begin{Verbatim}[commandchars=\\\{\}]
'The dataset contains:\textbackslash{}n- Float variables\textbackslash{}n- Integer variables\textbackslash{}n- One
categorical variable (object type)'
\end{Verbatim}
\end{tcolorbox}
        
    \begin{tcolorbox}[breakable, size=fbox, boxrule=1pt, pad at break*=1mm,colback=cellbackground, colframe=cellborder]
\prompt{In}{incolor}{18}{\boxspacing}
\begin{Verbatim}[commandchars=\\\{\}]
\PY{n}{df}\PY{o}{.}\PY{n}{info}\PY{p}{(}\PY{p}{)}
\end{Verbatim}
\end{tcolorbox}

    \begin{Verbatim}[commandchars=\\\{\}]
<class 'pandas.DataFrame'>
RangeIndex: 20640 entries, 0 to 20639
Data columns (total 10 columns):
 \#   Column              Non-Null Count  Dtype
---  ------              --------------  -----
 0   longitude           20640 non-null  float64
 1   latitude            20640 non-null  float64
 2   housing\_median\_age  20640 non-null  float64
 3   total\_rooms         20640 non-null  float64
 4   total\_bedrooms      20433 non-null  float64
 5   population          20640 non-null  float64
 6   households          20640 non-null  float64
 7   median\_income       20640 non-null  float64
 8   median\_house\_value  20640 non-null  float64
 9   ocean\_proximity     20640 non-null  str
dtypes: float64(9), str(1)
memory usage: 1.7 MB
    \end{Verbatim}

    \begin{tcolorbox}[breakable, size=fbox, boxrule=1pt, pad at break*=1mm,colback=cellbackground, colframe=cellborder]
\prompt{In}{incolor}{19}{\boxspacing}
\begin{Verbatim}[commandchars=\\\{\}]
\PY{l+s+sd}{\PYZsq{}\PYZsq{}\PYZsq{}The dataset contains both numerical and categorical variables}
\PY{l+s+sd}{Most columns are numeric, while ocean\PYZus{}proximity is the only categorical feature\PYZsq{}\PYZsq{}\PYZsq{}}
\end{Verbatim}
\end{tcolorbox}

            \begin{tcolorbox}[breakable, size=fbox, boxrule=.5pt, pad at break*=1mm, opacityfill=0]
\prompt{Out}{outcolor}{19}{\boxspacing}
\begin{Verbatim}[commandchars=\\\{\}]
'The dataset contains both numerical and categorical variables\textbackslash{}nMost columns are
numeric, while ocean\_proximity is the only categorical feature'
\end{Verbatim}
\end{tcolorbox}
        
    \begin{tcolorbox}[breakable, size=fbox, boxrule=1pt, pad at break*=1mm,colback=cellbackground, colframe=cellborder]
\prompt{In}{incolor}{20}{\boxspacing}
\begin{Verbatim}[commandchars=\\\{\}]
\PY{n}{df}\PY{o}{.}\PY{n}{shape}
\end{Verbatim}
\end{tcolorbox}

            \begin{tcolorbox}[breakable, size=fbox, boxrule=.5pt, pad at break*=1mm, opacityfill=0]
\prompt{Out}{outcolor}{20}{\boxspacing}
\begin{Verbatim}[commandchars=\\\{\}]
(20640, 10)
\end{Verbatim}
\end{tcolorbox}
        
    \begin{tcolorbox}[breakable, size=fbox, boxrule=1pt, pad at break*=1mm,colback=cellbackground, colframe=cellborder]
\prompt{In}{incolor}{21}{\boxspacing}
\begin{Verbatim}[commandchars=\\\{\}]
\PY{l+s+sd}{\PYZsq{}\PYZsq{}\PYZsq{}The dataset contains 20,640 observations and 10 features\PYZsq{}\PYZsq{}\PYZsq{}}
\end{Verbatim}
\end{tcolorbox}

            \begin{tcolorbox}[breakable, size=fbox, boxrule=.5pt, pad at break*=1mm, opacityfill=0]
\prompt{Out}{outcolor}{21}{\boxspacing}
\begin{Verbatim}[commandchars=\\\{\}]
'The dataset contains 20,640 observations and 10 features'
\end{Verbatim}
\end{tcolorbox}
        
    \begin{tcolorbox}[breakable, size=fbox, boxrule=1pt, pad at break*=1mm,colback=cellbackground, colframe=cellborder]
\prompt{In}{incolor}{22}{\boxspacing}
\begin{Verbatim}[commandchars=\\\{\}]
\PY{n}{df}\PY{o}{.}\PY{n}{describe}\PY{p}{(}\PY{p}{)}     
\end{Verbatim}
\end{tcolorbox}

            \begin{tcolorbox}[breakable, size=fbox, boxrule=.5pt, pad at break*=1mm, opacityfill=0]
\prompt{Out}{outcolor}{22}{\boxspacing}
\begin{Verbatim}[commandchars=\\\{\}]
          longitude      latitude  housing\_median\_age   total\_rooms  \textbackslash{}
count  20640.000000  20640.000000        20640.000000  20640.000000
mean    -119.569704     35.631861           28.639486   2635.763081
std        2.003532      2.135952           12.585558   2181.615252
min     -124.350000     32.540000            1.000000      2.000000
25\%     -121.800000     33.930000           18.000000   1447.750000
50\%     -118.490000     34.260000           29.000000   2127.000000
75\%     -118.010000     37.710000           37.000000   3148.000000
max     -114.310000     41.950000           52.000000  39320.000000

       total\_bedrooms    population    households  median\_income  \textbackslash{}
count    20433.000000  20640.000000  20640.000000   20640.000000
mean       537.870553   1425.476744    499.539680       3.870671
std        421.385070   1132.462122    382.329753       1.899822
min          1.000000      3.000000      1.000000       0.499900
25\%        296.000000    787.000000    280.000000       2.563400
50\%        435.000000   1166.000000    409.000000       3.534800
75\%        647.000000   1725.000000    605.000000       4.743250
max       6445.000000  35682.000000   6082.000000      15.000100

       median\_house\_value
count        20640.000000
mean        206855.816909
std         115395.615874
min          14999.000000
25\%         119600.000000
50\%         179700.000000
75\%         264725.000000
max         500001.000000
\end{Verbatim}
\end{tcolorbox}
        
    \begin{tcolorbox}[breakable, size=fbox, boxrule=1pt, pad at break*=1mm,colback=cellbackground, colframe=cellborder]
\prompt{In}{incolor}{23}{\boxspacing}
\begin{Verbatim}[commandchars=\\\{\}]
\PY{l+s+sd}{\PYZsq{}\PYZsq{}\PYZsq{}}
\PY{l+s+sd}{The summary statistics provide an overview of the distribution of each numerical feature}

\PY{l+s+sd}{Some variables, such as total\PYZus{}rooms, population, and total\PYZus{}bedrooms, have large ranges and high standard deviations, indicating considerable variation }
\PY{l+s+sd}{These features may also contain outliers, which will be explored during the Exploratory Data Analysis (EDA) stage\PYZsq{}\PYZsq{}\PYZsq{}}
\end{Verbatim}
\end{tcolorbox}

            \begin{tcolorbox}[breakable, size=fbox, boxrule=.5pt, pad at break*=1mm, opacityfill=0]
\prompt{Out}{outcolor}{23}{\boxspacing}
\begin{Verbatim}[commandchars=\\\{\}]
'\textbackslash{}nThe summary statistics provide an overview of the distribution of each
numerical feature\textbackslash{}n\textbackslash{}nSome variables, such as total\_rooms, population, and
total\_bedrooms, have large ranges and high standard deviations, indicating
considerable variation \textbackslash{}nThese features may also contain outliers, which will be
explored during the Exploratory Data Analysis (EDA) stage'
\end{Verbatim}
\end{tcolorbox}
        
    \begin{tcolorbox}[breakable, size=fbox, boxrule=1pt, pad at break*=1mm,colback=cellbackground, colframe=cellborder]
\prompt{In}{incolor}{24}{\boxspacing}
\begin{Verbatim}[commandchars=\\\{\}]
\PY{l+s+sd}{\PYZsq{}\PYZsq{}\PYZsq{}\PYZsh{}\PYZsh{} Description of Features}
\PY{l+s+sd}{longitude \PYZhy{} Geographic longitude of the housing block}
\PY{l+s+sd}{latitude \PYZhy{} Geographic latitude of the housing block}
\PY{l+s+sd}{housing\PYZus{}median\PYZus{}age \PYZhy{} Median age of houses in the block (years)}
\PY{l+s+sd}{total\PYZus{}rooms \PYZhy{}Total number of rooms within the block}
\PY{l+s+sd}{total\PYZus{}bedrooms \PYZhy{} Total number of bedrooms within the block}
\PY{l+s+sd}{population \PYZhy{} Total population living in the block}
\PY{l+s+sd}{households \PYZhy{} Total number of households in the block}
\PY{l+s+sd}{median\PYZus{}income \PYZhy{}Median household income  }
\PY{l+s+sd}{median\PYZus{}house\PYZus{}value \PYZhy{}Median house value (Target Variable}
\PY{l+s+sd}{)}
\PY{l+s+sd}{ocean\PYZus{}proximity \PYZhy{} Distance category from the ocean (categorical feature)\PYZsq{}\PYZsq{}\PYZsq{}}
\end{Verbatim}
\end{tcolorbox}

            \begin{tcolorbox}[breakable, size=fbox, boxrule=.5pt, pad at break*=1mm, opacityfill=0]
\prompt{Out}{outcolor}{24}{\boxspacing}
\begin{Verbatim}[commandchars=\\\{\}]
'\#\# Description of Features\textbackslash{}nlongitude - Geographic longitude of the housing
block\textbackslash{}nlatitude - Geographic latitude of the housing block\textbackslash{}nhousing\_median\_age -
Median age of houses in the block (years)\textbackslash{}ntotal\_rooms -Total number of rooms
within the block\textbackslash{}ntotal\_bedrooms - Total number of bedrooms within the
block\textbackslash{}npopulation - Total population living in the block\textbackslash{}nhouseholds - Total
number of households in the block\textbackslash{}nmedian\_income -Median household income
\textbackslash{}nmedian\_house\_value -Median house value (Target Variable\textbackslash{}n)\textbackslash{}nocean\_proximity -
Distance category from the ocean (categorical feature)'
\end{Verbatim}
\end{tcolorbox}
        
    \begin{tcolorbox}[breakable, size=fbox, boxrule=1pt, pad at break*=1mm,colback=cellbackground, colframe=cellborder]
\prompt{In}{incolor}{25}{\boxspacing}
\begin{Verbatim}[commandchars=\\\{\}]
\PY{n}{categorical\PYZus{}cols} \PY{o}{=} \PY{n}{df}\PY{o}{.}\PY{n}{select\PYZus{}dtypes}\PY{p}{(}\PY{n}{include}\PY{o}{=}\PY{p}{[}\PY{l+s+s1}{\PYZsq{}}\PY{l+s+s1}{object}\PY{l+s+s1}{\PYZsq{}}\PY{p}{,} \PY{l+s+s1}{\PYZsq{}}\PY{l+s+s1}{str}\PY{l+s+s1}{\PYZsq{}}\PY{p}{]}\PY{p}{)}\PY{o}{.}\PY{n}{columns}\PY{o}{.}\PY{n}{tolist}\PY{p}{(}\PY{p}{)}
\end{Verbatim}
\end{tcolorbox}

    \begin{tcolorbox}[breakable, size=fbox, boxrule=1pt, pad at break*=1mm,colback=cellbackground, colframe=cellborder]
\prompt{In}{incolor}{26}{\boxspacing}
\begin{Verbatim}[commandchars=\\\{\}]
\PY{n}{numerical\PYZus{}cols} \PY{o}{=} \PY{n}{df}\PY{o}{.}\PY{n}{select\PYZus{}dtypes}\PY{p}{(}\PY{n}{include}\PY{o}{=}\PY{p}{[}\PY{l+s+s1}{\PYZsq{}}\PY{l+s+s1}{int64}\PY{l+s+s1}{\PYZsq{}}\PY{p}{,}\PY{l+s+s1}{\PYZsq{}}\PY{l+s+s1}{float64}\PY{l+s+s1}{\PYZsq{}}\PY{p}{]}\PY{p}{)}\PY{o}{.}\PY{n}{columns}\PY{o}{.}\PY{n}{tolist}\PY{p}{(}\PY{p}{)}
\end{Verbatim}
\end{tcolorbox}

    \begin{tcolorbox}[breakable, size=fbox, boxrule=1pt, pad at break*=1mm,colback=cellbackground, colframe=cellborder]
\prompt{In}{incolor}{27}{\boxspacing}
\begin{Verbatim}[commandchars=\\\{\}]
\PY{n+nb}{print}\PY{p}{(}\PY{l+s+s2}{\PYZdq{}}\PY{l+s+s2}{Numerical:}\PY{l+s+s2}{\PYZdq{}}\PY{p}{,} \PY{n}{numerical\PYZus{}cols}\PY{p}{)}
\PY{n+nb}{print}\PY{p}{(}\PY{l+s+s2}{\PYZdq{}}\PY{l+s+s2}{Categorical:}\PY{l+s+s2}{\PYZdq{}}\PY{p}{,} \PY{n}{categorical\PYZus{}cols}\PY{p}{)}
\end{Verbatim}
\end{tcolorbox}

    \begin{Verbatim}[commandchars=\\\{\}]
Numerical: ['longitude', 'latitude', 'housing\_median\_age', 'total\_rooms',
'total\_bedrooms', 'population', 'households', 'median\_income',
'median\_house\_value']
Categorical: ['ocean\_proximity']
    \end{Verbatim}

    \begin{tcolorbox}[breakable, size=fbox, boxrule=1pt, pad at break*=1mm,colback=cellbackground, colframe=cellborder]
\prompt{In}{incolor}{28}{\boxspacing}
\begin{Verbatim}[commandchars=\\\{\}]
\PY{n}{df}\PY{o}{.}\PY{n}{describe}\PY{p}{(}\PY{n}{include}\PY{o}{=}\PY{l+s+s1}{\PYZsq{}}\PY{l+s+s1}{str}\PY{l+s+s1}{\PYZsq{}}\PY{p}{)}          
\end{Verbatim}
\end{tcolorbox}

            \begin{tcolorbox}[breakable, size=fbox, boxrule=.5pt, pad at break*=1mm, opacityfill=0]
\prompt{Out}{outcolor}{28}{\boxspacing}
\begin{Verbatim}[commandchars=\\\{\}]
       ocean\_proximity
count            20640
unique               5
top          <1H OCEAN
freq              9136
\end{Verbatim}
\end{tcolorbox}
        
    \begin{tcolorbox}[breakable, size=fbox, boxrule=1pt, pad at break*=1mm,colback=cellbackground, colframe=cellborder]
\prompt{In}{incolor}{29}{\boxspacing}
\begin{Verbatim}[commandchars=\\\{\}]
                                              \PY{c+c1}{\PYZsh{}PART\PYZhy{}2(Data Cleaning )}
\end{Verbatim}
\end{tcolorbox}

    \begin{tcolorbox}[breakable, size=fbox, boxrule=1pt, pad at break*=1mm,colback=cellbackground, colframe=cellborder]
\prompt{In}{incolor}{30}{\boxspacing}
\begin{Verbatim}[commandchars=\\\{\}]
\PY{c+c1}{\PYZsh{}checking missing values}
\PY{n}{df}\PY{o}{.}\PY{n}{isnull}\PY{p}{(}\PY{p}{)}\PY{o}{.}\PY{n}{sum}\PY{p}{(}\PY{p}{)}
\end{Verbatim}
\end{tcolorbox}

            \begin{tcolorbox}[breakable, size=fbox, boxrule=.5pt, pad at break*=1mm, opacityfill=0]
\prompt{Out}{outcolor}{30}{\boxspacing}
\begin{Verbatim}[commandchars=\\\{\}]
longitude               0
latitude                0
housing\_median\_age      0
total\_rooms             0
total\_bedrooms        207
population              0
households              0
median\_income           0
median\_house\_value      0
ocean\_proximity         0
dtype: int64
\end{Verbatim}
\end{tcolorbox}
        
    \begin{tcolorbox}[breakable, size=fbox, boxrule=1pt, pad at break*=1mm,colback=cellbackground, colframe=cellborder]
\prompt{In}{incolor}{31}{\boxspacing}
\begin{Verbatim}[commandchars=\\\{\}]
\PY{l+s+sd}{\PYZsq{}\PYZsq{}\PYZsq{}The dataset contains missing values only in the otal\PYZus{}bedrooms column (207 missing values)}
\PY{l+s+sd}{All other columns are complete\PYZsq{}\PYZsq{}\PYZsq{}}
\end{Verbatim}
\end{tcolorbox}

            \begin{tcolorbox}[breakable, size=fbox, boxrule=.5pt, pad at break*=1mm, opacityfill=0]
\prompt{Out}{outcolor}{31}{\boxspacing}
\begin{Verbatim}[commandchars=\\\{\}]
'The dataset contains missing values only in the otal\_bedrooms column (207
missing values)\textbackslash{}nAll other columns are complete'
\end{Verbatim}
\end{tcolorbox}
        
    \begin{tcolorbox}[breakable, size=fbox, boxrule=1pt, pad at break*=1mm,colback=cellbackground, colframe=cellborder]
\prompt{In}{incolor}{32}{\boxspacing}
\begin{Verbatim}[commandchars=\\\{\}]
\PY{c+c1}{\PYZsh{}handling missing values ,\PYZsh{}Fill mising value with meidan}
\PY{n}{df}\PY{p}{[}\PY{l+s+s2}{\PYZdq{}}\PY{l+s+s2}{total\PYZus{}bedrooms}\PY{l+s+s2}{\PYZdq{}}\PY{p}{]} \PY{o}{=} \PY{n}{df}\PY{p}{[}\PY{l+s+s2}{\PYZdq{}}\PY{l+s+s2}{total\PYZus{}bedrooms}\PY{l+s+s2}{\PYZdq{}}\PY{p}{]}\PY{o}{.}\PY{n}{fillna}\PY{p}{(}\PY{n}{df}\PY{p}{[}\PY{l+s+s2}{\PYZdq{}}\PY{l+s+s2}{total\PYZus{}bedrooms}\PY{l+s+s2}{\PYZdq{}}\PY{p}{]}\PY{o}{.}\PY{n}{median}\PY{p}{(}\PY{p}{)}\PY{p}{)}
\end{Verbatim}
\end{tcolorbox}

    \begin{tcolorbox}[breakable, size=fbox, boxrule=1pt, pad at break*=1mm,colback=cellbackground, colframe=cellborder]
\prompt{In}{incolor}{33}{\boxspacing}
\begin{Verbatim}[commandchars=\\\{\}]
\PY{c+c1}{\PYZsh{}verify missing values}

\PY{n+nb}{print}\PY{p}{(}\PY{n}{df}\PY{o}{.}\PY{n}{isnull}\PY{p}{(}\PY{p}{)}\PY{o}{.}\PY{n}{sum}\PY{p}{(}\PY{p}{)}\PY{p}{)}
\end{Verbatim}
\end{tcolorbox}

    \begin{Verbatim}[commandchars=\\\{\}]
longitude             0
latitude              0
housing\_median\_age    0
total\_rooms           0
total\_bedrooms        0
population            0
households            0
median\_income         0
median\_house\_value    0
ocean\_proximity       0
dtype: int64
    \end{Verbatim}

    \begin{tcolorbox}[breakable, size=fbox, boxrule=1pt, pad at break*=1mm,colback=cellbackground, colframe=cellborder]
\prompt{In}{incolor}{34}{\boxspacing}
\begin{Verbatim}[commandchars=\\\{\}]
\PY{l+s+sd}{\PYZsq{}\PYZsq{}\PYZsq{}The median was chosen because total\PYZus{}bedrooms contains extreme values and possible outliers}

\PY{l+s+sd}{Unlike the mean, the median is not affected by unusually large or small values, making it a more reliable replacement for missing values\PYZsq{}\PYZsq{}\PYZsq{}}
\end{Verbatim}
\end{tcolorbox}

            \begin{tcolorbox}[breakable, size=fbox, boxrule=.5pt, pad at break*=1mm, opacityfill=0]
\prompt{Out}{outcolor}{34}{\boxspacing}
\begin{Verbatim}[commandchars=\\\{\}]
'The median was chosen because total\_bedrooms contains extreme values and
possible outliers\textbackslash{}n\textbackslash{}nUnlike the mean, the median is not affected by unusually
large or small values, making it a more reliable replacement for missing values'
\end{Verbatim}
\end{tcolorbox}
        
    \begin{tcolorbox}[breakable, size=fbox, boxrule=1pt, pad at break*=1mm,colback=cellbackground, colframe=cellborder]
\prompt{In}{incolor}{35}{\boxspacing}
\begin{Verbatim}[commandchars=\\\{\}]
\PY{c+c1}{\PYZsh{}checking duplicates}
\PY{n}{duplicates} \PY{o}{=} \PY{n}{df}\PY{o}{.}\PY{n}{duplicated}\PY{p}{(}\PY{p}{)}\PY{o}{.}\PY{n}{sum}\PY{p}{(}\PY{p}{)}

\PY{n+nb}{print}\PY{p}{(}\PY{l+s+s2}{\PYZdq{}}\PY{l+s+s2}{Duplicate Rows:}\PY{l+s+s2}{\PYZdq{}}\PY{p}{,} \PY{n}{duplicates}\PY{p}{)}
\end{Verbatim}
\end{tcolorbox}

    \begin{Verbatim}[commandchars=\\\{\}]
Duplicate Rows: 0
    \end{Verbatim}

    \begin{tcolorbox}[breakable, size=fbox, boxrule=1pt, pad at break*=1mm,colback=cellbackground, colframe=cellborder]
\prompt{In}{incolor}{36}{\boxspacing}
\begin{Verbatim}[commandchars=\\\{\}]
\PY{l+s+sd}{\PYZsq{}\PYZsq{}\PYZsq{}No duplicate rows were found in the dataset\PYZsq{}\PYZsq{}\PYZsq{}}
\end{Verbatim}
\end{tcolorbox}

            \begin{tcolorbox}[breakable, size=fbox, boxrule=.5pt, pad at break*=1mm, opacityfill=0]
\prompt{Out}{outcolor}{36}{\boxspacing}
\begin{Verbatim}[commandchars=\\\{\}]
'No duplicate rows were found in the dataset'
\end{Verbatim}
\end{tcolorbox}
        
    \begin{tcolorbox}[breakable, size=fbox, boxrule=1pt, pad at break*=1mm,colback=cellbackground, colframe=cellborder]
\prompt{In}{incolor}{37}{\boxspacing}
\begin{Verbatim}[commandchars=\\\{\}]
\PY{c+c1}{\PYZsh{}removing duplicates}
\PY{n}{df}\PY{o}{.}\PY{n}{drop\PYZus{}duplicates}\PY{p}{(}\PY{n}{inplace}\PY{o}{=}\PY{k+kc}{True}\PY{p}{)}
\PY{n+nb}{print}\PY{p}{(}\PY{l+s+s2}{\PYZdq{}}\PY{l+s+s2}{Duplicates after removal:}\PY{l+s+s2}{\PYZdq{}}\PY{p}{,} \PY{n}{df}\PY{o}{.}\PY{n}{duplicated}\PY{p}{(}\PY{p}{)}\PY{o}{.}\PY{n}{sum}\PY{p}{(}\PY{p}{)}\PY{p}{)}
\end{Verbatim}
\end{tcolorbox}

    \begin{Verbatim}[commandchars=\\\{\}]
Duplicates after removal: 0
    \end{Verbatim}

    \begin{tcolorbox}[breakable, size=fbox, boxrule=1pt, pad at break*=1mm,colback=cellbackground, colframe=cellborder]
\prompt{In}{incolor}{38}{\boxspacing}
\begin{Verbatim}[commandchars=\\\{\}]
\PY{l+s+sd}{\PYZsq{}\PYZsq{}\PYZsq{}No duplicate rows were found in the dataset\PYZsq{}\PYZsq{}\PYZsq{}}
\end{Verbatim}
\end{tcolorbox}

            \begin{tcolorbox}[breakable, size=fbox, boxrule=.5pt, pad at break*=1mm, opacityfill=0]
\prompt{Out}{outcolor}{38}{\boxspacing}
\begin{Verbatim}[commandchars=\\\{\}]
'No duplicate rows were found in the dataset'
\end{Verbatim}
\end{tcolorbox}
        
    \begin{tcolorbox}[breakable, size=fbox, boxrule=1pt, pad at break*=1mm,colback=cellbackground, colframe=cellborder]
\prompt{In}{incolor}{39}{\boxspacing}
\begin{Verbatim}[commandchars=\\\{\}]
\PY{c+c1}{\PYZsh{}check incorrecccct values}
\PY{n+nb}{print}\PY{p}{(}\PY{l+s+s2}{\PYZdq{}}\PY{l+s+s2}{Housing Age \PYZlt{}=0 :}\PY{l+s+s2}{\PYZdq{}}\PY{p}{,} \PY{p}{(}\PY{n}{df}\PY{p}{[}\PY{l+s+s2}{\PYZdq{}}\PY{l+s+s2}{housing\PYZus{}median\PYZus{}age}\PY{l+s+s2}{\PYZdq{}}\PY{p}{]} \PY{o}{\PYZlt{}}\PY{o}{=} \PY{l+m+mi}{0}\PY{p}{)}\PY{o}{.}\PY{n}{sum}\PY{p}{(}\PY{p}{)}\PY{p}{)}

\PY{n+nb}{print}\PY{p}{(}\PY{l+s+s2}{\PYZdq{}}\PY{l+s+s2}{Median Income \PYZlt{}=0 :}\PY{l+s+s2}{\PYZdq{}}\PY{p}{,} \PY{p}{(}\PY{n}{df}\PY{p}{[}\PY{l+s+s2}{\PYZdq{}}\PY{l+s+s2}{median\PYZus{}income}\PY{l+s+s2}{\PYZdq{}}\PY{p}{]} \PY{o}{\PYZlt{}}\PY{o}{=} \PY{l+m+mi}{0}\PY{p}{)}\PY{o}{.}\PY{n}{sum}\PY{p}{(}\PY{p}{)}\PY{p}{)}

\PY{n+nb}{print}\PY{p}{(}\PY{l+s+s2}{\PYZdq{}}\PY{l+s+s2}{House Price \PYZlt{}=0 :}\PY{l+s+s2}{\PYZdq{}}\PY{p}{,} \PY{p}{(}\PY{n}{df}\PY{p}{[}\PY{l+s+s2}{\PYZdq{}}\PY{l+s+s2}{median\PYZus{}house\PYZus{}value}\PY{l+s+s2}{\PYZdq{}}\PY{p}{]} \PY{o}{\PYZlt{}}\PY{o}{=} \PY{l+m+mi}{0}\PY{p}{)}\PY{o}{.}\PY{n}{sum}\PY{p}{(}\PY{p}{)}\PY{p}{)}

\PY{n+nb}{print}\PY{p}{(}\PY{l+s+s2}{\PYZdq{}}\PY{l+s+s2}{Total Rooms \PYZlt{}=0 :}\PY{l+s+s2}{\PYZdq{}}\PY{p}{,} \PY{p}{(}\PY{n}{df}\PY{p}{[}\PY{l+s+s2}{\PYZdq{}}\PY{l+s+s2}{total\PYZus{}rooms}\PY{l+s+s2}{\PYZdq{}}\PY{p}{]} \PY{o}{\PYZlt{}}\PY{o}{=} \PY{l+m+mi}{0}\PY{p}{)}\PY{o}{.}\PY{n}{sum}\PY{p}{(}\PY{p}{)}\PY{p}{)}

\PY{n+nb}{print}\PY{p}{(}\PY{l+s+s2}{\PYZdq{}}\PY{l+s+s2}{Population \PYZlt{}=0 :}\PY{l+s+s2}{\PYZdq{}}\PY{p}{,} \PY{p}{(}\PY{n}{df}\PY{p}{[}\PY{l+s+s2}{\PYZdq{}}\PY{l+s+s2}{population}\PY{l+s+s2}{\PYZdq{}}\PY{p}{]} \PY{o}{\PYZlt{}}\PY{o}{=} \PY{l+m+mi}{0}\PY{p}{)}\PY{o}{.}\PY{n}{sum}\PY{p}{(}\PY{p}{)}\PY{p}{)}

\PY{n+nb}{print}\PY{p}{(}\PY{l+s+s2}{\PYZdq{}}\PY{l+s+s2}{Households \PYZlt{}=0 :}\PY{l+s+s2}{\PYZdq{}}\PY{p}{,} \PY{p}{(}\PY{n}{df}\PY{p}{[}\PY{l+s+s2}{\PYZdq{}}\PY{l+s+s2}{households}\PY{l+s+s2}{\PYZdq{}}\PY{p}{]} \PY{o}{\PYZlt{}}\PY{o}{=} \PY{l+m+mi}{0}\PY{p}{)}\PY{o}{.}\PY{n}{sum}\PY{p}{(}\PY{p}{)}\PY{p}{)}
\end{Verbatim}
\end{tcolorbox}

    \begin{Verbatim}[commandchars=\\\{\}]
Housing Age <=0 : 0
Median Income <=0 : 0
House Price <=0 : 0
Total Rooms <=0 : 0
Population <=0 : 0
Households <=0 : 0
    \end{Verbatim}

    \begin{tcolorbox}[breakable, size=fbox, boxrule=1pt, pad at break*=1mm,colback=cellbackground, colframe=cellborder]
\prompt{In}{incolor}{40}{\boxspacing}
\begin{Verbatim}[commandchars=\\\{\}]
\PY{l+s+sd}{\PYZsq{}\PYZsq{}\PYZsq{}No impossible values were detected.}

\PY{l+s+sd}{All numerical features contain valid positive values suitable for analysis\PYZsq{}\PYZsq{}\PYZsq{}}
\end{Verbatim}
\end{tcolorbox}

            \begin{tcolorbox}[breakable, size=fbox, boxrule=.5pt, pad at break*=1mm, opacityfill=0]
\prompt{Out}{outcolor}{40}{\boxspacing}
\begin{Verbatim}[commandchars=\\\{\}]
'No impossible values were detected.\textbackslash{}n\textbackslash{}nAll numerical features contain valid
positive values suitable for analysis'
\end{Verbatim}
\end{tcolorbox}
        
    \begin{tcolorbox}[breakable, size=fbox, boxrule=1pt, pad at break*=1mm,colback=cellbackground, colframe=cellborder]
\prompt{In}{incolor}{41}{\boxspacing}
\begin{Verbatim}[commandchars=\\\{\}]
\PY{c+c1}{\PYZsh{}checking negative values}
\PY{n}{numerical\PYZus{}cols} \PY{o}{=} \PY{n}{df}\PY{o}{.}\PY{n}{select\PYZus{}dtypes}\PY{p}{(}\PY{n}{include}\PY{o}{=}\PY{p}{[}\PY{l+s+s1}{\PYZsq{}}\PY{l+s+s1}{int64}\PY{l+s+s1}{\PYZsq{}}\PY{p}{,}\PY{l+s+s1}{\PYZsq{}}\PY{l+s+s1}{float64}\PY{l+s+s1}{\PYZsq{}}\PY{p}{]}\PY{p}{)}\PY{o}{.}\PY{n}{columns}

\PY{k}{for} \PY{n}{col} \PY{o+ow}{in} \PY{n}{numerical\PYZus{}cols}\PY{p}{:}
    \PY{n+nb}{print}\PY{p}{(}\PY{n}{col}\PY{p}{,} \PY{p}{(}\PY{n}{df}\PY{p}{[}\PY{n}{col}\PY{p}{]} \PY{o}{\PYZlt{}} \PY{l+m+mi}{0}\PY{p}{)}\PY{o}{.}\PY{n}{sum}\PY{p}{(}\PY{p}{)}\PY{p}{)}
\end{Verbatim}
\end{tcolorbox}

    \begin{Verbatim}[commandchars=\\\{\}]
longitude 20640
latitude 0
housing\_median\_age 0
total\_rooms 0
total\_bedrooms 0
population 0
households 0
median\_income 0
median\_house\_value 0
    \end{Verbatim}

    \begin{tcolorbox}[breakable, size=fbox, boxrule=1pt, pad at break*=1mm,colback=cellbackground, colframe=cellborder]
\prompt{In}{incolor}{42}{\boxspacing}
\begin{Verbatim}[commandchars=\\\{\}]
\PY{k}{for} \PY{n}{col} \PY{o+ow}{in} \PY{n}{numerical\PYZus{}cols}\PY{p}{:}
    \PY{n+nb}{print}\PY{p}{(}\PY{n}{col}\PY{p}{,} \PY{p}{(}\PY{n}{df}\PY{p}{[}\PY{n}{col}\PY{p}{]} \PY{o}{==} \PY{l+m+mi}{0}\PY{p}{)}\PY{o}{.}\PY{n}{sum}\PY{p}{(}\PY{p}{)}\PY{p}{)}
\end{Verbatim}
\end{tcolorbox}

    \begin{Verbatim}[commandchars=\\\{\}]
longitude 0
latitude 0
housing\_median\_age 0
total\_rooms 0
total\_bedrooms 0
population 0
households 0
median\_income 0
median\_house\_value 0
    \end{Verbatim}

    \begin{tcolorbox}[breakable, size=fbox, boxrule=1pt, pad at break*=1mm,colback=cellbackground, colframe=cellborder]
\prompt{In}{incolor}{43}{\boxspacing}
\begin{Verbatim}[commandchars=\\\{\}]
\PY{n}{numerical\PYZus{}cols} \PY{o}{=} \PY{n}{df}\PY{o}{.}\PY{n}{select\PYZus{}dtypes}\PY{p}{(}\PY{n}{include}\PY{o}{=}\PY{p}{[}\PY{l+s+s1}{\PYZsq{}}\PY{l+s+s1}{int64}\PY{l+s+s1}{\PYZsq{}}\PY{p}{,} \PY{l+s+s1}{\PYZsq{}}\PY{l+s+s1}{float64}\PY{l+s+s1}{\PYZsq{}}\PY{p}{]}\PY{p}{)}\PY{o}{.}\PY{n}{columns}
\PY{n+nb}{print}\PY{p}{(}\PY{n}{numerical\PYZus{}cols}\PY{p}{)}
\end{Verbatim}
\end{tcolorbox}

    \begin{Verbatim}[commandchars=\\\{\}]
Index(['longitude', 'latitude', 'housing\_median\_age', 'total\_rooms',
       'total\_bedrooms', 'population', 'households', 'median\_income',
       'median\_house\_value'],
      dtype='str')
    \end{Verbatim}

    \begin{tcolorbox}[breakable, size=fbox, boxrule=1pt, pad at break*=1mm,colback=cellbackground, colframe=cellborder]
\prompt{In}{incolor}{44}{\boxspacing}
\begin{Verbatim}[commandchars=\\\{\}]
\PY{n+nb}{print}\PY{p}{(}\PY{n}{df}\PY{o}{.}\PY{n}{dtypes}\PY{p}{)}
\end{Verbatim}
\end{tcolorbox}

    \begin{Verbatim}[commandchars=\\\{\}]
longitude             float64
latitude              float64
housing\_median\_age    float64
total\_rooms           float64
total\_bedrooms        float64
population            float64
households            float64
median\_income         float64
median\_house\_value    float64
ocean\_proximity           str
dtype: object
    \end{Verbatim}

    \begin{tcolorbox}[breakable, size=fbox, boxrule=1pt, pad at break*=1mm,colback=cellbackground, colframe=cellborder]
\prompt{In}{incolor}{45}{\boxspacing}
\begin{Verbatim}[commandchars=\\\{\}]
\PY{l+s+sd}{\PYZsq{}\PYZsq{}\PYZsq{}All columns have appropriate data types}

\PY{l+s+sd}{No conversion was required because:}

\PY{l+s+sd}{\PYZhy{} Numerical variables are stored as integers or floating\PYZhy{}point numbers.}
\PY{l+s+sd}{\PYZhy{} ocean\PYZus{}proximity is correctly stored as a categorical (object) variabl\PYZsq{}\PYZsq{}\PYZsq{}}
\end{Verbatim}
\end{tcolorbox}

            \begin{tcolorbox}[breakable, size=fbox, boxrule=.5pt, pad at break*=1mm, opacityfill=0]
\prompt{Out}{outcolor}{45}{\boxspacing}
\begin{Verbatim}[commandchars=\\\{\}]
'All columns have appropriate data types\textbackslash{}n\textbackslash{}nNo conversion was required
because:\textbackslash{}n\textbackslash{}n- Numerical variables are stored as integers or floating-point
numbers.\textbackslash{}n- ocean\_proximity is correctly stored as a categorical (object)
variabl'
\end{Verbatim}
\end{tcolorbox}
        
    \begin{tcolorbox}[breakable, size=fbox, boxrule=1pt, pad at break*=1mm,colback=cellbackground, colframe=cellborder]
\prompt{In}{incolor}{46}{\boxspacing}
\begin{Verbatim}[commandchars=\\\{\}]
\PY{n+nb}{print}\PY{p}{(}\PY{l+s+s2}{\PYZdq{}}\PY{l+s+s2}{Final Dataset Shape:}\PY{l+s+s2}{\PYZdq{}}\PY{p}{,} \PY{n}{df}\PY{o}{.}\PY{n}{shape}\PY{p}{)}
\end{Verbatim}
\end{tcolorbox}

    \begin{Verbatim}[commandchars=\\\{\}]
Final Dataset Shape: (20640, 10)
    \end{Verbatim}

    \begin{tcolorbox}[breakable, size=fbox, boxrule=1pt, pad at break*=1mm,colback=cellbackground, colframe=cellborder]
\prompt{In}{incolor}{47}{\boxspacing}
\begin{Verbatim}[commandchars=\\\{\}]
\PY{n+nb}{print}\PY{p}{(}\PY{n}{df}\PY{o}{.}\PY{n}{isnull}\PY{p}{(}\PY{p}{)}\PY{o}{.}\PY{n}{sum}\PY{p}{(}\PY{p}{)}\PY{p}{)}
\end{Verbatim}
\end{tcolorbox}

    \begin{Verbatim}[commandchars=\\\{\}]
longitude             0
latitude              0
housing\_median\_age    0
total\_rooms           0
total\_bedrooms        0
population            0
households            0
median\_income         0
median\_house\_value    0
ocean\_proximity       0
dtype: int64
    \end{Verbatim}

    \begin{tcolorbox}[breakable, size=fbox, boxrule=1pt, pad at break*=1mm,colback=cellbackground, colframe=cellborder]
\prompt{In}{incolor}{48}{\boxspacing}
\begin{Verbatim}[commandchars=\\\{\}]
\PY{c+c1}{\PYZsh{}\PYZsh{}\PYZsh{} Interpretation }

\PY{l+s+sd}{\PYZsq{}\PYZsq{}\PYZsq{}The dataset is now completely clean}

\PY{l+s+sd}{\PYZhy{} No missing values remain}
\PY{l+s+sd}{\PYZhy{} No duplicate records exist}
\PY{l+s+sd}{\PYZhy{} No invalid values were detected}
\PY{l+s+sd}{\PYZhy{} All variables have appropriate data types}

\PY{l+s+sd}{The dataset is now ready for Exploratory Data Analysis (EDA) and machine learning\PYZsq{}\PYZsq{}\PYZsq{}}
\end{Verbatim}
\end{tcolorbox}

            \begin{tcolorbox}[breakable, size=fbox, boxrule=.5pt, pad at break*=1mm, opacityfill=0]
\prompt{Out}{outcolor}{48}{\boxspacing}
\begin{Verbatim}[commandchars=\\\{\}]
'The dataset is now completely clean\textbackslash{}n\textbackslash{}n- No missing values remain\textbackslash{}n- No
duplicate records exist\textbackslash{}n- No invalid values were detected\textbackslash{}n- All variables have
appropriate data types\textbackslash{}n\textbackslash{}nThe dataset is now ready for Exploratory Data Analysis
(EDA) and machine learning'
\end{Verbatim}
\end{tcolorbox}
        
    \begin{tcolorbox}[breakable, size=fbox, boxrule=1pt, pad at break*=1mm,colback=cellbackground, colframe=cellborder]
\prompt{In}{incolor}{49}{\boxspacing}
\begin{Verbatim}[commandchars=\\\{\}]
                        \PY{c+c1}{\PYZsh{}PART \PYZhy{}3(Exploratory Data Analysis (EDA))}
\end{Verbatim}
\end{tcolorbox}

    \begin{tcolorbox}[breakable, size=fbox, boxrule=1pt, pad at break*=1mm,colback=cellbackground, colframe=cellborder]
\prompt{In}{incolor}{50}{\boxspacing}
\begin{Verbatim}[commandchars=\\\{\}]
\PY{c+c1}{\PYZsh{}calculate mean,median.mode}
\PY{n}{mean} \PY{o}{=} \PY{n}{df}\PY{p}{[}\PY{n}{numerical\PYZus{}cols}\PY{p}{]}\PY{o}{.}\PY{n}{mean}\PY{p}{(}\PY{p}{)}
\PY{n}{median} \PY{o}{=} \PY{n}{df}\PY{p}{[}\PY{n}{numerical\PYZus{}cols}\PY{p}{]}\PY{o}{.}\PY{n}{median}\PY{p}{(}\PY{p}{)}
\PY{n}{mode} \PY{o}{=} \PY{n}{df}\PY{p}{[}\PY{n}{numerical\PYZus{}cols}\PY{p}{]}\PY{o}{.}\PY{n}{mode}\PY{p}{(}\PY{p}{)}\PY{o}{.}\PY{n}{iloc}\PY{p}{[}\PY{l+m+mi}{0}\PY{p}{]}
\end{Verbatim}
\end{tcolorbox}

    \begin{tcolorbox}[breakable, size=fbox, boxrule=1pt, pad at break*=1mm,colback=cellbackground, colframe=cellborder]
\prompt{In}{incolor}{51}{\boxspacing}
\begin{Verbatim}[commandchars=\\\{\}]
\PY{n+nb}{print}\PY{p}{(}\PY{n}{mean}\PY{p}{)}
\PY{n+nb}{print}\PY{p}{(}\PY{n}{median}\PY{p}{)}
\PY{n+nb}{print}\PY{p}{(}\PY{n}{mode}\PY{p}{)}
\end{Verbatim}
\end{tcolorbox}

    \begin{Verbatim}[commandchars=\\\{\}]
longitude               -119.569704
latitude                  35.631861
housing\_median\_age        28.639486
total\_rooms             2635.763081
total\_bedrooms           536.838857
population              1425.476744
households               499.539680
median\_income              3.870671
median\_house\_value    206855.816909
dtype: float64
longitude               -118.4900
latitude                  34.2600
housing\_median\_age        29.0000
total\_rooms             2127.0000
total\_bedrooms           435.0000
population              1166.0000
households               409.0000
median\_income              3.5348
median\_house\_value    179700.0000
dtype: float64
longitude               -118.310
latitude                  34.060
housing\_median\_age        52.000
total\_rooms             1527.000
total\_bedrooms           435.000
population               891.000
households               306.000
median\_income              3.125
median\_house\_value    500001.000
Name: 0, dtype: float64
    \end{Verbatim}

    \begin{tcolorbox}[breakable, size=fbox, boxrule=1pt, pad at break*=1mm,colback=cellbackground, colframe=cellborder]
\prompt{In}{incolor}{52}{\boxspacing}
\begin{Verbatim}[commandchars=\\\{\}]
\PY{l+s+sd}{\PYZsq{}\PYZsq{}\PYZsq{}\PYZhy{} Mean represents the average value.}
\PY{l+s+sd}{\PYZhy{} Median represents the middle value.}
\PY{l+s+sd}{\PYZhy{} Moderepresents the most frequently occurring value.}

\PY{l+s+sd}{Comparing these measures helps identify whether a feature is symmetrically distributed or skewed\PYZsq{}\PYZsq{}\PYZsq{}}
\end{Verbatim}
\end{tcolorbox}

            \begin{tcolorbox}[breakable, size=fbox, boxrule=.5pt, pad at break*=1mm, opacityfill=0]
\prompt{Out}{outcolor}{52}{\boxspacing}
\begin{Verbatim}[commandchars=\\\{\}]
'- Mean represents the average value.\textbackslash{}n- Median represents the middle value.\textbackslash{}n-
Moderepresents the most frequently occurring value.\textbackslash{}n\textbackslash{}nComparing these measures
helps identify whether a feature is symmetrically distributed or skewed'
\end{Verbatim}
\end{tcolorbox}
        
    \begin{tcolorbox}[breakable, size=fbox, boxrule=1pt, pad at break*=1mm,colback=cellbackground, colframe=cellborder]
\prompt{In}{incolor}{53}{\boxspacing}
\begin{Verbatim}[commandchars=\\\{\}]
\PY{c+c1}{\PYZsh{}calculating variamce ,SD,IQR}
\PY{n}{variance} \PY{o}{=} \PY{n}{df}\PY{p}{[}\PY{n}{numerical\PYZus{}cols}\PY{p}{]}\PY{o}{.}\PY{n}{var}\PY{p}{(}\PY{p}{)}

\PY{n}{std} \PY{o}{=} \PY{n}{df}\PY{p}{[}\PY{n}{numerical\PYZus{}cols}\PY{p}{]}\PY{o}{.}\PY{n}{std}\PY{p}{(}\PY{p}{)}

\PY{n}{Q1} \PY{o}{=} \PY{n}{df}\PY{p}{[}\PY{n}{numerical\PYZus{}cols}\PY{p}{]}\PY{o}{.}\PY{n}{quantile}\PY{p}{(}\PY{l+m+mf}{0.25}\PY{p}{)}

\PY{n}{Q3} \PY{o}{=} \PY{n}{df}\PY{p}{[}\PY{n}{numerical\PYZus{}cols}\PY{p}{]}\PY{o}{.}\PY{n}{quantile}\PY{p}{(}\PY{l+m+mf}{0.75}\PY{p}{)}

\PY{n}{IQR} \PY{o}{=} \PY{n}{Q3} \PY{o}{\PYZhy{}} \PY{n}{Q1}
\end{Verbatim}
\end{tcolorbox}

    \begin{tcolorbox}[breakable, size=fbox, boxrule=1pt, pad at break*=1mm,colback=cellbackground, colframe=cellborder]
\prompt{In}{incolor}{54}{\boxspacing}
\begin{Verbatim}[commandchars=\\\{\}]
\PY{n}{summary\PYZus{}table} \PY{o}{=} \PY{n}{pd}\PY{o}{.}\PY{n}{DataFrame}\PY{p}{(}\PY{p}{\PYZob{}}
    \PY{l+s+s2}{\PYZdq{}}\PY{l+s+s2}{Mean}\PY{l+s+s2}{\PYZdq{}}\PY{p}{:} \PY{n}{mean}\PY{p}{,}
    \PY{l+s+s2}{\PYZdq{}}\PY{l+s+s2}{Median}\PY{l+s+s2}{\PYZdq{}}\PY{p}{:} \PY{n}{median}\PY{p}{,}
    \PY{l+s+s2}{\PYZdq{}}\PY{l+s+s2}{Mode}\PY{l+s+s2}{\PYZdq{}}\PY{p}{:} \PY{n}{mode}\PY{p}{,}
    \PY{l+s+s2}{\PYZdq{}}\PY{l+s+s2}{Variance}\PY{l+s+s2}{\PYZdq{}}\PY{p}{:} \PY{n}{variance}\PY{p}{,}
    \PY{l+s+s2}{\PYZdq{}}\PY{l+s+s2}{Standard Deviation}\PY{l+s+s2}{\PYZdq{}}\PY{p}{:} \PY{n}{std}\PY{p}{,}
    \PY{l+s+s2}{\PYZdq{}}\PY{l+s+s2}{Q1}\PY{l+s+s2}{\PYZdq{}}\PY{p}{:} \PY{n}{Q1}\PY{p}{,}
    \PY{l+s+s2}{\PYZdq{}}\PY{l+s+s2}{Q3}\PY{l+s+s2}{\PYZdq{}}\PY{p}{:} \PY{n}{Q3}\PY{p}{,}
    \PY{l+s+s2}{\PYZdq{}}\PY{l+s+s2}{IQR}\PY{l+s+s2}{\PYZdq{}}\PY{p}{:} \PY{n}{IQR}
\PY{p}{\PYZcb{}}\PY{p}{)}

\PY{n}{summary\PYZus{}table}
\end{Verbatim}
\end{tcolorbox}

            \begin{tcolorbox}[breakable, size=fbox, boxrule=.5pt, pad at break*=1mm, opacityfill=0]
\prompt{Out}{outcolor}{54}{\boxspacing}
\begin{Verbatim}[commandchars=\\\{\}]
                             Mean       Median        Mode      Variance  \textbackslash{}
longitude             -119.569704    -118.4900    -118.310  4.014139e+00
latitude                35.631861      34.2600      34.060  4.562293e+00
housing\_median\_age      28.639486      29.0000      52.000  1.583963e+02
total\_rooms           2635.763081    2127.0000    1527.000  4.759445e+06
total\_bedrooms         536.838857     435.0000     435.000  1.758895e+05
population            1425.476744    1166.0000     891.000  1.282470e+06
households             499.539680     409.0000     306.000  1.461760e+05
median\_income            3.870671       3.5348       3.125  3.609323e+00
median\_house\_value  206855.816909  179700.0000  500001.000  1.331615e+10

                    Standard Deviation           Q1            Q3  \textbackslash{}
longitude                     2.003532    -121.8000    -118.01000
latitude                      2.135952      33.9300      37.71000
housing\_median\_age           12.585558      18.0000      37.00000
total\_rooms                2181.615252    1447.7500    3148.00000
total\_bedrooms              419.391878     297.0000     643.25000
population                 1132.462122     787.0000    1725.00000
households                  382.329753     280.0000     605.00000
median\_income                 1.899822       2.5634       4.74325
median\_house\_value       115395.615874  119600.0000  264725.00000

                             IQR
longitude                3.79000
latitude                 3.78000
housing\_median\_age      19.00000
total\_rooms           1700.25000
total\_bedrooms         346.25000
population             938.00000
households             325.00000
median\_income            2.17985
median\_house\_value  145125.00000
\end{Verbatim}
\end{tcolorbox}
        
    \begin{tcolorbox}[breakable, size=fbox, boxrule=1pt, pad at break*=1mm,colback=cellbackground, colframe=cellborder]
\prompt{In}{incolor}{55}{\boxspacing}
\begin{Verbatim}[commandchars=\\\{\}]
\PY{c+c1}{\PYZsh{}identify skewed features}
\PY{k}{for} \PY{n}{col} \PY{o+ow}{in} \PY{n}{numerical\PYZus{}cols}\PY{p}{:}
    \PY{k}{if} \PY{n}{mean}\PY{p}{[}\PY{n}{col}\PY{p}{]} \PY{o}{\PYZgt{}} \PY{n}{median}\PY{p}{[}\PY{n}{col}\PY{p}{]}\PY{p}{:}
        \PY{n+nb}{print}\PY{p}{(}\PY{l+s+sa}{f}\PY{l+s+s2}{\PYZdq{}}\PY{l+s+si}{\PYZob{}}\PY{n}{col}\PY{l+s+si}{\PYZcb{}}\PY{l+s+s2}{: Positively (Right) Skewed}\PY{l+s+s2}{\PYZdq{}}\PY{p}{)}
    \PY{k}{elif} \PY{n}{mean}\PY{p}{[}\PY{n}{col}\PY{p}{]} \PY{o}{\PYZlt{}} \PY{n}{median}\PY{p}{[}\PY{n}{col}\PY{p}{]}\PY{p}{:}
        \PY{n+nb}{print}\PY{p}{(}\PY{l+s+sa}{f}\PY{l+s+s2}{\PYZdq{}}\PY{l+s+si}{\PYZob{}}\PY{n}{col}\PY{l+s+si}{\PYZcb{}}\PY{l+s+s2}{: Negatively (Left) Skewed}\PY{l+s+s2}{\PYZdq{}}\PY{p}{)}
    \PY{k}{else}\PY{p}{:}
        \PY{n+nb}{print}\PY{p}{(}\PY{l+s+sa}{f}\PY{l+s+s2}{\PYZdq{}}\PY{l+s+si}{\PYZob{}}\PY{n}{col}\PY{l+s+si}{\PYZcb{}}\PY{l+s+s2}{: Approximately Symmetric}\PY{l+s+s2}{\PYZdq{}}\PY{p}{)}
\end{Verbatim}
\end{tcolorbox}

    \begin{Verbatim}[commandchars=\\\{\}]
longitude: Negatively (Left) Skewed
latitude: Positively (Right) Skewed
housing\_median\_age: Negatively (Left) Skewed
total\_rooms: Positively (Right) Skewed
total\_bedrooms: Positively (Right) Skewed
population: Positively (Right) Skewed
households: Positively (Right) Skewed
median\_income: Positively (Right) Skewed
median\_house\_value: Positively (Right) Skewed
    \end{Verbatim}

    \begin{tcolorbox}[breakable, size=fbox, boxrule=1pt, pad at break*=1mm,colback=cellbackground, colframe=cellborder]
\prompt{In}{incolor}{56}{\boxspacing}
\begin{Verbatim}[commandchars=\\\{\}]
\PY{l+s+sd}{\PYZsq{}\PYZsq{}\PYZsq{}}
\PY{l+s+sd}{A feature is:}

\PY{l+s+sd}{\PYZhy{} Right\PYZhy{}skewed when Mean \PYZgt{} Median}
\PY{l+s+sd}{\PYZhy{} Left\PYZhy{}skewed when Mean \PYZlt{} Median}
\PY{l+s+sd}{\PYZhy{} Approximately symmetric when Mean ≈ Median}

\PY{l+s+sd}{This comparison provides a simple indication of the distribution shape\PYZsq{}\PYZsq{}\PYZsq{}}
\end{Verbatim}
\end{tcolorbox}

            \begin{tcolorbox}[breakable, size=fbox, boxrule=.5pt, pad at break*=1mm, opacityfill=0]
\prompt{Out}{outcolor}{56}{\boxspacing}
\begin{Verbatim}[commandchars=\\\{\}]
'\textbackslash{}nA feature is:\textbackslash{}n\textbackslash{}n- Right-skewed when Mean > Median\textbackslash{}n- Left-skewed when Mean <
Median\textbackslash{}n- Approximately symmetric when Mean ≈ Median\textbackslash{}n\textbackslash{}nThis comparison provides
a simple indication of the distribution shape'
\end{Verbatim}
\end{tcolorbox}
        
    \begin{tcolorbox}[breakable, size=fbox, boxrule=1pt, pad at break*=1mm,colback=cellbackground, colframe=cellborder]
\prompt{In}{incolor}{57}{\boxspacing}
\begin{Verbatim}[commandchars=\\\{\}]
\PY{c+c1}{\PYZsh{}calculating Actual Skewnwss}
\PY{n}{skewness} \PY{o}{=} \PY{n}{df}\PY{p}{[}\PY{n}{numerical\PYZus{}cols}\PY{p}{]}\PY{o}{.}\PY{n}{skew}\PY{p}{(}\PY{p}{)}

\PY{n+nb}{print}\PY{p}{(}\PY{n}{skewness}\PY{p}{)}
\end{Verbatim}
\end{tcolorbox}

    \begin{Verbatim}[commandchars=\\\{\}]
longitude            -0.297801
latitude              0.465953
housing\_median\_age    0.060331
total\_rooms           4.147343
total\_bedrooms        3.481141
population            4.935858
households            3.410438
median\_income         1.646657
median\_house\_value    0.977763
dtype: float64
    \end{Verbatim}

    \begin{tcolorbox}[breakable, size=fbox, boxrule=1pt, pad at break*=1mm,colback=cellbackground, colframe=cellborder]
\prompt{In}{incolor}{58}{\boxspacing}
\begin{Verbatim}[commandchars=\\\{\}]
\PY{c+c1}{\PYZsh{}fina; Summary table including skewness}
\PY{n}{summary\PYZus{}table}\PY{p}{[}\PY{l+s+s2}{\PYZdq{}}\PY{l+s+s2}{Skewness}\PY{l+s+s2}{\PYZdq{}}\PY{p}{]} \PY{o}{=} \PY{n}{df}\PY{p}{[}\PY{n}{numerical\PYZus{}cols}\PY{p}{]}\PY{o}{.}\PY{n}{skew}\PY{p}{(}\PY{p}{)}

\PY{n}{summary\PYZus{}table}
\end{Verbatim}
\end{tcolorbox}

            \begin{tcolorbox}[breakable, size=fbox, boxrule=.5pt, pad at break*=1mm, opacityfill=0]
\prompt{Out}{outcolor}{58}{\boxspacing}
\begin{Verbatim}[commandchars=\\\{\}]
                             Mean       Median        Mode      Variance  \textbackslash{}
longitude             -119.569704    -118.4900    -118.310  4.014139e+00
latitude                35.631861      34.2600      34.060  4.562293e+00
housing\_median\_age      28.639486      29.0000      52.000  1.583963e+02
total\_rooms           2635.763081    2127.0000    1527.000  4.759445e+06
total\_bedrooms         536.838857     435.0000     435.000  1.758895e+05
population            1425.476744    1166.0000     891.000  1.282470e+06
households             499.539680     409.0000     306.000  1.461760e+05
median\_income            3.870671       3.5348       3.125  3.609323e+00
median\_house\_value  206855.816909  179700.0000  500001.000  1.331615e+10

                    Standard Deviation           Q1            Q3  \textbackslash{}
longitude                     2.003532    -121.8000    -118.01000
latitude                      2.135952      33.9300      37.71000
housing\_median\_age           12.585558      18.0000      37.00000
total\_rooms                2181.615252    1447.7500    3148.00000
total\_bedrooms              419.391878     297.0000     643.25000
population                 1132.462122     787.0000    1725.00000
households                  382.329753     280.0000     605.00000
median\_income                 1.899822       2.5634       4.74325
median\_house\_value       115395.615874  119600.0000  264725.00000

                             IQR  Skewness
longitude                3.79000 -0.297801
latitude                 3.78000  0.465953
housing\_median\_age      19.00000  0.060331
total\_rooms           1700.25000  4.147343
total\_bedrooms         346.25000  3.481141
population             938.00000  4.935858
households             325.00000  3.410438
median\_income            2.17985  1.646657
median\_house\_value  145125.00000  0.977763
\end{Verbatim}
\end{tcolorbox}
        
    \begin{tcolorbox}[breakable, size=fbox, boxrule=1pt, pad at break*=1mm,colback=cellbackground, colframe=cellborder]
\prompt{In}{incolor}{59}{\boxspacing}
\begin{Verbatim}[commandchars=\\\{\}]
\PY{l+s+sd}{\PYZsq{}\PYZsq{}\PYZsq{}}
\PY{l+s+sd}{\PYZsh{}\PYZsh{}\PYZsh{} }

\PY{l+s+sd}{The summary statistics table combines all important descriptive measures into a single table}
\PY{l+s+sd}{It provides a comprehensive overview of the dataset and highlights differences in central tendency and variability across features\PYZsq{}\PYZsq{}\PYZsq{}}
\end{Verbatim}
\end{tcolorbox}

            \begin{tcolorbox}[breakable, size=fbox, boxrule=.5pt, pad at break*=1mm, opacityfill=0]
\prompt{Out}{outcolor}{59}{\boxspacing}
\begin{Verbatim}[commandchars=\\\{\}]
'\textbackslash{}n\#\#\# \textbackslash{}n\textbackslash{}nThe summary statistics table combines all important descriptive
measures into a single table\textbackslash{}nIt provides a comprehensive overview of the
dataset and highlights differences in central tendency and variability across
features'
\end{Verbatim}
\end{tcolorbox}
        
    \begin{tcolorbox}[breakable, size=fbox, boxrule=1pt, pad at break*=1mm,colback=cellbackground, colframe=cellborder]
\prompt{In}{incolor}{60}{\boxspacing}
\begin{Verbatim}[commandchars=\\\{\}]
\PY{c+c1}{\PYZsh{}Histogram}
\PY{n}{numerical\PYZus{}cols} \PY{o}{=} \PY{n}{df}\PY{o}{.}\PY{n}{select\PYZus{}dtypes}\PY{p}{(}\PY{n}{include}\PY{o}{=}\PY{p}{[}\PY{l+s+s1}{\PYZsq{}}\PY{l+s+s1}{int64}\PY{l+s+s1}{\PYZsq{}}\PY{p}{,} \PY{l+s+s1}{\PYZsq{}}\PY{l+s+s1}{float64}\PY{l+s+s1}{\PYZsq{}}\PY{p}{]}\PY{p}{)}\PY{o}{.}\PY{n}{columns}

\PY{k}{for} \PY{n}{col} \PY{o+ow}{in} \PY{n}{numerical\PYZus{}cols}\PY{p}{:}
    \PY{n}{plt}\PY{o}{.}\PY{n}{figure}\PY{p}{(}\PY{n}{figsize}\PY{o}{=}\PY{p}{(}\PY{l+m+mi}{8}\PY{p}{,}\PY{l+m+mi}{5}\PY{p}{)}\PY{p}{)}

    \PY{n}{plt}\PY{o}{.}\PY{n}{hist}\PY{p}{(}\PY{n}{df}\PY{p}{[}\PY{n}{col}\PY{p}{]}\PY{p}{,} \PY{n}{bins}\PY{o}{=}\PY{l+m+mi}{30}\PY{p}{,} \PY{n}{edgecolor}\PY{o}{=}\PY{l+s+s1}{\PYZsq{}}\PY{l+s+s1}{black}\PY{l+s+s1}{\PYZsq{}}\PY{p}{,} \PY{n}{color}\PY{o}{=}\PY{l+s+s1}{\PYZsq{}}\PY{l+s+s1}{skyblue}\PY{l+s+s1}{\PYZsq{}}\PY{p}{)}

    \PY{n}{plt}\PY{o}{.}\PY{n}{title}\PY{p}{(}\PY{l+s+sa}{f}\PY{l+s+s2}{\PYZdq{}}\PY{l+s+s2}{Distribution of }\PY{l+s+si}{\PYZob{}}\PY{n}{col}\PY{l+s+si}{\PYZcb{}}\PY{l+s+s2}{\PYZdq{}}\PY{p}{,} \PY{n}{fontsize}\PY{o}{=}\PY{l+m+mi}{14}\PY{p}{)}
    \PY{n}{plt}\PY{o}{.}\PY{n}{xlabel}\PY{p}{(}\PY{n}{col}\PY{p}{,} \PY{n}{fontsize}\PY{o}{=}\PY{l+m+mi}{12}\PY{p}{)}
    \PY{n}{plt}\PY{o}{.}\PY{n}{ylabel}\PY{p}{(}\PY{l+s+s2}{\PYZdq{}}\PY{l+s+s2}{Frequency}\PY{l+s+s2}{\PYZdq{}}\PY{p}{,} \PY{n}{fontsize}\PY{o}{=}\PY{l+m+mi}{12}\PY{p}{)}

    \PY{n}{plt}\PY{o}{.}\PY{n}{grid}\PY{p}{(}\PY{k+kc}{True}\PY{p}{)}
    \PY{n}{plt}\PY{o}{.}\PY{n}{show}\PY{p}{(}\PY{p}{)}
\end{Verbatim}
\end{tcolorbox}

    \begin{center}
    \adjustimage{max size={0.9\linewidth}{0.9\paperheight}}{output_60_0.png}
    \end{center}
    { \hspace*{\fill} \\}
    
    \begin{center}
    \adjustimage{max size={0.9\linewidth}{0.9\paperheight}}{output_60_1.png}
    \end{center}
    { \hspace*{\fill} \\}
    
    \begin{center}
    \adjustimage{max size={0.9\linewidth}{0.9\paperheight}}{output_60_2.png}
    \end{center}
    { \hspace*{\fill} \\}
    
    \begin{center}
    \adjustimage{max size={0.9\linewidth}{0.9\paperheight}}{output_60_3.png}
    \end{center}
    { \hspace*{\fill} \\}
    
    \begin{center}
    \adjustimage{max size={0.9\linewidth}{0.9\paperheight}}{output_60_4.png}
    \end{center}
    { \hspace*{\fill} \\}
    
    \begin{center}
    \adjustimage{max size={0.9\linewidth}{0.9\paperheight}}{output_60_5.png}
    \end{center}
    { \hspace*{\fill} \\}
    
    \begin{center}
    \adjustimage{max size={0.9\linewidth}{0.9\paperheight}}{output_60_6.png}
    \end{center}
    { \hspace*{\fill} \\}
    
    \begin{center}
    \adjustimage{max size={0.9\linewidth}{0.9\paperheight}}{output_60_7.png}
    \end{center}
    { \hspace*{\fill} \\}
    
    \begin{center}
    \adjustimage{max size={0.9\linewidth}{0.9\paperheight}}{output_60_8.png}
    \end{center}
    { \hspace*{\fill} \\}
    
    \begin{tcolorbox}[breakable, size=fbox, boxrule=1pt, pad at break*=1mm,colback=cellbackground, colframe=cellborder]
\prompt{In}{incolor}{61}{\boxspacing}
\begin{Verbatim}[commandchars=\\\{\}]
\PY{l+s+sd}{\PYZsq{}\PYZsq{}\PYZsq{}}
\PY{l+s+sd}{1] Longitude:}
\PY{l+s+sd}{The longitude distribution is not uniform, indicating that houses are concentrated in specific regions of California rather than being evenly distributed. This reflects the geographical clustering of residential areas}
\PY{l+s+sd}{2] Latitude:}
\PY{l+s+sd}{The latitude distribution shows that most housing blocks are concentrated within certain latitude ranges. This suggests that the population is unevenly distributed across California.}
\PY{l+s+sd}{3]Housing Median Age:}
\PY{l+s+sd}{The housing median age is moderately distributed, with many houses falling between 20 and 40 years old. The distribution indicates that California contains a mix of older and newer residential areas.}
\PY{l+s+sd}{4] Total Rooms:}
\PY{l+s+sd}{The distribution of total rooms is highly right\PYZhy{}skewed. Most housing blocks contain a relatively small number of rooms, while a few blocks have extremely large numbers of rooms. }
\PY{l+s+sd}{These extreme values may represent densely populated neighborhoods or unusually large residential developments.}
\PY{l+s+sd}{5] Total Bedrooms:}
\PY{l+s+sd}{The histogram of total bedrooms is also right\PYZhy{}skewed. Most observations have fewer bedrooms, while a small number of housing blocks contain exceptionally high bedroom counts.}
\PY{l+s+sd}{These may be considered potential outliers.}
\PY{l+s+sd}{6] Population:}
\PY{l+s+sd}{The population distribution is strongly right\PYZhy{}skewed. Most housing blocks have relatively small populations, whereas a few blocks contain very large populations. This indicates considerable variation in population size across different regions}
\PY{l+s+sd}{7] Households:}
\PY{l+s+sd}{The households feature follows a right\PYZhy{}skewed distribution. Most blocks have a moderate number of households, while a few blocks contain a significantly larger number of households.}
\PY{l+s+sd}{8] Median Income:}
\PY{l+s+sd}{Median income is positively skewed, with most housing blocks having moderate income levels and fewer blocks having very high incomes. This suggests income inequality across different regions of California.}
\PY{l+s+sd}{9]Median House Value (Target Variable):}
\PY{l+s+sd}{The target variable, median house value, is positively skewed. Most houses are priced in the lower to middle range, while fewer houses have very high values. }
\PY{l+s+sd}{The sharp cutoff near the maximum value indicates that house prices were capped in the dataset, meaning values above this limit were not recorded\PYZsq{}\PYZsq{}\PYZsq{}}
\end{Verbatim}
\end{tcolorbox}

            \begin{tcolorbox}[breakable, size=fbox, boxrule=.5pt, pad at break*=1mm, opacityfill=0]
\prompt{Out}{outcolor}{61}{\boxspacing}
\begin{Verbatim}[commandchars=\\\{\}]
'\textbackslash{}n1] Longitude:\textbackslash{}nThe longitude distribution is not uniform, indicating that
houses are concentrated in specific regions of California rather than being
evenly distributed. This reflects the geographical clustering of residential
areas\textbackslash{}n2] Latitude:\textbackslash{}nThe latitude distribution shows that most housing blocks
are concentrated within certain latitude ranges. This suggests that the
population is unevenly distributed across California.\textbackslash{}n3]Housing Median
Age:\textbackslash{}nThe housing median age is moderately distributed, with many houses falling
between 20 and 40 years old. The distribution indicates that California contains
a mix of older and newer residential areas.\textbackslash{}n4] Total Rooms:\textbackslash{}nThe distribution
of total rooms is highly right-skewed. Most housing blocks contain a relatively
small number of rooms, while a few blocks have extremely large numbers of rooms.
\textbackslash{}nThese extreme values may represent densely populated neighborhoods or
unusually large residential developments.\textbackslash{}n5] Total Bedrooms:\textbackslash{}nThe histogram of
total bedrooms is also right-skewed. Most observations have fewer bedrooms,
while a small number of housing blocks contain exceptionally high bedroom
counts.\textbackslash{}nThese may be considered potential outliers.\textbackslash{}n6] Population:\textbackslash{}nThe
population distribution is strongly right-skewed. Most housing blocks have
relatively small populations, whereas a few blocks contain very large
populations. This indicates considerable variation in population size across
different regions\textbackslash{}n7] Households:\textbackslash{}nThe households feature follows a right-skewed
distribution. Most blocks have a moderate number of households, while a few
blocks contain a significantly larger number of households.\textbackslash{}n8] Median
Income:\textbackslash{}nMedian income is positively skewed, with most housing blocks having
moderate income levels and fewer blocks having very high incomes. This suggests
income inequality across different regions of California.\textbackslash{}n9]Median House Value
(Target Variable):\textbackslash{}nThe target variable, median house value, is positively
skewed. Most houses are priced in the lower to middle range, while fewer houses
have very high values. \textbackslash{}nThe sharp cutoff near the maximum value indicates that
house prices were capped in the dataset, meaning values above this limit were
not recorded'
\end{Verbatim}
\end{tcolorbox}
        
    \begin{tcolorbox}[breakable, size=fbox, boxrule=1pt, pad at break*=1mm,colback=cellbackground, colframe=cellborder]
\prompt{In}{incolor}{62}{\boxspacing}
\begin{Verbatim}[commandchars=\\\{\}]
\PY{l+s+sd}{\PYZsq{}\PYZsq{}\PYZsq{}Overall Histogram Analysis (Good for Assignment):}

\PY{l+s+sd}{The histograms reveal that most numerical features are positively (right) skewed , meaning that the majority of observations are concentrated at lower values with a small number of extremely high values. Features such as total\PYZus{}rooms, **total\PYZus{}bedrooms, population, and households exhibit noticeable outliers}

\PY{l+s+sd}{The median\PYZus{}income feature also shows a right\PYZhy{}skewed distribution, suggesting that most households earn moderate incomes while only a small proportion have very high incomes}

\PY{l+s+sd}{The target variable, median\PYZus{}house\PYZus{}value, is also right\PYZhy{}skewed and shows an upper limit, indicating that house prices were capped in the dataset. These observations suggest that the dataset contains non\PYZhy{}normal distributions and outliers, which are common in real\PYZhy{}world housing data and should be considered during model development\PYZsq{}\PYZsq{}\PYZsq{}}
\end{Verbatim}
\end{tcolorbox}

            \begin{tcolorbox}[breakable, size=fbox, boxrule=.5pt, pad at break*=1mm, opacityfill=0]
\prompt{Out}{outcolor}{62}{\boxspacing}
\begin{Verbatim}[commandchars=\\\{\}]
'Overall Histogram Analysis (Good for Assignment):\textbackslash{}n\textbackslash{}nThe histograms reveal that
most numerical features are positively (right) skewed , meaning that the
majority of observations are concentrated at lower values with a small number of
extremely high values. Features such as total\_rooms, **total\_bedrooms,
population, and households exhibit noticeable outliers\textbackslash{}n\textbackslash{}nThe median\_income
feature also shows a right-skewed distribution, suggesting that most households
earn moderate incomes while only a small proportion have very high
incomes\textbackslash{}n\textbackslash{}nThe target variable, median\_house\_value, is also right-skewed and
shows an upper limit, indicating that house prices were capped in the dataset.
These observations suggest that the dataset contains non-normal distributions
and outliers, which are common in real-world housing data and should be
considered during model development'
\end{Verbatim}
\end{tcolorbox}
        
    \begin{tcolorbox}[breakable, size=fbox, boxrule=1pt, pad at break*=1mm,colback=cellbackground, colframe=cellborder]
\prompt{In}{incolor}{63}{\boxspacing}
\begin{Verbatim}[commandchars=\\\{\}]
\PY{c+c1}{\PYZsh{}Box Plots to identify Outliers}
\PY{k}{for} \PY{n}{col} \PY{o+ow}{in} \PY{n}{numerical\PYZus{}cols}\PY{p}{:}
    \PY{n}{plt}\PY{o}{.}\PY{n}{figure}\PY{p}{(}\PY{n}{figsize}\PY{o}{=}\PY{p}{(}\PY{l+m+mi}{8}\PY{p}{,}\PY{l+m+mi}{4}\PY{p}{)}\PY{p}{)}

    \PY{n}{sns}\PY{o}{.}\PY{n}{boxplot}\PY{p}{(}\PY{n}{x}\PY{o}{=}\PY{n}{df}\PY{p}{[}\PY{n}{col}\PY{p}{]}\PY{p}{,} \PY{n}{color}\PY{o}{=}\PY{l+s+s2}{\PYZdq{}}\PY{l+s+s2}{lightgreen}\PY{l+s+s2}{\PYZdq{}}\PY{p}{)}

    \PY{n}{plt}\PY{o}{.}\PY{n}{title}\PY{p}{(}\PY{l+s+sa}{f}\PY{l+s+s2}{\PYZdq{}}\PY{l+s+s2}{Box Plot of }\PY{l+s+si}{\PYZob{}}\PY{n}{col}\PY{l+s+si}{\PYZcb{}}\PY{l+s+s2}{\PYZdq{}}\PY{p}{,} \PY{n}{fontsize}\PY{o}{=}\PY{l+m+mi}{14}\PY{p}{)}
    \PY{n}{plt}\PY{o}{.}\PY{n}{xlabel}\PY{p}{(}\PY{n}{col}\PY{p}{)}

    \PY{n}{plt}\PY{o}{.}\PY{n}{show}\PY{p}{(}\PY{p}{)}
\end{Verbatim}
\end{tcolorbox}

    \begin{center}
    \adjustimage{max size={0.9\linewidth}{0.9\paperheight}}{output_63_0.png}
    \end{center}
    { \hspace*{\fill} \\}
    
    \begin{center}
    \adjustimage{max size={0.9\linewidth}{0.9\paperheight}}{output_63_1.png}
    \end{center}
    { \hspace*{\fill} \\}
    
    \begin{center}
    \adjustimage{max size={0.9\linewidth}{0.9\paperheight}}{output_63_2.png}
    \end{center}
    { \hspace*{\fill} \\}
    
    \begin{center}
    \adjustimage{max size={0.9\linewidth}{0.9\paperheight}}{output_63_3.png}
    \end{center}
    { \hspace*{\fill} \\}
    
    \begin{center}
    \adjustimage{max size={0.9\linewidth}{0.9\paperheight}}{output_63_4.png}
    \end{center}
    { \hspace*{\fill} \\}
    
    \begin{center}
    \adjustimage{max size={0.9\linewidth}{0.9\paperheight}}{output_63_5.png}
    \end{center}
    { \hspace*{\fill} \\}
    
    \begin{center}
    \adjustimage{max size={0.9\linewidth}{0.9\paperheight}}{output_63_6.png}
    \end{center}
    { \hspace*{\fill} \\}
    
    \begin{center}
    \adjustimage{max size={0.9\linewidth}{0.9\paperheight}}{output_63_7.png}
    \end{center}
    { \hspace*{\fill} \\}
    
    \begin{center}
    \adjustimage{max size={0.9\linewidth}{0.9\paperheight}}{output_63_8.png}
    \end{center}
    { \hspace*{\fill} \\}
    
    \begin{tcolorbox}[breakable, size=fbox, boxrule=1pt, pad at break*=1mm,colback=cellbackground, colframe=cellborder]
\prompt{In}{incolor}{64}{\boxspacing}
\begin{Verbatim}[commandchars=\\\{\}]
\PY{l+s+sd}{\PYZsq{}\PYZsq{}\PYZsq{}}
\PY{l+s+sd}{1] Longitude:}
\PY{l+s+sd}{The box plot of longitude shows a moderate spread of values with a few outliers. These outliers represent housing locations that are farther east or west compared to the majority of California housing blocks.}

\PY{l+s+sd}{2] Latitude:}
\PY{l+s+sd}{The latitude box plot indicates a reasonable distribution of housing locations with a few outliers. These outliers correspond to housing blocks located in less common northern or southern regions of California.}

\PY{l+s+sd}{3] Housing Median Age:}
\PY{l+s+sd}{The housing median age has relatively few outliers, indicating that most houses fall within a similar age range. The majority of homes are between 20 and 40 years old, showing a balanced distribution of housing ages.}

\PY{l+s+sd}{4] Total Rooms:}
\PY{l+s+sd}{The total rooms feature contains a large number of outliers beyond the upper whisker. These outliers represent housing blocks with exceptionally high numbers of rooms, indicating a highly right\PYZhy{}skewed distribution.}

\PY{l+s+sd}{5] Total Bedrooms:}
\PY{l+s+sd}{The total bedrooms feature also shows numerous outliers. Most housing blocks have a moderate number of bedrooms, while a small number contain exceptionally large bedroom counts, contributing to a positively skewed distribution.}

\PY{l+s+sd}{6] Population:}
\PY{l+s+sd}{The population feature exhibits many extreme outliers. While most housing blocks have moderate population sizes, a few blocks contain very large populations, resulting in a long upper tail.}

\PY{l+s+sd}{7] Households:}
\PY{l+s+sd}{The households feature has several outliers above the upper whisker. These represent housing blocks with unusually high numbers of households compared to the majority of the dataset.}

\PY{l+s+sd}{8] Median Income:}
\PY{l+s+sd}{The median income feature contains a few high\PYZhy{}income outliers. Most housing blocks have moderate income levels, while only a small number of areas report exceptionally high median incomes.}

\PY{l+s+sd}{9] Median House Value (Target Variable):}
\PY{l+s+sd}{The median house value feature shows several high\PYZhy{}value outliers and an upper limit near \PYZdl{}500,000. This suggests that house prices were capped in the dataset, causing the highest\PYZhy{}priced houses to be grouped at the maximum recorded value.}
\PY{l+s+sd}{\PYZsq{}\PYZsq{}\PYZsq{}}
\end{Verbatim}
\end{tcolorbox}

            \begin{tcolorbox}[breakable, size=fbox, boxrule=.5pt, pad at break*=1mm, opacityfill=0]
\prompt{Out}{outcolor}{64}{\boxspacing}
\begin{Verbatim}[commandchars=\\\{\}]
'\textbackslash{}n1] Longitude:\textbackslash{}nThe box plot of longitude shows a moderate spread of values
with a few outliers. These outliers represent housing locations that are farther
east or west compared to the majority of California housing blocks.\textbackslash{}n\textbackslash{}n2]
Latitude:\textbackslash{}nThe latitude box plot indicates a reasonable distribution of housing
locations with a few outliers. These outliers correspond to housing blocks
located in less common northern or southern regions of California.\textbackslash{}n\textbackslash{}n3] Housing
Median Age:\textbackslash{}nThe housing median age has relatively few outliers, indicating that
most houses fall within a similar age range. The majority of homes are between
20 and 40 years old, showing a balanced distribution of housing ages.\textbackslash{}n\textbackslash{}n4]
Total Rooms:\textbackslash{}nThe total rooms feature contains a large number of outliers beyond
the upper whisker. These outliers represent housing blocks with exceptionally
high numbers of rooms, indicating a highly right-skewed distribution.\textbackslash{}n\textbackslash{}n5]
Total Bedrooms:\textbackslash{}nThe total bedrooms feature also shows numerous outliers. Most
housing blocks have a moderate number of bedrooms, while a small number contain
exceptionally large bedroom counts, contributing to a positively skewed
distribution.\textbackslash{}n\textbackslash{}n6] Population:\textbackslash{}nThe population feature exhibits many extreme
outliers. While most housing blocks have moderate population sizes, a few blocks
contain very large populations, resulting in a long upper tail.\textbackslash{}n\textbackslash{}n7]
Households:\textbackslash{}nThe households feature has several outliers above the upper
whisker. These represent housing blocks with unusually high numbers of
households compared to the majority of the dataset.\textbackslash{}n\textbackslash{}n8] Median Income:\textbackslash{}nThe
median income feature contains a few high-income outliers. Most housing blocks
have moderate income levels, while only a small number of areas report
exceptionally high median incomes.\textbackslash{}n\textbackslash{}n9] Median House Value (Target
Variable):\textbackslash{}nThe median house value feature shows several high-value outliers and
an upper limit near \$500,000. This suggests that house prices were capped in the
dataset, causing the highest-priced houses to be grouped at the maximum recorded
value.\textbackslash{}n'
\end{Verbatim}
\end{tcolorbox}
        
    \begin{tcolorbox}[breakable, size=fbox, boxrule=1pt, pad at break*=1mm,colback=cellbackground, colframe=cellborder]
\prompt{In}{incolor}{65}{\boxspacing}
\begin{Verbatim}[commandchars=\\\{\}]
\PY{l+s+sd}{\PYZsq{}\PYZsq{}\PYZsq{}}
\PY{l+s+sd}{Overall Analysis:}
\PY{l+s+sd}{The box plots reveal that several numerical features contain significant outliers, particularly total\PYZus{}rooms, total\PYZus{}bedrooms, population, and households. These outliers indicate that some housing blocks have exceptionally large values compared to the majority of the dataset.}
\PY{l+s+sd}{Such outliers are common in real\PYZhy{}world housing datasets and may represent densely populated neighborhoods or unusually large residential developments. Since these values appear to be valid observations rather than data entry errors, they were retained for further analysis and model building.}
\PY{l+s+sd}{\PYZsq{}\PYZsq{}\PYZsq{}}
\end{Verbatim}
\end{tcolorbox}

            \begin{tcolorbox}[breakable, size=fbox, boxrule=.5pt, pad at break*=1mm, opacityfill=0]
\prompt{Out}{outcolor}{65}{\boxspacing}
\begin{Verbatim}[commandchars=\\\{\}]
'\textbackslash{}nOverall Analysis:\textbackslash{}nThe box plots reveal that several numerical features
contain significant outliers, particularly total\_rooms, total\_bedrooms,
population, and households. These outliers indicate that some housing blocks
have exceptionally large values compared to the majority of the dataset.\textbackslash{}nSuch
outliers are common in real-world housing datasets and may represent densely
populated neighborhoods or unusually large residential developments. Since these
values appear to be valid observations rather than data entry errors, they were
retained for further analysis and model building.\textbackslash{}n'
\end{Verbatim}
\end{tcolorbox}
        
    \begin{tcolorbox}[breakable, size=fbox, boxrule=1pt, pad at break*=1mm,colback=cellbackground, colframe=cellborder]
\prompt{In}{incolor}{66}{\boxspacing}
\begin{Verbatim}[commandchars=\\\{\}]
\PY{c+c1}{\PYZsh{}Scatter plots with Target Variable}
\PY{n}{target} \PY{o}{=} \PY{l+s+s2}{\PYZdq{}}\PY{l+s+s2}{median\PYZus{}house\PYZus{}value}\PY{l+s+s2}{\PYZdq{}}
\PY{n}{features} \PY{o}{=} \PY{n}{numerical\PYZus{}cols}\PY{o}{.}\PY{n}{drop}\PY{p}{(}\PY{n}{target}\PY{p}{)}

\PY{k}{for} \PY{n}{col} \PY{o+ow}{in} \PY{n}{features}\PY{p}{:}
    \PY{n}{plt}\PY{o}{.}\PY{n}{figure}\PY{p}{(}\PY{n}{figsize}\PY{o}{=}\PY{p}{(}\PY{l+m+mi}{7}\PY{p}{,}\PY{l+m+mi}{5}\PY{p}{)}\PY{p}{)}

    \PY{n}{plt}\PY{o}{.}\PY{n}{scatter}\PY{p}{(}\PY{n}{df}\PY{p}{[}\PY{n}{col}\PY{p}{]}\PY{p}{,}
                \PY{n}{df}\PY{p}{[}\PY{n}{target}\PY{p}{]}\PY{p}{,}
                \PY{n}{alpha}\PY{o}{=}\PY{l+m+mf}{0.5}\PY{p}{,}
                \PY{n}{color}\PY{o}{=}\PY{l+s+s2}{\PYZdq{}}\PY{l+s+s2}{tomato}\PY{l+s+s2}{\PYZdq{}}\PY{p}{)}

    \PY{n}{plt}\PY{o}{.}\PY{n}{title}\PY{p}{(}\PY{l+s+sa}{f}\PY{l+s+s2}{\PYZdq{}}\PY{l+s+si}{\PYZob{}}\PY{n}{col}\PY{l+s+si}{\PYZcb{}}\PY{l+s+s2}{ vs }\PY{l+s+si}{\PYZob{}}\PY{n}{target}\PY{l+s+si}{\PYZcb{}}\PY{l+s+s2}{\PYZdq{}}\PY{p}{,} \PY{n}{fontsize}\PY{o}{=}\PY{l+m+mi}{14}\PY{p}{)}
    \PY{n}{plt}\PY{o}{.}\PY{n}{xlabel}\PY{p}{(}\PY{n}{col}\PY{p}{)}
    \PY{n}{plt}\PY{o}{.}\PY{n}{ylabel}\PY{p}{(}\PY{n}{target}\PY{p}{)}

    \PY{n}{plt}\PY{o}{.}\PY{n}{grid}\PY{p}{(}\PY{k+kc}{True}\PY{p}{)}

    \PY{n}{plt}\PY{o}{.}\PY{n}{show}\PY{p}{(}\PY{p}{)}
\end{Verbatim}
\end{tcolorbox}

    \begin{center}
    \adjustimage{max size={0.9\linewidth}{0.9\paperheight}}{output_66_0.png}
    \end{center}
    { \hspace*{\fill} \\}
    
    \begin{center}
    \adjustimage{max size={0.9\linewidth}{0.9\paperheight}}{output_66_1.png}
    \end{center}
    { \hspace*{\fill} \\}
    
    \begin{center}
    \adjustimage{max size={0.9\linewidth}{0.9\paperheight}}{output_66_2.png}
    \end{center}
    { \hspace*{\fill} \\}
    
    \begin{center}
    \adjustimage{max size={0.9\linewidth}{0.9\paperheight}}{output_66_3.png}
    \end{center}
    { \hspace*{\fill} \\}
    
    \begin{center}
    \adjustimage{max size={0.9\linewidth}{0.9\paperheight}}{output_66_4.png}
    \end{center}
    { \hspace*{\fill} \\}
    
    \begin{center}
    \adjustimage{max size={0.9\linewidth}{0.9\paperheight}}{output_66_5.png}
    \end{center}
    { \hspace*{\fill} \\}
    
    \begin{center}
    \adjustimage{max size={0.9\linewidth}{0.9\paperheight}}{output_66_6.png}
    \end{center}
    { \hspace*{\fill} \\}
    
    \begin{center}
    \adjustimage{max size={0.9\linewidth}{0.9\paperheight}}{output_66_7.png}
    \end{center}
    { \hspace*{\fill} \\}
    
    \begin{tcolorbox}[breakable, size=fbox, boxrule=1pt, pad at break*=1mm,colback=cellbackground, colframe=cellborder]
\prompt{In}{incolor}{67}{\boxspacing}
\begin{Verbatim}[commandchars=\\\{\}]
\PY{l+s+sd}{\PYZsq{}\PYZsq{}\PYZsq{}}
\PY{l+s+sd}{1] Longitude:}
\PY{l+s+sd}{The scatter plot shows that house values vary across different longitudes. Houses located near the western coastal regions generally have higher median house values than those located inland, indicating that location significantly influences house prices.}

\PY{l+s+sd}{2] Latitude:}
\PY{l+s+sd}{The latitude scatter plot indicates that house values change across different geographical regions. Higher house prices are observed in certain latitude ranges, suggesting that location plays an important role in determining property values.}

\PY{l+s+sd}{3] Housing Median Age:}
\PY{l+s+sd}{The relationship between housing median age and house value is weak to moderate. Although some older houses have higher values, there is no clear linear trend, indicating that house age alone is not a strong predictor of house prices.}

\PY{l+s+sd}{4] Total Rooms:}
\PY{l+s+sd}{The scatter plot shows a weak positive relationship between total rooms and median house value. Houses with more rooms tend to have higher prices, but the points are widely scattered, indicating that other factors also influence house values.}

\PY{l+s+sd}{5] Total Bedrooms:}
\PY{l+s+sd}{The total bedrooms feature shows a weak positive relationship with median house value. While larger houses generally have higher prices, the wide spread of data points suggests that the relationship is not very strong.}

\PY{l+s+sd}{6] Population:}
\PY{l+s+sd}{The population feature exhibits a very weak relationship with house values. Housing blocks with similar populations have a wide range of house prices, indicating that population alone is not a reliable predictor.}

\PY{l+s+sd}{7] Households:}
\PY{l+s+sd}{The households feature shows a slight positive trend with house prices. However, the large variation in data points suggests that this feature has limited predictive power when used independently.}

\PY{l+s+sd}{8] Median Income:}
\PY{l+s+sd}{The scatter plot reveals a strong positive relationship between median income and median house value. As median income increases, house prices generally increase as well. This makes median income the strongest predictor of house prices in the dataset.}

\PY{l+s+sd}{Overall Analysis:}

\PY{l+s+sd}{The scatter plots indicate that median\PYZus{}income has the strongest positive relationship with the target variable (**median\PYZus{}house\PYZus{}value**), making it the best feature for the Simple Linear Regression model. Features such as **population**, **households**, **total\PYZus{}rooms**, and **total\PYZus{}bedrooms** exhibit weaker relationships, while **longitude** and **latitude** highlight the importance of geographical location in determining house prices.}
\PY{l+s+sd}{\PYZsq{}\PYZsq{}\PYZsq{}}
\end{Verbatim}
\end{tcolorbox}

            \begin{tcolorbox}[breakable, size=fbox, boxrule=.5pt, pad at break*=1mm, opacityfill=0]
\prompt{Out}{outcolor}{67}{\boxspacing}
\begin{Verbatim}[commandchars=\\\{\}]
'\textbackslash{}n1] Longitude:\textbackslash{}nThe scatter plot shows that house values vary across different
longitudes. Houses located near the western coastal regions generally have
higher median house values than those located inland, indicating that location
significantly influences house prices.\textbackslash{}n\textbackslash{}n2] Latitude:\textbackslash{}nThe latitude scatter
plot indicates that house values change across different geographical regions.
Higher house prices are observed in certain latitude ranges, suggesting that
location plays an important role in determining property values.\textbackslash{}n\textbackslash{}n3] Housing
Median Age:\textbackslash{}nThe relationship between housing median age and house value is weak
to moderate. Although some older houses have higher values, there is no clear
linear trend, indicating that house age alone is not a strong predictor of house
prices.\textbackslash{}n\textbackslash{}n4] Total Rooms:\textbackslash{}nThe scatter plot shows a weak positive relationship
between total rooms and median house value. Houses with more rooms tend to have
higher prices, but the points are widely scattered, indicating that other
factors also influence house values.\textbackslash{}n\textbackslash{}n5] Total Bedrooms:\textbackslash{}nThe total bedrooms
feature shows a weak positive relationship with median house value. While larger
houses generally have higher prices, the wide spread of data points suggests
that the relationship is not very strong.\textbackslash{}n\textbackslash{}n6] Population:\textbackslash{}nThe population
feature exhibits a very weak relationship with house values. Housing blocks with
similar populations have a wide range of house prices, indicating that
population alone is not a reliable predictor.\textbackslash{}n\textbackslash{}n7] Households:\textbackslash{}nThe households
feature shows a slight positive trend with house prices. However, the large
variation in data points suggests that this feature has limited predictive power
when used independently.\textbackslash{}n\textbackslash{}n8] Median Income:\textbackslash{}nThe scatter plot reveals a strong
positive relationship between median income and median house value. As median
income increases, house prices generally increase as well. This makes median
income the strongest predictor of house prices in the dataset.\textbackslash{}n\textbackslash{}nOverall
Analysis:\textbackslash{}n\textbackslash{}nThe scatter plots indicate that median\_income has the strongest
positive relationship with the target variable (**median\_house\_value**), making
it the best feature for the Simple Linear Regression model. Features such as
**population**, **households**, **total\_rooms**, and **total\_bedrooms** exhibit
weaker relationships, while **longitude** and **latitude** highlight the
importance of geographical location in determining house prices.\textbackslash{}n'
\end{Verbatim}
\end{tcolorbox}
        
    \begin{tcolorbox}[breakable, size=fbox, boxrule=1pt, pad at break*=1mm,colback=cellbackground, colframe=cellborder]
\prompt{In}{incolor}{68}{\boxspacing}
\begin{Verbatim}[commandchars=\\\{\}]
\PY{c+c1}{\PYZsh{}correlation heatmap}
\PY{n}{plt}\PY{o}{.}\PY{n}{figure}\PY{p}{(}\PY{n}{figsize}\PY{o}{=}\PY{p}{(}\PY{l+m+mi}{10}\PY{p}{,}\PY{l+m+mi}{8}\PY{p}{)}\PY{p}{)}

\PY{n}{correlation} \PY{o}{=} \PY{n}{df}\PY{p}{[}\PY{n}{numerical\PYZus{}cols}\PY{p}{]}\PY{o}{.}\PY{n}{corr}\PY{p}{(}\PY{p}{)}

\PY{n}{sns}\PY{o}{.}\PY{n}{heatmap}\PY{p}{(}\PY{n}{correlation}\PY{p}{,}
            \PY{n}{annot}\PY{o}{=}\PY{k+kc}{True}\PY{p}{,}
            \PY{n}{cmap}\PY{o}{=}\PY{l+s+s2}{\PYZdq{}}\PY{l+s+s2}{coolwarm}\PY{l+s+s2}{\PYZdq{}}\PY{p}{,}
            \PY{n}{fmt}\PY{o}{=}\PY{l+s+s2}{\PYZdq{}}\PY{l+s+s2}{.2f}\PY{l+s+s2}{\PYZdq{}}\PY{p}{)}

\PY{n}{plt}\PY{o}{.}\PY{n}{title}\PY{p}{(}\PY{l+s+s2}{\PYZdq{}}\PY{l+s+s2}{Correlation Heatmap}\PY{l+s+s2}{\PYZdq{}}\PY{p}{,} \PY{n}{fontsize}\PY{o}{=}\PY{l+m+mi}{16}\PY{p}{)}

\PY{n}{plt}\PY{o}{.}\PY{n}{show}\PY{p}{(}\PY{p}{)}
\end{Verbatim}
\end{tcolorbox}

    \begin{center}
    \adjustimage{max size={0.9\linewidth}{0.9\paperheight}}{output_68_0.png}
    \end{center}
    { \hspace*{\fill} \\}
    
    \begin{tcolorbox}[breakable, size=fbox, boxrule=1pt, pad at break*=1mm,colback=cellbackground, colframe=cellborder]
\prompt{In}{incolor}{69}{\boxspacing}
\begin{Verbatim}[commandchars=\\\{\}]
\PY{l+s+sd}{\PYZsq{}\PYZsq{}\PYZsq{}}
\PY{l+s+sd}{Correlation Heatmap Analysis:}

\PY{l+s+sd}{1] The correlation heatmap shows the strength and direction of the relationship between all numerical features.}

\PY{l+s+sd}{2] Median Income has the strongest positive correlation with Median House Value. This indicates that areas with higher median incomes generally have higher house prices.}

\PY{l+s+sd}{3] Longitude has a negative correlation with Median House Value, suggesting that house prices tend to decrease as longitude increases. This reflects geographical differences in housing prices across California.}

\PY{l+s+sd}{4] Latitude also shows a negative correlation with Median House Value. Certain latitude regions have lower house values, while others, particularly coastal regions, have higher values.}

\PY{l+s+sd}{5] Total Rooms, Total Bedrooms, Population, and Households are strongly positively correlated with each other. This is expected because housing blocks with more rooms generally contain more bedrooms, households, and residents.}

\PY{l+s+sd}{6] Housing Median Age has a weak positive correlation with Median House Value, indicating that older houses may have slightly higher prices, but the relationship is not very strong.}

\PY{l+s+sd}{7] Population has a very weak correlation with Median House Value. This suggests that population alone is not a good predictor of house prices.}

\PY{l+s+sd}{8] The heatmap helps identify the most useful features for machine learning. Based on the correlation values, Median Income is selected as the best feature for the Simple Linear Regression model because it has the highest correlation with the target variable.}

\PY{l+s+sd}{Overall Analysis:}

\PY{l+s+sd}{The correlation heatmap provides a clear overview of relationships among the numerical features. It confirms that Median Income is the strongest predictor of Median House Value, while geographical features such as Longitude and Latitude also influence housing prices. Features with strong correlations are useful for building predictive models and selecting important variables for regression analysis.}
\PY{l+s+sd}{\PYZsq{}\PYZsq{}\PYZsq{}}
\end{Verbatim}
\end{tcolorbox}

            \begin{tcolorbox}[breakable, size=fbox, boxrule=.5pt, pad at break*=1mm, opacityfill=0]
\prompt{Out}{outcolor}{69}{\boxspacing}
\begin{Verbatim}[commandchars=\\\{\}]
'\textbackslash{}nCorrelation Heatmap Analysis:\textbackslash{}n\textbackslash{}n1] The correlation heatmap shows the
strength and direction of the relationship between all numerical features.\textbackslash{}n\textbackslash{}n2]
Median Income has the strongest positive correlation with Median House Value.
This indicates that areas with higher median incomes generally have higher house
prices.\textbackslash{}n\textbackslash{}n3] Longitude has a negative correlation with Median House Value,
suggesting that house prices tend to decrease as longitude increases. This
reflects geographical differences in housing prices across California.\textbackslash{}n\textbackslash{}n4]
Latitude also shows a negative correlation with Median House Value. Certain
latitude regions have lower house values, while others, particularly coastal
regions, have higher values.\textbackslash{}n\textbackslash{}n5] Total Rooms, Total Bedrooms, Population, and
Households are strongly positively correlated with each other. This is expected
because housing blocks with more rooms generally contain more bedrooms,
households, and residents.\textbackslash{}n\textbackslash{}n6] Housing Median Age has a weak positive
correlation with Median House Value, indicating that older houses may have
slightly higher prices, but the relationship is not very strong.\textbackslash{}n\textbackslash{}n7]
Population has a very weak correlation with Median House Value. This suggests
that population alone is not a good predictor of house prices.\textbackslash{}n\textbackslash{}n8] The heatmap
helps identify the most useful features for machine learning. Based on the
correlation values, Median Income is selected as the best feature for the Simple
Linear Regression model because it has the highest correlation with the target
variable.\textbackslash{}n\textbackslash{}nOverall Analysis:\textbackslash{}n\textbackslash{}nThe correlation heatmap provides a clear
overview of relationships among the numerical features. It confirms that Median
Income is the strongest predictor of Median House Value, while geographical
features such as Longitude and Latitude also influence housing prices. Features
with strong correlations are useful for building predictive models and selecting
important variables for regression analysis.\textbackslash{}n'
\end{Verbatim}
\end{tcolorbox}
        
    \begin{tcolorbox}[breakable, size=fbox, boxrule=1pt, pad at break*=1mm,colback=cellbackground, colframe=cellborder]
\prompt{In}{incolor}{70}{\boxspacing}
\begin{Verbatim}[commandchars=\\\{\}]
\PY{k+kn}{import}\PY{+w}{ }\PY{n+nn}{seaborn}\PY{+w}{ }\PY{k}{as}\PY{+w}{ }\PY{n+nn}{sns}
\PY{k+kn}{import}\PY{+w}{ }\PY{n+nn}{matplotlib}\PY{n+nn}{.}\PY{n+nn}{pyplot}\PY{+w}{ }\PY{k}{as}\PY{+w}{ }\PY{n+nn}{plt}

\PY{n}{important\PYZus{}features} \PY{o}{=} \PY{p}{[}
    \PY{l+s+s2}{\PYZdq{}}\PY{l+s+s2}{median\PYZus{}income}\PY{l+s+s2}{\PYZdq{}}\PY{p}{,}
    \PY{l+s+s2}{\PYZdq{}}\PY{l+s+s2}{housing\PYZus{}median\PYZus{}age}\PY{l+s+s2}{\PYZdq{}}\PY{p}{,}
    \PY{l+s+s2}{\PYZdq{}}\PY{l+s+s2}{total\PYZus{}rooms}\PY{l+s+s2}{\PYZdq{}}\PY{p}{,}
    \PY{l+s+s2}{\PYZdq{}}\PY{l+s+s2}{median\PYZus{}house\PYZus{}value}\PY{l+s+s2}{\PYZdq{}}
\PY{p}{]}

\PY{n}{sns}\PY{o}{.}\PY{n}{pairplot}\PY{p}{(}
    \PY{n}{df}\PY{p}{[}\PY{n}{important\PYZus{}features}\PY{p}{]}\PY{p}{,}
    \PY{n}{diag\PYZus{}kind}\PY{o}{=}\PY{l+s+s2}{\PYZdq{}}\PY{l+s+s2}{hist}\PY{l+s+s2}{\PYZdq{}}\PY{p}{,}
    \PY{n}{plot\PYZus{}kws}\PY{o}{=}\PY{p}{\PYZob{}}\PY{l+s+s2}{\PYZdq{}}\PY{l+s+s2}{alpha}\PY{l+s+s2}{\PYZdq{}}\PY{p}{:} \PY{l+m+mf}{0.5}\PY{p}{,} \PY{l+s+s2}{\PYZdq{}}\PY{l+s+s2}{color}\PY{l+s+s2}{\PYZdq{}}\PY{p}{:} \PY{l+s+s2}{\PYZdq{}}\PY{l+s+s2}{blue}\PY{l+s+s2}{\PYZdq{}}\PY{p}{\PYZcb{}}\PY{p}{,}
    \PY{n}{diag\PYZus{}kws}\PY{o}{=}\PY{p}{\PYZob{}}\PY{l+s+s2}{\PYZdq{}}\PY{l+s+s2}{color}\PY{l+s+s2}{\PYZdq{}}\PY{p}{:} \PY{l+s+s2}{\PYZdq{}}\PY{l+s+s2}{skyblue}\PY{l+s+s2}{\PYZdq{}}\PY{p}{\PYZcb{}}
\PY{p}{)}

\PY{n}{plt}\PY{o}{.}\PY{n}{suptitle}\PY{p}{(}\PY{l+s+s2}{\PYZdq{}}\PY{l+s+s2}{Pair Plot of Important Housing Features}\PY{l+s+s2}{\PYZdq{}}\PY{p}{,} \PY{n}{y}\PY{o}{=}\PY{l+m+mf}{1.02}\PY{p}{,} \PY{n}{fontsize}\PY{o}{=}\PY{l+m+mi}{16}\PY{p}{)}

\PY{n}{plt}\PY{o}{.}\PY{n}{show}\PY{p}{(}\PY{p}{)}
\end{Verbatim}
\end{tcolorbox}

    \begin{center}
    \adjustimage{max size={0.9\linewidth}{0.9\paperheight}}{output_70_0.png}
    \end{center}
    { \hspace*{\fill} \\}
    
    \begin{tcolorbox}[breakable, size=fbox, boxrule=1pt, pad at break*=1mm,colback=cellbackground, colframe=cellborder]
\prompt{In}{incolor}{71}{\boxspacing}
\begin{Verbatim}[commandchars=\\\{\}]
\PY{l+s+sd}{\PYZsq{}\PYZsq{}\PYZsq{}}
\PY{l+s+sd}{Pair Plot Analysis:}

\PY{l+s+sd}{1] Median Income:}
\PY{l+s+sd}{The pair plot shows a strong positive relationship between median income and median house value. }
\PY{l+s+sd}{As the median income increases, the house prices generally increase, confirming that median income is the strongest predictor of house prices.}

\PY{l+s+sd}{2] Housing Median Age:}
\PY{l+s+sd}{Housing median age exhibits a weak to moderate relationship with median house value.}
\PY{l+s+sd}{Although some older houses have higher values, the relationship is not strongly linear, indicating that age alone does not determine house prices.}

\PY{l+s+sd}{3] Total Rooms:}
\PY{l+s+sd}{The total rooms feature shows a slight positive relationship with median house value.}
\PY{l+s+sd}{However, the data points are widely scattered, suggesting that the number of rooms alone is not sufficient to accurately predict house prices.}

\PY{l+s+sd}{4] Median House Value:}
\PY{l+s+sd}{The target variable displays a positive association with median income but weaker relationships with housing median age and total rooms. }
\PY{l+s+sd}{The distribution also shows a concentration of values near the upper limit, indicating that house prices were capped in the dataset.}

\PY{l+s+sd}{5] Distribution of Individual Features:}
\PY{l+s+sd}{The diagonal histograms display the distribution of each feature. Median income and median house value are positively skewed, while total rooms contains several extreme values, indicating the presence of outliers.}

\PY{l+s+sd}{Overall Analysis:}

\PY{l+s+sd}{The pair plot provides a comprehensive visualization of the relationships among the selected features.}
\PY{l+s+sd}{It confirms that median income has the strongest positive relationship with median house value, making it the most suitable feature for Simple Linear Regression. }
\PY{l+s+sd}{Housing median age and total rooms show weaker relationships with the target variable. The pair plot also highlights the presence of skewed distributions and outliers, which are important considerations during model development.}
\PY{l+s+sd}{\PYZsq{}\PYZsq{}\PYZsq{}}
\end{Verbatim}
\end{tcolorbox}

            \begin{tcolorbox}[breakable, size=fbox, boxrule=.5pt, pad at break*=1mm, opacityfill=0]
\prompt{Out}{outcolor}{71}{\boxspacing}
\begin{Verbatim}[commandchars=\\\{\}]
'\textbackslash{}nPair Plot Analysis:\textbackslash{}n\textbackslash{}n1] Median Income:\textbackslash{}nThe pair plot shows a strong
positive relationship between median income and median house value. \textbackslash{}nAs the
median income increases, the house prices generally increase, confirming that
median income is the strongest predictor of house prices.\textbackslash{}n\textbackslash{}n2] Housing Median
Age:\textbackslash{}nHousing median age exhibits a weak to moderate relationship with median
house value.\textbackslash{}nAlthough some older houses have higher values, the relationship is
not strongly linear, indicating that age alone does not determine house
prices.\textbackslash{}n\textbackslash{}n3] Total Rooms:\textbackslash{}nThe total rooms feature shows a slight positive
relationship with median house value.\textbackslash{}nHowever, the data points are widely
scattered, suggesting that the number of rooms alone is not sufficient to
accurately predict house prices.\textbackslash{}n\textbackslash{}n4] Median House Value:\textbackslash{}nThe target variable
displays a positive association with median income but weaker relationships with
housing median age and total rooms. \textbackslash{}nThe distribution also shows a
concentration of values near the upper limit, indicating that house prices were
capped in the dataset.\textbackslash{}n\textbackslash{}n5] Distribution of Individual Features:\textbackslash{}nThe diagonal
histograms display the distribution of each feature. Median income and median
house value are positively skewed, while total rooms contains several extreme
values, indicating the presence of outliers.\textbackslash{}n\textbackslash{}nOverall Analysis:\textbackslash{}n\textbackslash{}nThe pair
plot provides a comprehensive visualization of the relationships among the
selected features.\textbackslash{}nIt confirms that median income has the strongest positive
relationship with median house value, making it the most suitable feature for
Simple Linear Regression. \textbackslash{}nHousing median age and total rooms show weaker
relationships with the target variable. The pair plot also highlights the
presence of skewed distributions and outliers, which are important
considerations during model development.\textbackslash{}n'
\end{Verbatim}
\end{tcolorbox}
        
    \begin{tcolorbox}[breakable, size=fbox, boxrule=1pt, pad at break*=1mm,colback=cellbackground, colframe=cellborder]
\prompt{In}{incolor}{72}{\boxspacing}
\begin{Verbatim}[commandchars=\\\{\}]
                    \PY{c+c1}{\PYZsh{} PART \PYZhy{}4(Insights \PYZam{} Observations )}
\end{Verbatim}
\end{tcolorbox}

    \begin{tcolorbox}[breakable, size=fbox, boxrule=1pt, pad at break*=1mm,colback=cellbackground, colframe=cellborder]
\prompt{In}{incolor}{73}{\boxspacing}
\begin{Verbatim}[commandchars=\\\{\}]
\PY{l+s+sd}{\PYZsq{}\PYZsq{}\PYZsq{}}
\PY{l+s+sd}{Observation 1}
\PY{l+s+sd}{Median Income has the strongest positive correlation with Median House Value. }
\PY{l+s+sd}{This indicates that areas with higher household incomes generally have higher house prices, making median income the best predictor for the regression model.}

\PY{l+s+sd}{Observation 2}
\PY{l+s+sd}{The distribution of Median House Value is positively (right) skewed, with most houses priced in the low to middle range}
\PY{l+s+sd}{The distribution also shows an upper limit (around \PYZdl{}500,000), indicating that house prices were capped in the dataset}

\PY{l+s+sd}{Observation 3}
\PY{l+s+sd}{Features such as Total Rooms, Total Bedrooms, Population, and Households are highly right\PYZhy{}skewed}
\PY{l+s+sd}{Most housing blocks have relatively small values, while a few blocks contain exceptionally large values}

\PY{l+s+sd}{Observation 4}
\PY{l+s+sd}{The box plots reveal many outliers in Total Rooms, Total Bedrooms, Population, and Households}
\PY{l+s+sd}{These outliers likely represent densely populated housing blocks or unusually large residential areas rather than data entry errors}
\PY{l+s+sd}{Observation 5}
\PY{l+s+sd}{Longitude and Latitude both influence house prices. }
\PY{l+s+sd}{Houses located in coastal regions generally have higher median house values compared to inland regions, highlighting the importance of geographical location.}

\PY{l+s+sd}{Observation 6}
\PY{l+s+sd}{Population has a very weak correlation with Median House Value. This suggests that simply having a larger population does not necessarily result in higher housing prices.}

\PY{l+s+sd}{Observation 7}
\PY{l+s+sd}{Housing Median Age has a weak positive relationship with house prices}
\PY{l+s+sd}{While some older houses have higher values, house age alone is not a strong predictor of housing prices}
\PY{l+s+sd}{    }
\PY{l+s+sd}{Observation 8}
\PY{l+s+sd}{The pair plot confirms that Median Income has the clearest positive linear relationship with Median House Value. Other features show weaker and more scattered relationships with the target variable.}

\PY{l+s+sd}{Observation 9}
\PY{l+s+sd}{Several numerical features, particularly Total Rooms, Total Bedrooms, Population, and Households, are strongly correlated with one another }
\PY{l+s+sd}{his suggests possible multicollinearity, which should be considered when building a Multiple Linear Regression model}
\PY{l+s+sd}{    }
\PY{l+s+sd}{Observation 10}
\PY{l+s+sd}{the exploratory data analysis indicates that economic and geographical factors have a greater influence on house prices than demographic factors}
\PY{l+s+sd}{    Therefore, features such as Median Income, Latitude, and Longitude are expected to improve prediction accuracy in regression models\PYZsq{}\PYZsq{}\PYZsq{}}
\end{Verbatim}
\end{tcolorbox}

            \begin{tcolorbox}[breakable, size=fbox, boxrule=.5pt, pad at break*=1mm, opacityfill=0]
\prompt{Out}{outcolor}{73}{\boxspacing}
\begin{Verbatim}[commandchars=\\\{\}]
'\textbackslash{}nObservation 1\textbackslash{}nMedian Income has the strongest positive correlation with
Median House Value. \textbackslash{}nThis indicates that areas with higher household incomes
generally have higher house prices, making median income the best predictor for
the regression model.\textbackslash{}n\textbackslash{}nObservation 2\textbackslash{}nThe distribution of Median House Value
is positively (right) skewed, with most houses priced in the low to middle
range\textbackslash{}nThe distribution also shows an upper limit (around \$500,000), indicating
that house prices were capped in the dataset\textbackslash{}n\textbackslash{}nObservation 3\textbackslash{}nFeatures such as
Total Rooms, Total Bedrooms, Population, and Households are highly right-
skewed\textbackslash{}nMost housing blocks have relatively small values, while a few blocks
contain exceptionally large values\textbackslash{}n\textbackslash{}nObservation 4\textbackslash{}nThe box plots reveal many
outliers in Total Rooms, Total Bedrooms, Population, and Households\textbackslash{}nThese
outliers likely represent densely populated housing blocks or unusually large
residential areas rather than data entry errors\textbackslash{}nObservation 5\textbackslash{}nLongitude and
Latitude both influence house prices. \textbackslash{}nHouses located in coastal regions
generally have higher median house values compared to inland regions,
highlighting the importance of geographical location.\textbackslash{}n\textbackslash{}nObservation
6\textbackslash{}nPopulation has a very weak correlation with Median House Value. This suggests
that simply having a larger population does not necessarily result in higher
housing prices.\textbackslash{}n\textbackslash{}nObservation 7\textbackslash{}nHousing Median Age has a weak positive
relationship with house prices\textbackslash{}nWhile some older houses have higher values,
house age alone is not a strong predictor of housing prices\textbackslash{}n\textbackslash{}nObservation
8\textbackslash{}nThe pair plot confirms that Median Income has the clearest positive linear
relationship with Median House Value. Other features show weaker and more
scattered relationships with the target variable.\textbackslash{}n\textbackslash{}nObservation 9\textbackslash{}nSeveral
numerical features, particularly Total Rooms, Total Bedrooms, Population, and
Households, are strongly correlated with one another \textbackslash{}nhis suggests possible
multicollinearity, which should be considered when building a Multiple Linear
Regression model\textbackslash{}n\textbackslash{}nObservation 10\textbackslash{}nthe exploratory data analysis indicates that
economic and geographical factors have a greater influence on house prices than
demographic factors\textbackslash{}n    Therefore, features such as Median Income, Latitude,
and Longitude are expected to improve prediction accuracy in regression models'
\end{Verbatim}
\end{tcolorbox}
        
    \begin{tcolorbox}[breakable, size=fbox, boxrule=1pt, pad at break*=1mm,colback=cellbackground, colframe=cellborder]
\prompt{In}{incolor}{74}{\boxspacing}
\begin{Verbatim}[commandchars=\\\{\}]
                          \PY{c+c1}{\PYZsh{}PART\PYZhy{}5(Simple Linear Regression )}
\end{Verbatim}
\end{tcolorbox}

    \begin{tcolorbox}[breakable, size=fbox, boxrule=1pt, pad at break*=1mm,colback=cellbackground, colframe=cellborder]
\prompt{In}{incolor}{75}{\boxspacing}
\begin{Verbatim}[commandchars=\\\{\}]
\PY{k+kn}{from}\PY{+w}{ }\PY{n+nn}{sklearn}\PY{n+nn}{.}\PY{n+nn}{model\PYZus{}selection}\PY{+w}{ }\PY{k+kn}{import} \PY{n}{train\PYZus{}test\PYZus{}split}
\PY{k+kn}{from}\PY{+w}{ }\PY{n+nn}{sklearn}\PY{n+nn}{.}\PY{n+nn}{linear\PYZus{}model}\PY{+w}{ }\PY{k+kn}{import} \PY{n}{LinearRegression}

\PY{k+kn}{from}\PY{+w}{ }\PY{n+nn}{sklearn}\PY{n+nn}{.}\PY{n+nn}{metrics}\PY{+w}{ }\PY{k+kn}{import} \PY{p}{(}
    \PY{n}{r2\PYZus{}score}\PY{p}{,}
    \PY{n}{mean\PYZus{}absolute\PYZus{}error}\PY{p}{,}
    \PY{n}{root\PYZus{}mean\PYZus{}squared\PYZus{}error}
\PY{p}{)}
\end{Verbatim}
\end{tcolorbox}

    \begin{tcolorbox}[breakable, size=fbox, boxrule=1pt, pad at break*=1mm,colback=cellbackground, colframe=cellborder]
\prompt{In}{incolor}{76}{\boxspacing}
\begin{Verbatim}[commandchars=\\\{\}]
\PY{k+kn}{from}\PY{+w}{ }\PY{n+nn}{sklearn}\PY{n+nn}{.}\PY{n+nn}{metrics}\PY{+w}{ }\PY{k+kn}{import} \PY{n}{mean\PYZus{}squared\PYZus{}error}
\end{Verbatim}
\end{tcolorbox}

    \begin{tcolorbox}[breakable, size=fbox, boxrule=1pt, pad at break*=1mm,colback=cellbackground, colframe=cellborder]
\prompt{In}{incolor}{77}{\boxspacing}
\begin{Verbatim}[commandchars=\\\{\}]
\PY{n}{X} \PY{o}{=} \PY{n}{df}\PY{p}{[}\PY{p}{[}\PY{l+s+s2}{\PYZdq{}}\PY{l+s+s2}{median\PYZus{}income}\PY{l+s+s2}{\PYZdq{}}\PY{p}{]}\PY{p}{]} \PY{c+c1}{\PYZsh{}feature}
\PY{n}{y} \PY{o}{=} \PY{n}{df}\PY{p}{[}\PY{l+s+s2}{\PYZdq{}}\PY{l+s+s2}{median\PYZus{}house\PYZus{}value}\PY{l+s+s2}{\PYZdq{}}\PY{p}{]} \PY{c+c1}{\PYZsh{}target}
\end{Verbatim}
\end{tcolorbox}

    \begin{tcolorbox}[breakable, size=fbox, boxrule=1pt, pad at break*=1mm,colback=cellbackground, colframe=cellborder]
\prompt{In}{incolor}{78}{\boxspacing}
\begin{Verbatim}[commandchars=\\\{\}]
\PY{l+s+sd}{\PYZsq{}\PYZsq{}\PYZsq{}From the correlation heatmap, median\PYZus{}income showed the highest positive correlation with median\PYZus{}house\PYZus{}value}
\PY{l+s+sd}{This indicates that neighborhoods with higher median incomes generally have higher house prices.}
\PY{l+s+sd}{Therefore, median\PYZus{}income was selected as the independent variable for the Simple Linear Regression model\PYZsq{}\PYZsq{}\PYZsq{}}
\end{Verbatim}
\end{tcolorbox}

            \begin{tcolorbox}[breakable, size=fbox, boxrule=.5pt, pad at break*=1mm, opacityfill=0]
\prompt{Out}{outcolor}{78}{\boxspacing}
\begin{Verbatim}[commandchars=\\\{\}]
'From the correlation heatmap, median\_income showed the highest positive
correlation with median\_house\_value\textbackslash{}nThis indicates that neighborhoods with
higher median incomes generally have higher house prices.\textbackslash{}nTherefore,
median\_income was selected as the independent variable for the Simple Linear
Regression model'
\end{Verbatim}
\end{tcolorbox}
        
    \begin{tcolorbox}[breakable, size=fbox, boxrule=1pt, pad at break*=1mm,colback=cellbackground, colframe=cellborder]
\prompt{In}{incolor}{79}{\boxspacing}
\begin{Verbatim}[commandchars=\\\{\}]
\PY{n}{X\PYZus{}train}\PY{p}{,} \PY{n}{X\PYZus{}test}\PY{p}{,} \PY{n}{y\PYZus{}train}\PY{p}{,} \PY{n}{y\PYZus{}test} \PY{o}{=} \PY{n}{train\PYZus{}test\PYZus{}split}\PY{p}{(}    \PY{c+c1}{\PYZsh{}split data}
    \PY{n}{X}\PY{p}{,}
    \PY{n}{y}\PY{p}{,}
    \PY{n}{test\PYZus{}size}\PY{o}{=}\PY{l+m+mf}{0.20}\PY{p}{,}
    \PY{n}{random\PYZus{}state}\PY{o}{=}\PY{l+m+mi}{42}
\PY{p}{)}
\end{Verbatim}
\end{tcolorbox}

    \begin{tcolorbox}[breakable, size=fbox, boxrule=1pt, pad at break*=1mm,colback=cellbackground, colframe=cellborder]
\prompt{In}{incolor}{80}{\boxspacing}
\begin{Verbatim}[commandchars=\\\{\}]
\PY{n+nb}{print}\PY{p}{(}\PY{l+s+s2}{\PYZdq{}}\PY{l+s+s2}{Training samples :}\PY{l+s+s2}{\PYZdq{}}\PY{p}{,} \PY{n+nb}{len}\PY{p}{(}\PY{n}{X\PYZus{}train}\PY{p}{)}\PY{p}{)}
\PY{n+nb}{print}\PY{p}{(}\PY{l+s+s2}{\PYZdq{}}\PY{l+s+s2}{Testing samples :}\PY{l+s+s2}{\PYZdq{}}\PY{p}{,} \PY{n+nb}{len}\PY{p}{(}\PY{n}{X\PYZus{}test}\PY{p}{)}\PY{p}{)}
\end{Verbatim}
\end{tcolorbox}

    \begin{Verbatim}[commandchars=\\\{\}]
Training samples : 16512
Testing samples : 4128
    \end{Verbatim}

    \begin{tcolorbox}[breakable, size=fbox, boxrule=1pt, pad at break*=1mm,colback=cellbackground, colframe=cellborder]
\prompt{In}{incolor}{81}{\boxspacing}
\begin{Verbatim}[commandchars=\\\{\}]
\PY{c+c1}{\PYZsh{}train model}
\PY{n}{model} \PY{o}{=} \PY{n}{LinearRegression}\PY{p}{(}\PY{p}{)}

\PY{n}{model}\PY{o}{.}\PY{n}{fit}\PY{p}{(}\PY{n}{X\PYZus{}train}\PY{p}{,} \PY{n}{y\PYZus{}train}\PY{p}{)}
\end{Verbatim}
\end{tcolorbox}

            \begin{tcolorbox}[breakable, size=fbox, boxrule=.5pt, pad at break*=1mm, opacityfill=0]
\prompt{Out}{outcolor}{81}{\boxspacing}
\begin{Verbatim}[commandchars=\\\{\}]
LinearRegression()
\end{Verbatim}
\end{tcolorbox}
        
    \begin{tcolorbox}[breakable, size=fbox, boxrule=1pt, pad at break*=1mm,colback=cellbackground, colframe=cellborder]
\prompt{In}{incolor}{82}{\boxspacing}
\begin{Verbatim}[commandchars=\\\{\}]
\PY{c+c1}{\PYZsh{}diplay slope and intercept}
\PY{n+nb}{print}\PY{p}{(}\PY{l+s+s2}{\PYZdq{}}\PY{l+s+s2}{Slope (Coefficient):}\PY{l+s+s2}{\PYZdq{}}\PY{p}{,} \PY{n}{model}\PY{o}{.}\PY{n}{coef\PYZus{}}\PY{p}{[}\PY{l+m+mi}{0}\PY{p}{]}\PY{p}{)}

\PY{n+nb}{print}\PY{p}{(}\PY{l+s+s2}{\PYZdq{}}\PY{l+s+s2}{Intercept:}\PY{l+s+s2}{\PYZdq{}}\PY{p}{,} \PY{n}{model}\PY{o}{.}\PY{n}{intercept\PYZus{}}\PY{p}{)}
\end{Verbatim}
\end{tcolorbox}

    \begin{Verbatim}[commandchars=\\\{\}]
Slope (Coefficient): 41933.84939381274
Intercept: 44459.729169078666
    \end{Verbatim}

    \begin{tcolorbox}[breakable, size=fbox, boxrule=1pt, pad at break*=1mm,colback=cellbackground, colframe=cellborder]
\prompt{In}{incolor}{83}{\boxspacing}
\begin{Verbatim}[commandchars=\\\{\}]
\PY{c+c1}{\PYZsh{}prediction}
\PY{n}{y\PYZus{}pred} \PY{o}{=} \PY{n}{model}\PY{o}{.}\PY{n}{predict}\PY{p}{(}\PY{n}{X\PYZus{}test}\PY{p}{)}
\end{Verbatim}
\end{tcolorbox}

    \begin{tcolorbox}[breakable, size=fbox, boxrule=1pt, pad at break*=1mm,colback=cellbackground, colframe=cellborder]
\prompt{In}{incolor}{84}{\boxspacing}
\begin{Verbatim}[commandchars=\\\{\}]
\PY{c+c1}{\PYZsh{}scatter plot with regression}
\PY{n}{plt}\PY{o}{.}\PY{n}{figure}\PY{p}{(}\PY{n}{figsize}\PY{o}{=}\PY{p}{(}\PY{l+m+mi}{8}\PY{p}{,}\PY{l+m+mi}{6}\PY{p}{)}\PY{p}{)}

\PY{n}{plt}\PY{o}{.}\PY{n}{scatter}\PY{p}{(}\PY{n}{X\PYZus{}test}\PY{p}{,} \PY{n}{y\PYZus{}test}\PY{p}{,} \PY{n}{color}\PY{o}{=}\PY{l+s+s2}{\PYZdq{}}\PY{l+s+s2}{blue}\PY{l+s+s2}{\PYZdq{}}\PY{p}{,} \PY{n}{alpha}\PY{o}{=}\PY{l+m+mf}{0.5}\PY{p}{,} \PY{n}{label}\PY{o}{=}\PY{l+s+s2}{\PYZdq{}}\PY{l+s+s2}{Actual Data}\PY{l+s+s2}{\PYZdq{}}\PY{p}{)}
\PY{n}{plt}\PY{o}{.}\PY{n}{plot}\PY{p}{(}\PY{n}{X\PYZus{}test}\PY{p}{,} \PY{n}{y\PYZus{}pred}\PY{p}{,} \PY{n}{color}\PY{o}{=}\PY{l+s+s2}{\PYZdq{}}\PY{l+s+s2}{red}\PY{l+s+s2}{\PYZdq{}}\PY{p}{,} \PY{n}{linewidth}\PY{o}{=}\PY{l+m+mi}{3}\PY{p}{,} \PY{n}{label}\PY{o}{=}\PY{l+s+s2}{\PYZdq{}}\PY{l+s+s2}{Regression Line}\PY{l+s+s2}{\PYZdq{}}\PY{p}{)}

\PY{n}{plt}\PY{o}{.}\PY{n}{title}\PY{p}{(}\PY{l+s+s2}{\PYZdq{}}\PY{l+s+s2}{Simple Linear Regression}\PY{l+s+s2}{\PYZdq{}}\PY{p}{)}
\PY{n}{plt}\PY{o}{.}\PY{n}{xlabel}\PY{p}{(}\PY{l+s+s2}{\PYZdq{}}\PY{l+s+s2}{Median Income}\PY{l+s+s2}{\PYZdq{}}\PY{p}{)}
\PY{n}{plt}\PY{o}{.}\PY{n}{ylabel}\PY{p}{(}\PY{l+s+s2}{\PYZdq{}}\PY{l+s+s2}{Median House Value}\PY{l+s+s2}{\PYZdq{}}\PY{p}{)}
\PY{n}{plt}\PY{o}{.}\PY{n}{legend}\PY{p}{(}\PY{p}{)}

\PY{n}{plt}\PY{o}{.}\PY{n}{show}\PY{p}{(}\PY{p}{)}
\end{Verbatim}
\end{tcolorbox}

    \begin{center}
    \adjustimage{max size={0.9\linewidth}{0.9\paperheight}}{output_84_0.png}
    \end{center}
    { \hspace*{\fill} \\}
    
    \begin{tcolorbox}[breakable, size=fbox, boxrule=1pt, pad at break*=1mm,colback=cellbackground, colframe=cellborder]
\prompt{In}{incolor}{85}{\boxspacing}
\begin{Verbatim}[commandchars=\\\{\}]
\PY{l+s+sd}{\PYZsq{}\PYZsq{}\PYZsq{}The regression line shows a positive relationship between median income and median house value}

\PY{l+s+sd}{Although the data points are scattered around the line, an increasing trend is clearly visible. }
\PY{l+s+sd}{This indicates that higher median income is generally associated with higher house prices\PYZsq{}\PYZsq{}\PYZsq{}}
\end{Verbatim}
\end{tcolorbox}

            \begin{tcolorbox}[breakable, size=fbox, boxrule=.5pt, pad at break*=1mm, opacityfill=0]
\prompt{Out}{outcolor}{85}{\boxspacing}
\begin{Verbatim}[commandchars=\\\{\}]
'The regression line shows a positive relationship between median income and
median house value\textbackslash{}n\textbackslash{}nAlthough the data points are scattered around the line, an
increasing trend is clearly visible. \textbackslash{}nThis indicates that higher median income
is generally associated with higher house prices'
\end{Verbatim}
\end{tcolorbox}
        
    \begin{tcolorbox}[breakable, size=fbox, boxrule=1pt, pad at break*=1mm,colback=cellbackground, colframe=cellborder]
\prompt{In}{incolor}{86}{\boxspacing}
\begin{Verbatim}[commandchars=\\\{\}]
\PY{c+c1}{\PYZsh{}Residual plot}
\PY{n}{residuals} \PY{o}{=} \PY{n}{y\PYZus{}test} \PY{o}{\PYZhy{}} \PY{n}{y\PYZus{}pred}

\PY{n}{plt}\PY{o}{.}\PY{n}{figure}\PY{p}{(}\PY{n}{figsize}\PY{o}{=}\PY{p}{(}\PY{l+m+mi}{8}\PY{p}{,}\PY{l+m+mi}{6}\PY{p}{)}\PY{p}{)}
\PY{n}{plt}\PY{o}{.}\PY{n}{scatter}\PY{p}{(}\PY{n}{y\PYZus{}pred}\PY{p}{,} \PY{n}{residuals}\PY{p}{,} \PY{n}{color}\PY{o}{=}\PY{l+s+s2}{\PYZdq{}}\PY{l+s+s2}{green}\PY{l+s+s2}{\PYZdq{}}\PY{p}{,} \PY{n}{alpha}\PY{o}{=}\PY{l+m+mf}{0.5}\PY{p}{)}

\PY{n}{plt}\PY{o}{.}\PY{n}{axhline}\PY{p}{(}\PY{n}{y}\PY{o}{=}\PY{l+m+mi}{0}\PY{p}{,} \PY{n}{color}\PY{o}{=}\PY{l+s+s2}{\PYZdq{}}\PY{l+s+s2}{red}\PY{l+s+s2}{\PYZdq{}}\PY{p}{,} \PY{n}{linestyle}\PY{o}{=}\PY{l+s+s2}{\PYZdq{}}\PY{l+s+s2}{\PYZhy{}\PYZhy{}}\PY{l+s+s2}{\PYZdq{}}\PY{p}{)}

\PY{n}{plt}\PY{o}{.}\PY{n}{title}\PY{p}{(}\PY{l+s+s2}{\PYZdq{}}\PY{l+s+s2}{Residual Plot}\PY{l+s+s2}{\PYZdq{}}\PY{p}{)}
\PY{n}{plt}\PY{o}{.}\PY{n}{xlabel}\PY{p}{(}\PY{l+s+s2}{\PYZdq{}}\PY{l+s+s2}{Predicted House Value}\PY{l+s+s2}{\PYZdq{}}\PY{p}{)}
\PY{n}{plt}\PY{o}{.}\PY{n}{ylabel}\PY{p}{(}\PY{l+s+s2}{\PYZdq{}}\PY{l+s+s2}{Residuals}\PY{l+s+s2}{\PYZdq{}}\PY{p}{)}

\PY{n}{plt}\PY{o}{.}\PY{n}{show}\PY{p}{(}\PY{p}{)}
\end{Verbatim}
\end{tcolorbox}

    \begin{center}
    \adjustimage{max size={0.9\linewidth}{0.9\paperheight}}{output_86_0.png}
    \end{center}
    { \hspace*{\fill} \\}
    
    \begin{tcolorbox}[breakable, size=fbox, boxrule=1pt, pad at break*=1mm,colback=cellbackground, colframe=cellborder]
\prompt{In}{incolor}{87}{\boxspacing}
\begin{Verbatim}[commandchars=\\\{\}]
\PY{l+s+sd}{\PYZsq{}\PYZsq{}\PYZsq{}Most residuals are randomly distributed around zero, indicating that the model captures the general trend of the data. }
\PY{l+s+sd}{However, some larger residuals suggest that a single predictor cannot explain all variations in house prices\PYZsq{}\PYZsq{}\PYZsq{}}
\end{Verbatim}
\end{tcolorbox}

            \begin{tcolorbox}[breakable, size=fbox, boxrule=.5pt, pad at break*=1mm, opacityfill=0]
\prompt{Out}{outcolor}{87}{\boxspacing}
\begin{Verbatim}[commandchars=\\\{\}]
'Most residuals are randomly distributed around zero, indicating that the model
captures the general trend of the data. \textbackslash{}nHowever, some larger residuals suggest
that a single predictor cannot explain all variations in house prices'
\end{Verbatim}
\end{tcolorbox}
        
    \begin{tcolorbox}[breakable, size=fbox, boxrule=1pt, pad at break*=1mm,colback=cellbackground, colframe=cellborder]
\prompt{In}{incolor}{88}{\boxspacing}
\begin{Verbatim}[commandchars=\\\{\}]
\PY{c+c1}{\PYZsh{}evalution metrics}
\PY{n}{r2} \PY{o}{=} \PY{n}{r2\PYZus{}score}\PY{p}{(}\PY{n}{y\PYZus{}test}\PY{p}{,} \PY{n}{y\PYZus{}pred}\PY{p}{)}
\PY{n}{mae} \PY{o}{=} \PY{n}{mean\PYZus{}absolute\PYZus{}error}\PY{p}{(}\PY{n}{y\PYZus{}test}\PY{p}{,} \PY{n}{y\PYZus{}pred}\PY{p}{)}
\PY{n}{rmse} \PY{o}{=} \PY{n}{np}\PY{o}{.}\PY{n}{sqrt}\PY{p}{(}\PY{n}{mean\PYZus{}squared\PYZus{}error}\PY{p}{(}\PY{n}{y\PYZus{}test}\PY{p}{,} \PY{n}{y\PYZus{}pred}\PY{p}{)}\PY{p}{)}
\PY{n+nb}{print}\PY{p}{(}\PY{l+s+s2}{\PYZdq{}}\PY{l+s+s2}{R² Score:}\PY{l+s+s2}{\PYZdq{}}\PY{p}{,} \PY{n}{r2}\PY{p}{)}
\PY{n+nb}{print}\PY{p}{(}\PY{l+s+s2}{\PYZdq{}}\PY{l+s+s2}{MAE:}\PY{l+s+s2}{\PYZdq{}}\PY{p}{,} \PY{n}{mae}\PY{p}{)}
\PY{n+nb}{print}\PY{p}{(}\PY{l+s+s2}{\PYZdq{}}\PY{l+s+s2}{RMSE:}\PY{l+s+s2}{\PYZdq{}}\PY{p}{,} \PY{n}{rmse}\PY{p}{)}
\end{Verbatim}
\end{tcolorbox}

    \begin{Verbatim}[commandchars=\\\{\}]
R² Score: 0.45885918903846656
MAE: 62990.86530093761
RMSE: 84209.01241414454
    \end{Verbatim}

    \begin{tcolorbox}[breakable, size=fbox, boxrule=1pt, pad at break*=1mm,colback=cellbackground, colframe=cellborder]
\prompt{In}{incolor}{89}{\boxspacing}
\begin{Verbatim}[commandchars=\\\{\}]
                 \PY{c+c1}{\PYZsh{}PART 6(Multiple Linear Regression)}
\end{Verbatim}
\end{tcolorbox}

    \begin{tcolorbox}[breakable, size=fbox, boxrule=1pt, pad at break*=1mm,colback=cellbackground, colframe=cellborder]
\prompt{In}{incolor}{90}{\boxspacing}
\begin{Verbatim}[commandchars=\\\{\}]
\PY{n}{features} \PY{o}{=} \PY{p}{[}
    \PY{l+s+s2}{\PYZdq{}}\PY{l+s+s2}{median\PYZus{}income}\PY{l+s+s2}{\PYZdq{}}\PY{p}{,}
    \PY{l+s+s2}{\PYZdq{}}\PY{l+s+s2}{housing\PYZus{}median\PYZus{}age}\PY{l+s+s2}{\PYZdq{}}\PY{p}{,}
    \PY{l+s+s2}{\PYZdq{}}\PY{l+s+s2}{total\PYZus{}rooms}\PY{l+s+s2}{\PYZdq{}}\PY{p}{,}
    \PY{l+s+s2}{\PYZdq{}}\PY{l+s+s2}{latitude}\PY{l+s+s2}{\PYZdq{}}\PY{p}{,}
    \PY{l+s+s2}{\PYZdq{}}\PY{l+s+s2}{longitude}\PY{l+s+s2}{\PYZdq{}}
\PY{p}{]}

\PY{n}{X} \PY{o}{=} \PY{n}{df}\PY{p}{[}\PY{n}{features}\PY{p}{]}
\PY{n}{y} \PY{o}{=} \PY{n}{df}\PY{p}{[}\PY{l+s+s2}{\PYZdq{}}\PY{l+s+s2}{median\PYZus{}house\PYZus{}value}\PY{l+s+s2}{\PYZdq{}}\PY{p}{]}
\end{Verbatim}
\end{tcolorbox}

    \begin{tcolorbox}[breakable, size=fbox, boxrule=1pt, pad at break*=1mm,colback=cellbackground, colframe=cellborder]
\prompt{In}{incolor}{91}{\boxspacing}
\begin{Verbatim}[commandchars=\\\{\}]
\PY{l+s+sd}{\PYZsq{}\PYZsq{}\PYZsq{}Feature Selection}

\PY{l+s+sd}{The following features were selected because they have meaningful relationships}
\PY{l+s+sd}{with the target variable (median\PYZus{}house\PYZus{}value).}

\PY{l+s+sd}{• median\PYZus{}income – Highest positive correlation with house value.}
\PY{l+s+sd}{• housing\PYZus{}median\PYZus{}age – Older neighborhoods may have higher property values.}
\PY{l+s+sd}{• total\PYZus{}rooms – Represents the size of houses.}
\PY{l+s+sd}{• latitude – Captures geographical location.}
\PY{l+s+sd}{• longitude – Captures geographical location.}

\PY{l+s+sd}{These features combine economic, structural, and geographical information,}
\PY{l+s+sd}{which are important factors affecting house prices.\PYZsq{}\PYZsq{}\PYZsq{}}
\end{Verbatim}
\end{tcolorbox}

            \begin{tcolorbox}[breakable, size=fbox, boxrule=.5pt, pad at break*=1mm, opacityfill=0]
\prompt{Out}{outcolor}{91}{\boxspacing}
\begin{Verbatim}[commandchars=\\\{\}]
'Feature Selection\textbackslash{}n\textbackslash{}nThe following features were selected because they have
meaningful relationships\textbackslash{}nwith the target variable (median\_house\_value).\textbackslash{}n\textbackslash{}n•
median\_income – Highest positive correlation with house value.\textbackslash{}n•
housing\_median\_age – Older neighborhoods may have higher property values.\textbackslash{}n•
total\_rooms – Represents the size of houses.\textbackslash{}n• latitude – Captures geographical
location.\textbackslash{}n• longitude – Captures geographical location.\textbackslash{}n\textbackslash{}nThese features
combine economic, structural, and geographical information,\textbackslash{}nwhich are important
factors affecting house prices.'
\end{Verbatim}
\end{tcolorbox}
        
    \begin{tcolorbox}[breakable, size=fbox, boxrule=1pt, pad at break*=1mm,colback=cellbackground, colframe=cellborder]
\prompt{In}{incolor}{92}{\boxspacing}
\begin{Verbatim}[commandchars=\\\{\}]
\PY{c+c1}{\PYZsh{}\PYZsh{} Conclusion}

\PY{l+s+sd}{\PYZsq{}\PYZsq{}\PYZsq{}The Simple Linear Regression model successfully predicts median house values using median income.}

\PY{l+s+sd}{Median income was selected because it had the highest positive correlation with the target variable.}

\PY{l+s+sd}{The model demonstrates a clear positive relationship between income and house prices. However, because only one feature is used, the model cannot capture all factors influencing house values. Therefore, its prediction accuracy is moderate.}

\PY{l+s+sd}{A Multiple Linear Regression model using additional features such as longitude, latitude, housing median age, and total rooms is expected to achieve better performance\PYZsq{}\PYZsq{}\PYZsq{}}
\end{Verbatim}
\end{tcolorbox}

            \begin{tcolorbox}[breakable, size=fbox, boxrule=.5pt, pad at break*=1mm, opacityfill=0]
\prompt{Out}{outcolor}{92}{\boxspacing}
\begin{Verbatim}[commandchars=\\\{\}]
'The Simple Linear Regression model successfully predicts median house values
using median income.\textbackslash{}n\textbackslash{}nMedian income was selected because it had the highest
positive correlation with the target variable.\textbackslash{}n\textbackslash{}nThe model demonstrates a clear
positive relationship between income and house prices. However, because only one
feature is used, the model cannot capture all factors influencing house values.
Therefore, its prediction accuracy is moderate.\textbackslash{}n\textbackslash{}nA Multiple Linear Regression
model using additional features such as longitude, latitude, housing median age,
and total rooms is expected to achieve better performance'
\end{Verbatim}
\end{tcolorbox}
        
    \begin{tcolorbox}[breakable, size=fbox, boxrule=1pt, pad at break*=1mm,colback=cellbackground, colframe=cellborder]
\prompt{In}{incolor}{93}{\boxspacing}
\begin{Verbatim}[commandchars=\\\{\}]
\PY{k+kn}{from}\PY{+w}{ }\PY{n+nn}{sklearn}\PY{n+nn}{.}\PY{n+nn}{model\PYZus{}selection}\PY{+w}{ }\PY{k+kn}{import} \PY{n}{train\PYZus{}test\PYZus{}split}

\PY{n}{X\PYZus{}train}\PY{p}{,} \PY{n}{X\PYZus{}test}\PY{p}{,} \PY{n}{y\PYZus{}train}\PY{p}{,} \PY{n}{y\PYZus{}test} \PY{o}{=} \PY{n}{train\PYZus{}test\PYZus{}split}\PY{p}{(}
    \PY{n}{X}\PY{p}{,}
    \PY{n}{y}\PY{p}{,}
    \PY{n}{test\PYZus{}size}\PY{o}{=}\PY{l+m+mf}{0.2}\PY{p}{,}
    \PY{n}{random\PYZus{}state}\PY{o}{=}\PY{l+m+mi}{42}
\PY{p}{)}
\end{Verbatim}
\end{tcolorbox}

    \begin{tcolorbox}[breakable, size=fbox, boxrule=1pt, pad at break*=1mm,colback=cellbackground, colframe=cellborder]
\prompt{In}{incolor}{94}{\boxspacing}
\begin{Verbatim}[commandchars=\\\{\}]
\PY{k+kn}{from}\PY{+w}{ }\PY{n+nn}{sklearn}\PY{n+nn}{.}\PY{n+nn}{linear\PYZus{}model}\PY{+w}{ }\PY{k+kn}{import} \PY{n}{LinearRegression}

\PY{n}{model} \PY{o}{=} \PY{n}{LinearRegression}\PY{p}{(}\PY{p}{)}

\PY{n}{model}\PY{o}{.}\PY{n}{fit}\PY{p}{(}\PY{n}{X\PYZus{}train}\PY{p}{,} \PY{n}{y\PYZus{}train}\PY{p}{)}
\end{Verbatim}
\end{tcolorbox}

            \begin{tcolorbox}[breakable, size=fbox, boxrule=.5pt, pad at break*=1mm, opacityfill=0]
\prompt{Out}{outcolor}{94}{\boxspacing}
\begin{Verbatim}[commandchars=\\\{\}]
LinearRegression()
\end{Verbatim}
\end{tcolorbox}
        
    \begin{tcolorbox}[breakable, size=fbox, boxrule=1pt, pad at break*=1mm,colback=cellbackground, colframe=cellborder]
\prompt{In}{incolor}{95}{\boxspacing}
\begin{Verbatim}[commandchars=\\\{\}]
\PY{n+nb}{print}\PY{p}{(}\PY{n}{model}\PY{o}{.}\PY{n}{coef\PYZus{}}\PY{o}{.}\PY{n}{shape}\PY{p}{)}
\end{Verbatim}
\end{tcolorbox}

    \begin{Verbatim}[commandchars=\\\{\}]
(5,)
    \end{Verbatim}

    \begin{tcolorbox}[breakable, size=fbox, boxrule=1pt, pad at break*=1mm,colback=cellbackground, colframe=cellborder]
\prompt{In}{incolor}{96}{\boxspacing}
\begin{Verbatim}[commandchars=\\\{\}]
\PY{n}{coefficients} \PY{o}{=} \PY{n}{pd}\PY{o}{.}\PY{n}{DataFrame}\PY{p}{(}\PY{p}{\PYZob{}}
    \PY{l+s+s2}{\PYZdq{}}\PY{l+s+s2}{Feature}\PY{l+s+s2}{\PYZdq{}}\PY{p}{:} \PY{n}{X}\PY{o}{.}\PY{n}{columns}\PY{p}{,}
    \PY{l+s+s2}{\PYZdq{}}\PY{l+s+s2}{Coefficient}\PY{l+s+s2}{\PYZdq{}}\PY{p}{:} \PY{n}{model}\PY{o}{.}\PY{n}{coef\PYZus{}}
\PY{p}{\PYZcb{}}\PY{p}{)}

\PY{n+nb}{print}\PY{p}{(}\PY{n}{coefficients}\PY{p}{)}
\end{Verbatim}
\end{tcolorbox}

    \begin{Verbatim}[commandchars=\\\{\}]
              Feature   Coefficient
0       median\_income  37535.980826
1  housing\_median\_age   1167.869663
2         total\_rooms      3.383466
3            latitude -43887.440531
4           longitude -44682.959059
    \end{Verbatim}

    \begin{tcolorbox}[breakable, size=fbox, boxrule=1pt, pad at break*=1mm,colback=cellbackground, colframe=cellborder]
\prompt{In}{incolor}{97}{\boxspacing}
\begin{Verbatim}[commandchars=\\\{\}]
\PY{n+nb}{print}\PY{p}{(}\PY{l+s+s2}{\PYZdq{}}\PY{l+s+s2}{X shape:}\PY{l+s+s2}{\PYZdq{}}\PY{p}{,} \PY{n}{X}\PY{o}{.}\PY{n}{shape}\PY{p}{)}
\PY{n+nb}{print}\PY{p}{(}\PY{l+s+s2}{\PYZdq{}}\PY{l+s+s2}{Number of features:}\PY{l+s+s2}{\PYZdq{}}\PY{p}{,} \PY{n+nb}{len}\PY{p}{(}\PY{n}{X}\PY{o}{.}\PY{n}{columns}\PY{p}{)}\PY{p}{)}
\PY{n+nb}{print}\PY{p}{(}\PY{l+s+s2}{\PYZdq{}}\PY{l+s+s2}{Feature names:}\PY{l+s+s2}{\PYZdq{}}\PY{p}{,} \PY{n}{X}\PY{o}{.}\PY{n}{columns}\PY{o}{.}\PY{n}{tolist}\PY{p}{(}\PY{p}{)}\PY{p}{)}

\PY{n+nb}{print}\PY{p}{(}\PY{l+s+s2}{\PYZdq{}}\PY{l+s+s2}{Coefficient shape:}\PY{l+s+s2}{\PYZdq{}}\PY{p}{,} \PY{n}{model}\PY{o}{.}\PY{n}{coef\PYZus{}}\PY{o}{.}\PY{n}{shape}\PY{p}{)}
\PY{n+nb}{print}\PY{p}{(}\PY{l+s+s2}{\PYZdq{}}\PY{l+s+s2}{Coefficients:}\PY{l+s+s2}{\PYZdq{}}\PY{p}{,} \PY{n}{model}\PY{o}{.}\PY{n}{coef\PYZus{}}\PY{p}{)}

\PY{n+nb}{print}\PY{p}{(}\PY{l+s+s2}{\PYZdq{}}\PY{l+s+s2}{Intercept:}\PY{l+s+s2}{\PYZdq{}}\PY{p}{,} \PY{n}{model}\PY{o}{.}\PY{n}{intercept\PYZus{}}\PY{p}{)}
\end{Verbatim}
\end{tcolorbox}

    \begin{Verbatim}[commandchars=\\\{\}]
X shape: (20640, 5)
Number of features: 5
Feature names: ['median\_income', 'housing\_median\_age', 'total\_rooms',
'latitude', 'longitude']
Coefficient shape: (5,)
Coefficients: [ 3.75359808e+04  1.16786966e+03  3.38346589e+00 -4.38874405e+04
 -4.46829591e+04]
Intercept: -3759827.101363706
    \end{Verbatim}

    \begin{tcolorbox}[breakable, size=fbox, boxrule=1pt, pad at break*=1mm,colback=cellbackground, colframe=cellborder]
\prompt{In}{incolor}{98}{\boxspacing}
\begin{Verbatim}[commandchars=\\\{\}]
\PY{n+nb}{print}\PY{p}{(}\PY{n}{model}\PY{o}{.}\PY{n}{coef\PYZus{}}\PY{p}{)}
\end{Verbatim}
\end{tcolorbox}

    \begin{Verbatim}[commandchars=\\\{\}]
[ 3.75359808e+04  1.16786966e+03  3.38346589e+00 -4.38874405e+04
 -4.46829591e+04]
    \end{Verbatim}

    \begin{tcolorbox}[breakable, size=fbox, boxrule=1pt, pad at break*=1mm,colback=cellbackground, colframe=cellborder]
\prompt{In}{incolor}{99}{\boxspacing}
\begin{Verbatim}[commandchars=\\\{\}]
\PY{c+c1}{\PYZsh{}Evalution metrics}
\PY{k+kn}{from}\PY{+w}{ }\PY{n+nn}{sklearn}\PY{n+nn}{.}\PY{n+nn}{metrics}\PY{+w}{ }\PY{k+kn}{import} \PY{n}{r2\PYZus{}score}\PY{p}{,} \PY{n}{mean\PYZus{}absolute\PYZus{}error}\PY{p}{,} \PY{n}{mean\PYZus{}squared\PYZus{}error}
\PY{k+kn}{import}\PY{+w}{ }\PY{n+nn}{numpy}\PY{+w}{ }\PY{k}{as}\PY{+w}{ }\PY{n+nn}{np}

\PY{n}{y\PYZus{}pred} \PY{o}{=} \PY{n}{model}\PY{o}{.}\PY{n}{predict}\PY{p}{(}\PY{n}{X\PYZus{}test}\PY{p}{)}

\PY{n}{r2} \PY{o}{=} \PY{n}{r2\PYZus{}score}\PY{p}{(}\PY{n}{y\PYZus{}test}\PY{p}{,} \PY{n}{y\PYZus{}pred}\PY{p}{)}
\PY{n}{mae} \PY{o}{=} \PY{n}{mean\PYZus{}absolute\PYZus{}error}\PY{p}{(}\PY{n}{y\PYZus{}test}\PY{p}{,} \PY{n}{y\PYZus{}pred}\PY{p}{)}
\PY{n}{rmse} \PY{o}{=} \PY{n}{np}\PY{o}{.}\PY{n}{sqrt}\PY{p}{(}\PY{n}{mean\PYZus{}squared\PYZus{}error}\PY{p}{(}\PY{n}{y\PYZus{}test}\PY{p}{,} \PY{n}{y\PYZus{}pred}\PY{p}{)}\PY{p}{)}

\PY{n+nb}{print}\PY{p}{(}\PY{l+s+s2}{\PYZdq{}}\PY{l+s+s2}{R² Score:}\PY{l+s+s2}{\PYZdq{}}\PY{p}{,} \PY{n}{r2}\PY{p}{)}
\PY{n+nb}{print}\PY{p}{(}\PY{l+s+s2}{\PYZdq{}}\PY{l+s+s2}{MAE:}\PY{l+s+s2}{\PYZdq{}}\PY{p}{,} \PY{n}{mae}\PY{p}{)}
\PY{n+nb}{print}\PY{p}{(}\PY{l+s+s2}{\PYZdq{}}\PY{l+s+s2}{RMSE:}\PY{l+s+s2}{\PYZdq{}}\PY{p}{,} \PY{n}{rmse}\PY{p}{)}
\end{Verbatim}
\end{tcolorbox}

    \begin{Verbatim}[commandchars=\\\{\}]
R² Score: 0.584448731061169
MAE: 54230.7444767925
RMSE: 73793.09681378353
    \end{Verbatim}

    \begin{tcolorbox}[breakable, size=fbox, boxrule=1pt, pad at break*=1mm,colback=cellbackground, colframe=cellborder]
\prompt{In}{incolor}{100}{\boxspacing}
\begin{Verbatim}[commandchars=\\\{\}]
\PY{k+kn}{from}\PY{+w}{ }\PY{n+nn}{sklearn}\PY{n+nn}{.}\PY{n+nn}{model\PYZus{}selection}\PY{+w}{ }\PY{k+kn}{import} \PY{n}{train\PYZus{}test\PYZus{}split}
\PY{k+kn}{from}\PY{+w}{ }\PY{n+nn}{sklearn}\PY{n+nn}{.}\PY{n+nn}{linear\PYZus{}model}\PY{+w}{ }\PY{k+kn}{import} \PY{n}{LinearRegression}
\PY{k+kn}{from}\PY{+w}{ }\PY{n+nn}{sklearn}\PY{n+nn}{.}\PY{n+nn}{metrics}\PY{+w}{ }\PY{k+kn}{import} \PY{n}{r2\PYZus{}score}\PY{p}{,} \PY{n}{mean\PYZus{}absolute\PYZus{}error}\PY{p}{,} \PY{n}{mean\PYZus{}squared\PYZus{}error}
\PY{k+kn}{import}\PY{+w}{ }\PY{n+nn}{numpy}\PY{+w}{ }\PY{k}{as}\PY{+w}{ }\PY{n+nn}{np}

\PY{c+c1}{\PYZsh{} Simple Linear Regression}
\PY{n}{X\PYZus{}simple} \PY{o}{=} \PY{n}{df}\PY{p}{[}\PY{p}{[}\PY{l+s+s2}{\PYZdq{}}\PY{l+s+s2}{median\PYZus{}income}\PY{l+s+s2}{\PYZdq{}}\PY{p}{]}\PY{p}{]}
\PY{n}{y} \PY{o}{=} \PY{n}{df}\PY{p}{[}\PY{l+s+s2}{\PYZdq{}}\PY{l+s+s2}{median\PYZus{}house\PYZus{}value}\PY{l+s+s2}{\PYZdq{}}\PY{p}{]}

\PY{n}{X\PYZus{}train\PYZus{}simple}\PY{p}{,} \PY{n}{X\PYZus{}test\PYZus{}simple}\PY{p}{,} \PY{n}{y\PYZus{}train\PYZus{}simple}\PY{p}{,} \PY{n}{y\PYZus{}test\PYZus{}simple} \PY{o}{=} \PY{n}{train\PYZus{}test\PYZus{}split}\PY{p}{(}
    \PY{n}{X\PYZus{}simple}\PY{p}{,} \PY{n}{y}\PY{p}{,} \PY{n}{test\PYZus{}size}\PY{o}{=}\PY{l+m+mf}{0.2}\PY{p}{,} \PY{n}{random\PYZus{}state}\PY{o}{=}\PY{l+m+mi}{42}
\PY{p}{)}

\PY{n}{simple\PYZus{}model} \PY{o}{=} \PY{n}{LinearRegression}\PY{p}{(}\PY{p}{)}
\PY{n}{simple\PYZus{}model}\PY{o}{.}\PY{n}{fit}\PY{p}{(}\PY{n}{X\PYZus{}train\PYZus{}simple}\PY{p}{,} \PY{n}{y\PYZus{}train\PYZus{}simple}\PY{p}{)}

\PY{n}{y\PYZus{}pred\PYZus{}simple} \PY{o}{=} \PY{n}{simple\PYZus{}model}\PY{o}{.}\PY{n}{predict}\PY{p}{(}\PY{n}{X\PYZus{}test\PYZus{}simple}\PY{p}{)}

\PY{c+c1}{\PYZsh{} Save metrics}
\PY{n}{simple\PYZus{}r2} \PY{o}{=} \PY{n}{r2\PYZus{}score}\PY{p}{(}\PY{n}{y\PYZus{}test\PYZus{}simple}\PY{p}{,} \PY{n}{y\PYZus{}pred\PYZus{}simple}\PY{p}{)}
\PY{n}{simple\PYZus{}mae} \PY{o}{=} \PY{n}{mean\PYZus{}absolute\PYZus{}error}\PY{p}{(}\PY{n}{y\PYZus{}test\PYZus{}simple}\PY{p}{,} \PY{n}{y\PYZus{}pred\PYZus{}simple}\PY{p}{)}
\PY{n}{simple\PYZus{}rmse} \PY{o}{=} \PY{n}{np}\PY{o}{.}\PY{n}{sqrt}\PY{p}{(}\PY{n}{mean\PYZus{}squared\PYZus{}error}\PY{p}{(}\PY{n}{y\PYZus{}test\PYZus{}simple}\PY{p}{,} \PY{n}{y\PYZus{}pred\PYZus{}simple}\PY{p}{)}\PY{p}{)}

\PY{n+nb}{print}\PY{p}{(}\PY{l+s+s2}{\PYZdq{}}\PY{l+s+s2}{Simple R²:}\PY{l+s+s2}{\PYZdq{}}\PY{p}{,} \PY{n}{simple\PYZus{}r2}\PY{p}{)}
\PY{n+nb}{print}\PY{p}{(}\PY{l+s+s2}{\PYZdq{}}\PY{l+s+s2}{Simple MAE:}\PY{l+s+s2}{\PYZdq{}}\PY{p}{,} \PY{n}{simple\PYZus{}mae}\PY{p}{)}
\PY{n+nb}{print}\PY{p}{(}\PY{l+s+s2}{\PYZdq{}}\PY{l+s+s2}{Simple RMSE:}\PY{l+s+s2}{\PYZdq{}}\PY{p}{,} \PY{n}{simple\PYZus{}rmse}\PY{p}{)}
\end{Verbatim}
\end{tcolorbox}

    \begin{Verbatim}[commandchars=\\\{\}]
Simple R²: 0.45885918903846656
Simple MAE: 62990.86530093761
Simple RMSE: 84209.01241414454
    \end{Verbatim}

    \begin{tcolorbox}[breakable, size=fbox, boxrule=1pt, pad at break*=1mm,colback=cellbackground, colframe=cellborder]
\prompt{In}{incolor}{101}{\boxspacing}
\begin{Verbatim}[commandchars=\\\{\}]
\PY{n}{multiple\PYZus{}r2} \PY{o}{=} \PY{n}{r2}
\PY{n}{multiple\PYZus{}mae} \PY{o}{=} \PY{n}{mae}
\PY{n}{multiple\PYZus{}rmse} \PY{o}{=} \PY{n}{rmse}
\end{Verbatim}
\end{tcolorbox}

    \begin{tcolorbox}[breakable, size=fbox, boxrule=1pt, pad at break*=1mm,colback=cellbackground, colframe=cellborder]
\prompt{In}{incolor}{102}{\boxspacing}
\begin{Verbatim}[commandchars=\\\{\}]
\PY{n}{y\PYZus{}pred} \PY{o}{=} \PY{n}{model}\PY{o}{.}\PY{n}{predict}\PY{p}{(}\PY{n}{X\PYZus{}test}\PY{p}{)}

\PY{n}{multiple\PYZus{}r2} \PY{o}{=} \PY{n}{r2\PYZus{}score}\PY{p}{(}\PY{n}{y\PYZus{}test}\PY{p}{,} \PY{n}{y\PYZus{}pred}\PY{p}{)}
\PY{n}{multiple\PYZus{}mae} \PY{o}{=} \PY{n}{mean\PYZus{}absolute\PYZus{}error}\PY{p}{(}\PY{n}{y\PYZus{}test}\PY{p}{,} \PY{n}{y\PYZus{}pred}\PY{p}{)}
\PY{n}{multiple\PYZus{}rmse} \PY{o}{=} \PY{n}{np}\PY{o}{.}\PY{n}{sqrt}\PY{p}{(}\PY{n}{mean\PYZus{}squared\PYZus{}error}\PY{p}{(}\PY{n}{y\PYZus{}test}\PY{p}{,} \PY{n}{y\PYZus{}pred}\PY{p}{)}\PY{p}{)}
\end{Verbatim}
\end{tcolorbox}

    \begin{tcolorbox}[breakable, size=fbox, boxrule=1pt, pad at break*=1mm,colback=cellbackground, colframe=cellborder]
\prompt{In}{incolor}{103}{\boxspacing}
\begin{Verbatim}[commandchars=\\\{\}]
\PY{n}{comparison} \PY{o}{=} \PY{n}{pd}\PY{o}{.}\PY{n}{DataFrame}\PY{p}{(}\PY{p}{\PYZob{}}
    \PY{l+s+s2}{\PYZdq{}}\PY{l+s+s2}{Metric}\PY{l+s+s2}{\PYZdq{}}\PY{p}{:} \PY{p}{[}\PY{l+s+s2}{\PYZdq{}}\PY{l+s+s2}{R² Score}\PY{l+s+s2}{\PYZdq{}}\PY{p}{,} \PY{l+s+s2}{\PYZdq{}}\PY{l+s+s2}{MAE}\PY{l+s+s2}{\PYZdq{}}\PY{p}{,} \PY{l+s+s2}{\PYZdq{}}\PY{l+s+s2}{RMSE}\PY{l+s+s2}{\PYZdq{}}\PY{p}{]}\PY{p}{,}
    \PY{l+s+s2}{\PYZdq{}}\PY{l+s+s2}{Simple LR}\PY{l+s+s2}{\PYZdq{}}\PY{p}{:} \PY{p}{[}\PY{n}{simple\PYZus{}r2}\PY{p}{,} \PY{n}{simple\PYZus{}mae}\PY{p}{,} \PY{n}{simple\PYZus{}rmse}\PY{p}{]}\PY{p}{,}
    \PY{l+s+s2}{\PYZdq{}}\PY{l+s+s2}{Multiple LR}\PY{l+s+s2}{\PYZdq{}}\PY{p}{:} \PY{p}{[}\PY{n}{multiple\PYZus{}r2}\PY{p}{,} \PY{n}{multiple\PYZus{}mae}\PY{p}{,} \PY{n}{multiple\PYZus{}rmse}\PY{p}{]}
\PY{p}{\PYZcb{}}\PY{p}{)}

\PY{n+nb}{print}\PY{p}{(}\PY{n}{comparison}\PY{p}{)}
\end{Verbatim}
\end{tcolorbox}

    \begin{Verbatim}[commandchars=\\\{\}]
     Metric     Simple LR   Multiple LR
0  R² Score      0.458859      0.584449
1       MAE  62990.865301  54230.744477
2      RMSE  84209.012414  73793.096814
    \end{Verbatim}

    \begin{tcolorbox}[breakable, size=fbox, boxrule=1pt, pad at break*=1mm,colback=cellbackground, colframe=cellborder]
\prompt{In}{incolor}{104}{\boxspacing}
\begin{Verbatim}[commandchars=\\\{\}]
\PY{n}{comparison}\PY{p}{[}\PY{l+s+s2}{\PYZdq{}}\PY{l+s+s2}{Improvement}\PY{l+s+s2}{\PYZdq{}}\PY{p}{]} \PY{o}{=} \PY{p}{[}
    \PY{l+s+sa}{f}\PY{l+s+s2}{\PYZdq{}}\PY{l+s+si}{\PYZob{}}\PY{p}{(}\PY{n}{multiple\PYZus{}r2}\PY{o}{\PYZhy{}}\PY{n}{simple\PYZus{}r2}\PY{p}{)}\PY{o}{*}\PY{l+m+mi}{100}\PY{l+s+si}{:}\PY{l+s+s2}{.2f}\PY{l+s+si}{\PYZcb{}}\PY{l+s+s2}{\PYZpc{}}\PY{l+s+s2}{\PYZdq{}}\PY{p}{,}
    \PY{l+s+sa}{f}\PY{l+s+s2}{\PYZdq{}}\PY{l+s+si}{\PYZob{}}\PY{p}{(}\PY{p}{(}\PY{n}{simple\PYZus{}mae}\PY{o}{\PYZhy{}}\PY{n}{multiple\PYZus{}mae}\PY{p}{)}\PY{o}{/}\PY{n}{simple\PYZus{}mae}\PY{p}{)}\PY{o}{*}\PY{l+m+mi}{100}\PY{l+s+si}{:}\PY{l+s+s2}{.2f}\PY{l+s+si}{\PYZcb{}}\PY{l+s+s2}{\PYZpc{} lower}\PY{l+s+s2}{\PYZdq{}}\PY{p}{,}
    \PY{l+s+sa}{f}\PY{l+s+s2}{\PYZdq{}}\PY{l+s+si}{\PYZob{}}\PY{p}{(}\PY{p}{(}\PY{n}{simple\PYZus{}rmse}\PY{o}{\PYZhy{}}\PY{n}{multiple\PYZus{}rmse}\PY{p}{)}\PY{o}{/}\PY{n}{simple\PYZus{}rmse}\PY{p}{)}\PY{o}{*}\PY{l+m+mi}{100}\PY{l+s+si}{:}\PY{l+s+s2}{.2f}\PY{l+s+si}{\PYZcb{}}\PY{l+s+s2}{\PYZpc{} lower}\PY{l+s+s2}{\PYZdq{}}
\PY{p}{]}

\PY{n}{comparison}
\end{Verbatim}
\end{tcolorbox}

            \begin{tcolorbox}[breakable, size=fbox, boxrule=.5pt, pad at break*=1mm, opacityfill=0]
\prompt{Out}{outcolor}{104}{\boxspacing}
\begin{Verbatim}[commandchars=\\\{\}]
     Metric     Simple LR   Multiple LR   Improvement
0  R² Score      0.458859      0.584449        12.56\%
1       MAE  62990.865301  54230.744477  13.91\% lower
2      RMSE  84209.012414  73793.096814  12.37\% lower
\end{Verbatim}
\end{tcolorbox}
        
    \begin{tcolorbox}[breakable, size=fbox, boxrule=1pt, pad at break*=1mm,colback=cellbackground, colframe=cellborder]
\prompt{In}{incolor}{105}{\boxspacing}
\begin{Verbatim}[commandchars=\\\{\}]
\PY{l+s+sd}{\PYZsq{}\PYZsq{}\PYZsq{}The Multiple Linear Regression model performs better than the Simple Linear Regression model}

\PY{l+s+sd}{\PYZhy{} The R² Score increased, indicating that the model explains a larger percentage of the variation in house prices}
\PY{l+s+sd}{\PYZhy{} Both MAE and RMSE decreased, meaning the prediction errors became smaller}
\PY{l+s+sd}{\PYZhy{} Adding more relevant features improved the model\PYZsq{}s predictive performance}

\PY{l+s+sd}{Overall, Multiple Linear Regression provides more accurate house price predictions than Simple Linear Regression\PYZsq{}\PYZsq{}\PYZsq{}}
\end{Verbatim}
\end{tcolorbox}

            \begin{tcolorbox}[breakable, size=fbox, boxrule=.5pt, pad at break*=1mm, opacityfill=0]
\prompt{Out}{outcolor}{105}{\boxspacing}
\begin{Verbatim}[commandchars=\\\{\}]
"The Multiple Linear Regression model performs better than the Simple Linear
Regression model\textbackslash{}n\textbackslash{}n- The R² Score increased, indicating that the model explains
a larger percentage of the variation in house prices\textbackslash{}n- Both MAE and RMSE
decreased, meaning the prediction errors became smaller\textbackslash{}n- Adding more relevant
features improved the model's predictive performance\textbackslash{}n\textbackslash{}nOverall, Multiple Linear
Regression provides more accurate house price predictions than Simple Linear
Regression"
\end{Verbatim}
\end{tcolorbox}
        
    \begin{tcolorbox}[breakable, size=fbox, boxrule=1pt, pad at break*=1mm,colback=cellbackground, colframe=cellborder]
\prompt{In}{incolor}{106}{\boxspacing}
\begin{Verbatim}[commandchars=\\\{\}]
\PY{l+s+sd}{\PYZsq{}\PYZsq{}\PYZsq{}\PYZsh{}\PYZsh{} Feature Importance}

\PY{l+s+sd}{Based on the regression coefficients:}

\PY{l+s+sd}{1. Median Income}
\PY{l+s+sd}{2. Longitude}
\PY{l+s+sd}{3. Latitude}

\PY{l+s+sd}{These three features have the largest impact on predicting house prices}

\PY{l+s+sd}{Median Income contributes positively, while Latitude and Longitude capture the geographical influence on property values\PYZsq{}\PYZsq{}\PYZsq{}}
\end{Verbatim}
\end{tcolorbox}

            \begin{tcolorbox}[breakable, size=fbox, boxrule=.5pt, pad at break*=1mm, opacityfill=0]
\prompt{Out}{outcolor}{106}{\boxspacing}
\begin{Verbatim}[commandchars=\\\{\}]
'\#\# Feature Importance\textbackslash{}n\textbackslash{}nBased on the regression coefficients:\textbackslash{}n\textbackslash{}n1. Median
Income\textbackslash{}n2. Longitude\textbackslash{}n3. Latitude\textbackslash{}n\textbackslash{}nThese three features have the largest
impact on predicting house prices\textbackslash{}n\textbackslash{}nMedian Income contributes positively, while
Latitude and Longitude capture the geographical influence on property values'
\end{Verbatim}
\end{tcolorbox}
        
    \begin{tcolorbox}[breakable, size=fbox, boxrule=1pt, pad at break*=1mm,colback=cellbackground, colframe=cellborder]
\prompt{In}{incolor}{107}{\boxspacing}
\begin{Verbatim}[commandchars=\\\{\}]
\PY{c+c1}{\PYZsh{}\PYZsh{} If Performance Did Not Improve}

\PY{l+s+sd}{\PYZsq{}\PYZsq{}\PYZsq{}If the Multiple Linear Regression model does not perform better than the Simple Linear Regression model, the possible reasons are:}

\PY{l+s+sd}{\PYZhy{} The selected features may not be very useful for predicting house prices}
\PY{l+s+sd}{\PYZhy{} Some features may contain similar information, so adding them does not improve the model much}
\PY{l+s+sd}{\PYZhy{} House prices may depend on complex factors that cannot be explained by a simple linear relationship}
\PY{l+s+sd}{\PYZhy{} The dataset may contain outliers or unusual values that affect the model\PYZsq{}s accuracy}
\PY{l+s+sd}{\PYZhy{} Using a more advanced model, such as Random Forest, may produce better predictions because it can capture more complex patterns in the data\PYZsq{}\PYZsq{}\PYZsq{}}
\end{Verbatim}
\end{tcolorbox}

            \begin{tcolorbox}[breakable, size=fbox, boxrule=.5pt, pad at break*=1mm, opacityfill=0]
\prompt{Out}{outcolor}{107}{\boxspacing}
\begin{Verbatim}[commandchars=\\\{\}]
"If the Multiple Linear Regression model does not perform better than the Simple
Linear Regression model, the possible reasons are:\textbackslash{}n\textbackslash{}n- The selected features
may not be very useful for predicting house prices\textbackslash{}n- Some features may contain
similar information, so adding them does not improve the model much\textbackslash{}n- House
prices may depend on complex factors that cannot be explained by a simple linear
relationship\textbackslash{}n- The dataset may contain outliers or unusual values that affect
the model's accuracy\textbackslash{}n- Using a more advanced model, such as Random Forest, may
produce better predictions because it can capture more complex patterns in the
data"
\end{Verbatim}
\end{tcolorbox}
        
    \begin{tcolorbox}[breakable, size=fbox, boxrule=1pt, pad at break*=1mm,colback=cellbackground, colframe=cellborder]
\prompt{In}{incolor}{108}{\boxspacing}
\begin{Verbatim}[commandchars=\\\{\}]
                     \PY{c+c1}{\PYZsh{}PART\PYZhy{}7(Visualization Portfolio )}
\end{Verbatim}
\end{tcolorbox}

    \begin{tcolorbox}[breakable, size=fbox, boxrule=1pt, pad at break*=1mm,colback=cellbackground, colframe=cellborder]
\prompt{In}{incolor}{109}{\boxspacing}
\begin{Verbatim}[commandchars=\\\{\}]
\PY{c+c1}{\PYZsh{}histogram}
\PY{k+kn}{import}\PY{+w}{ }\PY{n+nn}{matplotlib}\PY{n+nn}{.}\PY{n+nn}{pyplot}\PY{+w}{ }\PY{k}{as}\PY{+w}{ }\PY{n+nn}{plt}

\PY{n}{features} \PY{o}{=} \PY{p}{[}\PY{l+s+s2}{\PYZdq{}}\PY{l+s+s2}{median\PYZus{}house\PYZus{}value}\PY{l+s+s2}{\PYZdq{}}\PY{p}{,} \PY{l+s+s2}{\PYZdq{}}\PY{l+s+s2}{median\PYZus{}income}\PY{l+s+s2}{\PYZdq{}}\PY{p}{,} \PY{l+s+s2}{\PYZdq{}}\PY{l+s+s2}{housing\PYZus{}median\PYZus{}age}\PY{l+s+s2}{\PYZdq{}}\PY{p}{]}

\PY{k}{for} \PY{n}{feature} \PY{o+ow}{in} \PY{n}{features}\PY{p}{:}
    \PY{n}{plt}\PY{o}{.}\PY{n}{figure}\PY{p}{(}\PY{n}{figsize}\PY{o}{=}\PY{p}{(}\PY{l+m+mi}{8}\PY{p}{,}\PY{l+m+mi}{5}\PY{p}{)}\PY{p}{)}

    \PY{n}{plt}\PY{o}{.}\PY{n}{hist}\PY{p}{(}\PY{n}{df}\PY{p}{[}\PY{n}{feature}\PY{p}{]}\PY{p}{,}
             \PY{n}{bins}\PY{o}{=}\PY{l+m+mi}{30}\PY{p}{,}
             \PY{n}{color}\PY{o}{=}\PY{l+s+s2}{\PYZdq{}}\PY{l+s+s2}{skyblue}\PY{l+s+s2}{\PYZdq{}}\PY{p}{,}
             \PY{n}{edgecolor}\PY{o}{=}\PY{l+s+s2}{\PYZdq{}}\PY{l+s+s2}{black}\PY{l+s+s2}{\PYZdq{}}\PY{p}{)}

    \PY{n}{plt}\PY{o}{.}\PY{n}{title}\PY{p}{(}\PY{l+s+sa}{f}\PY{l+s+s2}{\PYZdq{}}\PY{l+s+s2}{Distribution of }\PY{l+s+si}{\PYZob{}}\PY{n}{feature}\PY{l+s+si}{\PYZcb{}}\PY{l+s+s2}{\PYZdq{}}\PY{p}{,} \PY{n}{fontsize}\PY{o}{=}\PY{l+m+mi}{14}\PY{p}{)}
    \PY{n}{plt}\PY{o}{.}\PY{n}{xlabel}\PY{p}{(}\PY{n}{feature}\PY{p}{)}
    \PY{n}{plt}\PY{o}{.}\PY{n}{ylabel}\PY{p}{(}\PY{l+s+s2}{\PYZdq{}}\PY{l+s+s2}{Frequency}\PY{l+s+s2}{\PYZdq{}}\PY{p}{)}

    \PY{n}{plt}\PY{o}{.}\PY{n}{grid}\PY{p}{(}\PY{k+kc}{True}\PY{p}{,} \PY{n}{linestyle}\PY{o}{=}\PY{l+s+s2}{\PYZdq{}}\PY{l+s+s2}{\PYZhy{}\PYZhy{}}\PY{l+s+s2}{\PYZdq{}}\PY{p}{,} \PY{n}{alpha}\PY{o}{=}\PY{l+m+mf}{0.6}\PY{p}{)}

    \PY{n}{plt}\PY{o}{.}\PY{n}{show}\PY{p}{(}\PY{p}{)}
\end{Verbatim}
\end{tcolorbox}

    \begin{center}
    \adjustimage{max size={0.9\linewidth}{0.9\paperheight}}{output_109_0.png}
    \end{center}
    { \hspace*{\fill} \\}
    
    \begin{center}
    \adjustimage{max size={0.9\linewidth}{0.9\paperheight}}{output_109_1.png}
    \end{center}
    { \hspace*{\fill} \\}
    
    \begin{center}
    \adjustimage{max size={0.9\linewidth}{0.9\paperheight}}{output_109_2.png}
    \end{center}
    { \hspace*{\fill} \\}
    
    \begin{tcolorbox}[breakable, size=fbox, boxrule=1pt, pad at break*=1mm,colback=cellbackground, colframe=cellborder]
\prompt{In}{incolor}{110}{\boxspacing}
\begin{Verbatim}[commandchars=\\\{\}]
\PY{l+s+sd}{\PYZsq{}\PYZsq{}\PYZsq{}Median House Value: Most house values are concentrated in the middle range, with a cap at the upper limit.}
\PY{l+s+sd}{Median Income: Slightly right\PYZhy{}skewed, indicating fewer high\PYZhy{}income areas.}
\PY{l+s+sd}{Housing Median Age: Most houses are between 15–35 years old\PYZsq{}\PYZsq{}\PYZsq{}}
\end{Verbatim}
\end{tcolorbox}

            \begin{tcolorbox}[breakable, size=fbox, boxrule=.5pt, pad at break*=1mm, opacityfill=0]
\prompt{Out}{outcolor}{110}{\boxspacing}
\begin{Verbatim}[commandchars=\\\{\}]
'Median House Value: Most house values are concentrated in the middle range,
with a cap at the upper limit.\textbackslash{}nMedian Income: Slightly right-skewed, indicating
fewer high-income areas.\textbackslash{}nHousing Median Age: Most houses are between 15–35
years old'
\end{Verbatim}
\end{tcolorbox}
        
    \begin{tcolorbox}[breakable, size=fbox, boxrule=1pt, pad at break*=1mm,colback=cellbackground, colframe=cellborder]
\prompt{In}{incolor}{ }{\boxspacing}
\begin{Verbatim}[commandchars=\\\{\}]

\end{Verbatim}
\end{tcolorbox}

    \begin{tcolorbox}[breakable, size=fbox, boxrule=1pt, pad at break*=1mm,colback=cellbackground, colframe=cellborder]
\prompt{In}{incolor}{111}{\boxspacing}
\begin{Verbatim}[commandchars=\\\{\}]
\PY{c+c1}{\PYZsh{}boxplot}
\PY{k+kn}{import}\PY{+w}{ }\PY{n+nn}{seaborn}\PY{+w}{ }\PY{k}{as}\PY{+w}{ }\PY{n+nn}{sns}
\PY{k+kn}{import}\PY{+w}{ }\PY{n+nn}{matplotlib}\PY{n+nn}{.}\PY{n+nn}{pyplot}\PY{+w}{ }\PY{k}{as}\PY{+w}{ }\PY{n+nn}{plt}

\PY{n}{features} \PY{o}{=} \PY{p}{[}\PY{l+s+s2}{\PYZdq{}}\PY{l+s+s2}{total\PYZus{}rooms}\PY{l+s+s2}{\PYZdq{}}\PY{p}{,} \PY{l+s+s2}{\PYZdq{}}\PY{l+s+s2}{population}\PY{l+s+s2}{\PYZdq{}}\PY{p}{,} \PY{l+s+s2}{\PYZdq{}}\PY{l+s+s2}{median\PYZus{}income}\PY{l+s+s2}{\PYZdq{}}\PY{p}{]}

\PY{k}{for} \PY{n}{feature} \PY{o+ow}{in} \PY{n}{features}\PY{p}{:}

    \PY{n}{plt}\PY{o}{.}\PY{n}{figure}\PY{p}{(}\PY{n}{figsize}\PY{o}{=}\PY{p}{(}\PY{l+m+mi}{8}\PY{p}{,}\PY{l+m+mi}{5}\PY{p}{)}\PY{p}{)}

    \PY{n}{sns}\PY{o}{.}\PY{n}{boxplot}\PY{p}{(}
        \PY{n}{x}\PY{o}{=}\PY{n}{df}\PY{p}{[}\PY{n}{feature}\PY{p}{]}\PY{p}{,}
        \PY{n}{color}\PY{o}{=}\PY{l+s+s2}{\PYZdq{}}\PY{l+s+s2}{lightgreen}\PY{l+s+s2}{\PYZdq{}}
    \PY{p}{)}

    \PY{n}{plt}\PY{o}{.}\PY{n}{title}\PY{p}{(}\PY{l+s+sa}{f}\PY{l+s+s2}{\PYZdq{}}\PY{l+s+s2}{Box Plot of }\PY{l+s+si}{\PYZob{}}\PY{n}{feature}\PY{l+s+si}{\PYZcb{}}\PY{l+s+s2}{\PYZdq{}}\PY{p}{,} \PY{n}{fontsize}\PY{o}{=}\PY{l+m+mi}{14}\PY{p}{)}
    \PY{n}{plt}\PY{o}{.}\PY{n}{xlabel}\PY{p}{(}\PY{n}{feature}\PY{p}{)}

    \PY{n}{plt}\PY{o}{.}\PY{n}{grid}\PY{p}{(}\PY{k+kc}{True}\PY{p}{,} \PY{n}{linestyle}\PY{o}{=}\PY{l+s+s2}{\PYZdq{}}\PY{l+s+s2}{\PYZhy{}\PYZhy{}}\PY{l+s+s2}{\PYZdq{}}\PY{p}{,} \PY{n}{alpha}\PY{o}{=}\PY{l+m+mf}{0.5}\PY{p}{)}

    \PY{n}{plt}\PY{o}{.}\PY{n}{show}\PY{p}{(}\PY{p}{)}
\end{Verbatim}
\end{tcolorbox}

    \begin{center}
    \adjustimage{max size={0.9\linewidth}{0.9\paperheight}}{output_112_0.png}
    \end{center}
    { \hspace*{\fill} \\}
    
    \begin{center}
    \adjustimage{max size={0.9\linewidth}{0.9\paperheight}}{output_112_1.png}
    \end{center}
    { \hspace*{\fill} \\}
    
    \begin{center}
    \adjustimage{max size={0.9\linewidth}{0.9\paperheight}}{output_112_2.png}
    \end{center}
    { \hspace*{\fill} \\}
    
    \begin{tcolorbox}[breakable, size=fbox, boxrule=1pt, pad at break*=1mm,colback=cellbackground, colframe=cellborder]
\prompt{In}{incolor}{112}{\boxspacing}
\begin{Verbatim}[commandchars=\\\{\}]
\PY{l+s+sd}{\PYZsq{}\PYZsq{}\PYZsq{}total\PYZus{}rooms contains several high\PYZhy{}value outliers.}
\PY{l+s+sd}{population has many extreme observations.}
\PY{l+s+sd}{median\PYZus{}income contains fewer outliers than the other variables.\PYZsq{}\PYZsq{}\PYZsq{}}
\end{Verbatim}
\end{tcolorbox}

            \begin{tcolorbox}[breakable, size=fbox, boxrule=.5pt, pad at break*=1mm, opacityfill=0]
\prompt{Out}{outcolor}{112}{\boxspacing}
\begin{Verbatim}[commandchars=\\\{\}]
'total\_rooms contains several high-value outliers.\textbackslash{}npopulation has many extreme
observations.\textbackslash{}nmedian\_income contains fewer outliers than the other variables.'
\end{Verbatim}
\end{tcolorbox}
        
    \begin{tcolorbox}[breakable, size=fbox, boxrule=1pt, pad at break*=1mm,colback=cellbackground, colframe=cellborder]
\prompt{In}{incolor}{113}{\boxspacing}
\begin{Verbatim}[commandchars=\\\{\}]
\PY{c+c1}{\PYZsh{}Scatterplot}
\PY{n}{plt}\PY{o}{.}\PY{n}{figure}\PY{p}{(}\PY{n}{figsize}\PY{o}{=}\PY{p}{(}\PY{l+m+mi}{8}\PY{p}{,}\PY{l+m+mi}{6}\PY{p}{)}\PY{p}{)}

\PY{n}{sns}\PY{o}{.}\PY{n}{regplot}\PY{p}{(}
    \PY{n}{x}\PY{o}{=}\PY{l+s+s2}{\PYZdq{}}\PY{l+s+s2}{median\PYZus{}income}\PY{l+s+s2}{\PYZdq{}}\PY{p}{,}
    \PY{n}{y}\PY{o}{=}\PY{l+s+s2}{\PYZdq{}}\PY{l+s+s2}{median\PYZus{}house\PYZus{}value}\PY{l+s+s2}{\PYZdq{}}\PY{p}{,}
    \PY{n}{data}\PY{o}{=}\PY{n}{df}\PY{p}{,}
    \PY{n}{scatter\PYZus{}kws}\PY{o}{=}\PY{p}{\PYZob{}}\PY{l+s+s2}{\PYZdq{}}\PY{l+s+s2}{color}\PY{l+s+s2}{\PYZdq{}}\PY{p}{:}\PY{l+s+s2}{\PYZdq{}}\PY{l+s+s2}{blue}\PY{l+s+s2}{\PYZdq{}}\PY{p}{,}\PY{l+s+s2}{\PYZdq{}}\PY{l+s+s2}{alpha}\PY{l+s+s2}{\PYZdq{}}\PY{p}{:}\PY{l+m+mf}{0.4}\PY{p}{\PYZcb{}}\PY{p}{,}
    \PY{n}{line\PYZus{}kws}\PY{o}{=}\PY{p}{\PYZob{}}\PY{l+s+s2}{\PYZdq{}}\PY{l+s+s2}{color}\PY{l+s+s2}{\PYZdq{}}\PY{p}{:}\PY{l+s+s2}{\PYZdq{}}\PY{l+s+s2}{red}\PY{l+s+s2}{\PYZdq{}}\PY{p}{\PYZcb{}}
\PY{p}{)}

\PY{n}{plt}\PY{o}{.}\PY{n}{title}\PY{p}{(}\PY{l+s+s2}{\PYZdq{}}\PY{l+s+s2}{Median Income vs Median House Value}\PY{l+s+s2}{\PYZdq{}}\PY{p}{,} \PY{n}{fontsize}\PY{o}{=}\PY{l+m+mi}{15}\PY{p}{)}

\PY{n}{plt}\PY{o}{.}\PY{n}{xlabel}\PY{p}{(}\PY{l+s+s2}{\PYZdq{}}\PY{l+s+s2}{Median Income}\PY{l+s+s2}{\PYZdq{}}\PY{p}{)}

\PY{n}{plt}\PY{o}{.}\PY{n}{ylabel}\PY{p}{(}\PY{l+s+s2}{\PYZdq{}}\PY{l+s+s2}{Median House Value}\PY{l+s+s2}{\PYZdq{}}\PY{p}{)}

\PY{n}{plt}\PY{o}{.}\PY{n}{grid}\PY{p}{(}\PY{k+kc}{True}\PY{p}{)}

\PY{n}{plt}\PY{o}{.}\PY{n}{show}\PY{p}{(}\PY{p}{)}
\end{Verbatim}
\end{tcolorbox}

    \begin{center}
    \adjustimage{max size={0.9\linewidth}{0.9\paperheight}}{output_114_0.png}
    \end{center}
    { \hspace*{\fill} \\}
    
    \begin{tcolorbox}[breakable, size=fbox, boxrule=1pt, pad at break*=1mm,colback=cellbackground, colframe=cellborder]
\prompt{In}{incolor}{114}{\boxspacing}
\begin{Verbatim}[commandchars=\\\{\}]
\PY{l+s+sd}{\PYZsq{}\PYZsq{}\PYZsq{}A positive relationship exists between income and house value.}
\PY{l+s+sd}{As income increases, house prices generally increase.}
\PY{l+s+sd}{The regression line confirms the positive trend.\PYZsq{}\PYZsq{}\PYZsq{}}
\end{Verbatim}
\end{tcolorbox}

            \begin{tcolorbox}[breakable, size=fbox, boxrule=.5pt, pad at break*=1mm, opacityfill=0]
\prompt{Out}{outcolor}{114}{\boxspacing}
\begin{Verbatim}[commandchars=\\\{\}]
'A positive relationship exists between income and house value.\textbackslash{}nAs income
increases, house prices generally increase.\textbackslash{}nThe regression line confirms the
positive trend.'
\end{Verbatim}
\end{tcolorbox}
        
    \begin{tcolorbox}[breakable, size=fbox, boxrule=1pt, pad at break*=1mm,colback=cellbackground, colframe=cellborder]
\prompt{In}{incolor}{115}{\boxspacing}
\begin{Verbatim}[commandchars=\\\{\}]
\PY{c+c1}{\PYZsh{}correlationmap}
\PY{n}{plt}\PY{o}{.}\PY{n}{figure}\PY{p}{(}\PY{n}{figsize}\PY{o}{=}\PY{p}{(}\PY{l+m+mi}{10}\PY{p}{,}\PY{l+m+mi}{8}\PY{p}{)}\PY{p}{)}

\PY{n}{corr} \PY{o}{=} \PY{n}{df}\PY{o}{.}\PY{n}{select\PYZus{}dtypes}\PY{p}{(}\PY{n}{include}\PY{o}{=}\PY{l+s+s2}{\PYZdq{}}\PY{l+s+s2}{number}\PY{l+s+s2}{\PYZdq{}}\PY{p}{)}\PY{o}{.}\PY{n}{corr}\PY{p}{(}\PY{p}{)}

\PY{n}{sns}\PY{o}{.}\PY{n}{heatmap}\PY{p}{(}
    \PY{n}{corr}\PY{p}{,}
    \PY{n}{annot}\PY{o}{=}\PY{k+kc}{True}\PY{p}{,}
    \PY{n}{cmap}\PY{o}{=}\PY{l+s+s2}{\PYZdq{}}\PY{l+s+s2}{coolwarm}\PY{l+s+s2}{\PYZdq{}}\PY{p}{,}
    \PY{n}{fmt}\PY{o}{=}\PY{l+s+s2}{\PYZdq{}}\PY{l+s+s2}{.2f}\PY{l+s+s2}{\PYZdq{}}\PY{p}{,}
    \PY{n}{linewidths}\PY{o}{=}\PY{l+m+mf}{0.5}
\PY{p}{)}

\PY{n}{plt}\PY{o}{.}\PY{n}{title}\PY{p}{(}\PY{l+s+s2}{\PYZdq{}}\PY{l+s+s2}{Correlation Heatmap}\PY{l+s+s2}{\PYZdq{}}\PY{p}{,} \PY{n}{fontsize}\PY{o}{=}\PY{l+m+mi}{16}\PY{p}{)}

\PY{n}{plt}\PY{o}{.}\PY{n}{show}\PY{p}{(}\PY{p}{)}
\end{Verbatim}
\end{tcolorbox}

    \begin{center}
    \adjustimage{max size={0.9\linewidth}{0.9\paperheight}}{output_116_0.png}
    \end{center}
    { \hspace*{\fill} \\}
    
    \begin{tcolorbox}[breakable, size=fbox, boxrule=1pt, pad at break*=1mm,colback=cellbackground, colframe=cellborder]
\prompt{In}{incolor}{116}{\boxspacing}
\begin{Verbatim}[commandchars=\\\{\}]
\PY{l+s+sd}{\PYZsq{}\PYZsq{}\PYZsq{}}
\PY{l+s+sd}{Analysis}
\PY{l+s+sd}{median\PYZus{}income has the strongest positive correlation with median\PYZus{}house\PYZus{}value.}
\PY{l+s+sd}{total\PYZus{}rooms and total\PYZus{}bedrooms are highly correlated.}
\PY{l+s+sd}{population and households are also strongly correlated.\PYZsq{}\PYZsq{}\PYZsq{}}
\end{Verbatim}
\end{tcolorbox}

            \begin{tcolorbox}[breakable, size=fbox, boxrule=.5pt, pad at break*=1mm, opacityfill=0]
\prompt{Out}{outcolor}{116}{\boxspacing}
\begin{Verbatim}[commandchars=\\\{\}]
'\textbackslash{}nAnalysis\textbackslash{}nmedian\_income has the strongest positive correlation with
median\_house\_value.\textbackslash{}ntotal\_rooms and total\_bedrooms are highly
correlated.\textbackslash{}npopulation and households are also strongly correlated.'
\end{Verbatim}
\end{tcolorbox}
        
    \begin{tcolorbox}[breakable, size=fbox, boxrule=1pt, pad at break*=1mm,colback=cellbackground, colframe=cellborder]
\prompt{In}{incolor}{117}{\boxspacing}
\begin{Verbatim}[commandchars=\\\{\}]
\PY{c+c1}{\PYZsh{}Regression line plot}
\PY{n}{plt}\PY{o}{.}\PY{n}{figure}\PY{p}{(}\PY{n}{figsize}\PY{o}{=}\PY{p}{(}\PY{l+m+mi}{8}\PY{p}{,}\PY{l+m+mi}{6}\PY{p}{)}\PY{p}{)}

\PY{n}{plt}\PY{o}{.}\PY{n}{scatter}\PY{p}{(}
    \PY{n}{X\PYZus{}test\PYZus{}simple}\PY{p}{,}
    \PY{n}{y\PYZus{}test\PYZus{}simple}\PY{p}{,}
    \PY{n}{color}\PY{o}{=}\PY{l+s+s2}{\PYZdq{}}\PY{l+s+s2}{blue}\PY{l+s+s2}{\PYZdq{}}\PY{p}{,}
    \PY{n}{alpha}\PY{o}{=}\PY{l+m+mf}{0.5}\PY{p}{,}
    \PY{n}{label}\PY{o}{=}\PY{l+s+s2}{\PYZdq{}}\PY{l+s+s2}{Actual Data}\PY{l+s+s2}{\PYZdq{}}
\PY{p}{)}

\PY{n}{plt}\PY{o}{.}\PY{n}{plot}\PY{p}{(}
    \PY{n}{X\PYZus{}test\PYZus{}simple}\PY{p}{,}
    \PY{n}{y\PYZus{}pred\PYZus{}simple}\PY{p}{,}
    \PY{n}{color}\PY{o}{=}\PY{l+s+s2}{\PYZdq{}}\PY{l+s+s2}{red}\PY{l+s+s2}{\PYZdq{}}\PY{p}{,}
    \PY{n}{linewidth}\PY{o}{=}\PY{l+m+mi}{3}\PY{p}{,}
    \PY{n}{label}\PY{o}{=}\PY{l+s+s2}{\PYZdq{}}\PY{l+s+s2}{Regression Line}\PY{l+s+s2}{\PYZdq{}}
\PY{p}{)}

\PY{n}{plt}\PY{o}{.}\PY{n}{title}\PY{p}{(}\PY{l+s+s2}{\PYZdq{}}\PY{l+s+s2}{Simple Linear Regression}\PY{l+s+s2}{\PYZdq{}}\PY{p}{,} \PY{n}{fontsize}\PY{o}{=}\PY{l+m+mi}{15}\PY{p}{)}

\PY{n}{plt}\PY{o}{.}\PY{n}{xlabel}\PY{p}{(}\PY{l+s+s2}{\PYZdq{}}\PY{l+s+s2}{Median Income}\PY{l+s+s2}{\PYZdq{}}\PY{p}{)}

\PY{n}{plt}\PY{o}{.}\PY{n}{ylabel}\PY{p}{(}\PY{l+s+s2}{\PYZdq{}}\PY{l+s+s2}{Median House Value}\PY{l+s+s2}{\PYZdq{}}\PY{p}{)}

\PY{n}{plt}\PY{o}{.}\PY{n}{legend}\PY{p}{(}\PY{p}{)}

\PY{n}{plt}\PY{o}{.}\PY{n}{grid}\PY{p}{(}\PY{k+kc}{True}\PY{p}{)}

\PY{n}{plt}\PY{o}{.}\PY{n}{show}\PY{p}{(}\PY{p}{)}
\end{Verbatim}
\end{tcolorbox}

    \begin{center}
    \adjustimage{max size={0.9\linewidth}{0.9\paperheight}}{output_118_0.png}
    \end{center}
    { \hspace*{\fill} \\}
    
    \begin{tcolorbox}[breakable, size=fbox, boxrule=1pt, pad at break*=1mm,colback=cellbackground, colframe=cellborder]
\prompt{In}{incolor}{118}{\boxspacing}
\begin{Verbatim}[commandchars=\\\{\}]
\PY{l+s+sd}{\PYZsq{}\PYZsq{}\PYZsq{}The regression line fits the general upward trend.}
\PY{l+s+sd}{Some observations lie far from the line, indicating prediction errors\PYZsq{}\PYZsq{}\PYZsq{}}
\end{Verbatim}
\end{tcolorbox}

            \begin{tcolorbox}[breakable, size=fbox, boxrule=.5pt, pad at break*=1mm, opacityfill=0]
\prompt{Out}{outcolor}{118}{\boxspacing}
\begin{Verbatim}[commandchars=\\\{\}]
'The regression line fits the general upward trend.\textbackslash{}nSome observations lie far
from the line, indicating prediction errors'
\end{Verbatim}
\end{tcolorbox}
        
    \begin{tcolorbox}[breakable, size=fbox, boxrule=1pt, pad at break*=1mm,colback=cellbackground, colframe=cellborder]
\prompt{In}{incolor}{119}{\boxspacing}
\begin{Verbatim}[commandchars=\\\{\}]
\PY{c+c1}{\PYZsh{}Residual Plot}
\PY{n}{residuals} \PY{o}{=} \PY{n}{y\PYZus{}test\PYZus{}simple} \PY{o}{\PYZhy{}} \PY{n}{y\PYZus{}pred\PYZus{}simple}

\PY{n}{plt}\PY{o}{.}\PY{n}{figure}\PY{p}{(}\PY{n}{figsize}\PY{o}{=}\PY{p}{(}\PY{l+m+mi}{8}\PY{p}{,}\PY{l+m+mi}{6}\PY{p}{)}\PY{p}{)}

\PY{n}{plt}\PY{o}{.}\PY{n}{scatter}\PY{p}{(}
    \PY{n}{y\PYZus{}pred\PYZus{}simple}\PY{p}{,}
    \PY{n}{residuals}\PY{p}{,}
    \PY{n}{color}\PY{o}{=}\PY{l+s+s2}{\PYZdq{}}\PY{l+s+s2}{purple}\PY{l+s+s2}{\PYZdq{}}\PY{p}{,}
    \PY{n}{alpha}\PY{o}{=}\PY{l+m+mf}{0.5}
\PY{p}{)}

\PY{n}{plt}\PY{o}{.}\PY{n}{axhline}\PY{p}{(}
    \PY{n}{y}\PY{o}{=}\PY{l+m+mi}{0}\PY{p}{,}
    \PY{n}{color}\PY{o}{=}\PY{l+s+s2}{\PYZdq{}}\PY{l+s+s2}{red}\PY{l+s+s2}{\PYZdq{}}\PY{p}{,}
    \PY{n}{linestyle}\PY{o}{=}\PY{l+s+s2}{\PYZdq{}}\PY{l+s+s2}{\PYZhy{}\PYZhy{}}\PY{l+s+s2}{\PYZdq{}}\PY{p}{,}
    \PY{n}{linewidth}\PY{o}{=}\PY{l+m+mi}{2}
\PY{p}{)}

\PY{n}{plt}\PY{o}{.}\PY{n}{title}\PY{p}{(}\PY{l+s+s2}{\PYZdq{}}\PY{l+s+s2}{Residual Plot}\PY{l+s+s2}{\PYZdq{}}\PY{p}{,} \PY{n}{fontsize}\PY{o}{=}\PY{l+m+mi}{15}\PY{p}{)}

\PY{n}{plt}\PY{o}{.}\PY{n}{xlabel}\PY{p}{(}\PY{l+s+s2}{\PYZdq{}}\PY{l+s+s2}{Predicted House Value}\PY{l+s+s2}{\PYZdq{}}\PY{p}{)}

\PY{n}{plt}\PY{o}{.}\PY{n}{ylabel}\PY{p}{(}\PY{l+s+s2}{\PYZdq{}}\PY{l+s+s2}{Residuals}\PY{l+s+s2}{\PYZdq{}}\PY{p}{)}

\PY{n}{plt}\PY{o}{.}\PY{n}{grid}\PY{p}{(}\PY{k+kc}{True}\PY{p}{)}

\PY{n}{plt}\PY{o}{.}\PY{n}{show}\PY{p}{(}\PY{p}{)}
\end{Verbatim}
\end{tcolorbox}

    \begin{center}
    \adjustimage{max size={0.9\linewidth}{0.9\paperheight}}{output_120_0.png}
    \end{center}
    { \hspace*{\fill} \\}
    
    \begin{tcolorbox}[breakable, size=fbox, boxrule=1pt, pad at break*=1mm,colback=cellbackground, colframe=cellborder]
\prompt{In}{incolor}{120}{\boxspacing}
\begin{Verbatim}[commandchars=\\\{\}]
\PY{l+s+sd}{\PYZsq{}\PYZsq{}\PYZsq{}Residuals are scattered around zero, suggesting the model captures the overall trend}
\PY{l+s+sd}{Some large residuals remain, indicating prediction errors for certain observations\PYZsq{}\PYZsq{}\PYZsq{}}
\end{Verbatim}
\end{tcolorbox}

            \begin{tcolorbox}[breakable, size=fbox, boxrule=.5pt, pad at break*=1mm, opacityfill=0]
\prompt{Out}{outcolor}{120}{\boxspacing}
\begin{Verbatim}[commandchars=\\\{\}]
'Residuals are scattered around zero, suggesting the model captures the overall
trend\textbackslash{}nSome large residuals remain, indicating prediction errors for certain
observations'
\end{Verbatim}
\end{tcolorbox}
        
    \begin{tcolorbox}[breakable, size=fbox, boxrule=1pt, pad at break*=1mm,colback=cellbackground, colframe=cellborder]
\prompt{In}{incolor}{121}{\boxspacing}
\begin{Verbatim}[commandchars=\\\{\}]
\PY{l+s+sd}{\PYZsq{}\PYZsq{}\PYZsq{}The visualizations provide a clear understanding of the California Housing dataset. Histograms show that several numerical features are right\PYZhy{}skewed, while box plots reveal the presence of outliers in variables such as total\PYZus{}rooms and population. The scatter plot and regression line indicate a strong positive relationship between median\PYZus{}income and median\PYZus{}house\PYZus{}value, making it the most influential predictor. The correlation heatmap confirms this relationship and highlights correlations among other features. Finally, the residual plot shows that the linear regression model captures the overall trend, although some prediction errors remain, suggesting that additional features or more complex models could further improve predictive performance.\PYZsq{}\PYZsq{}\PYZsq{}}
\end{Verbatim}
\end{tcolorbox}

            \begin{tcolorbox}[breakable, size=fbox, boxrule=.5pt, pad at break*=1mm, opacityfill=0]
\prompt{Out}{outcolor}{121}{\boxspacing}
\begin{Verbatim}[commandchars=\\\{\}]
'The visualizations provide a clear understanding of the California Housing
dataset. Histograms show that several numerical features are right-skewed, while
box plots reveal the presence of outliers in variables such as total\_rooms and
population. The scatter plot and regression line indicate a strong positive
relationship between median\_income and median\_house\_value, making it the most
influential predictor. The correlation heatmap confirms this relationship and
highlights correlations among other features. Finally, the residual plot shows
that the linear regression model captures the overall trend, although some
prediction errors remain, suggesting that additional features or more complex
models could further improve predictive performance.'
\end{Verbatim}
\end{tcolorbox}
        
    \begin{tcolorbox}[breakable, size=fbox, boxrule=1pt, pad at break*=1mm,colback=cellbackground, colframe=cellborder]
\prompt{In}{incolor}{122}{\boxspacing}
\begin{Verbatim}[commandchars=\\\{\}]
                  \PY{c+c1}{\PYZsh{}PART\PYZhy{}8(Conclusion \PYZam{} Reflection )}
\end{Verbatim}
\end{tcolorbox}

    \begin{tcolorbox}[breakable, size=fbox, boxrule=1pt, pad at break*=1mm,colback=cellbackground, colframe=cellborder]
\prompt{In}{incolor}{123}{\boxspacing}
\begin{Verbatim}[commandchars=\\\{\}]
\PY{l+s+sd}{\PYZsq{}\PYZsq{}\PYZsq{}}

\PY{l+s+sd}{1. What did you learn from this assignment?}

\PY{l+s+sd}{\PYZhy{} I learned how to import and inspect a dataset using Python.}
\PY{l+s+sd}{\PYZhy{} I learned how to clean data by handling missing values and checking for duplicates.}
\PY{l+s+sd}{\PYZhy{} I learned how to perform Exploratory Data Analysis (EDA) using summary statistics and visualizations.}
\PY{l+s+sd}{\PYZhy{} I learned how to build Simple Linear Regression and Multiple Linear Regression models.}
\PY{l+s+sd}{\PYZhy{} I learned how to evaluate model performance using R² Score, MAE, and RMSE.}


\PY{l+s+sd}{2. What was the most surprising finding in your data?}

\PY{l+s+sd}{\PYZhy{} The most surprising finding was that Median Income had the strongest relationship with Median House Value.}
\PY{l+s+sd}{\PYZhy{} I expected features like Total Rooms or Housing Median Age to have a stronger impact, but income was the best predictor.}
\PY{l+s+sd}{\PYZhy{} I also found that geographical features such as Latitude and Longitude significantly influence house prices.}


\PY{l+s+sd}{ 3. What challenges did you face and how did you overcome them?}

\PY{l+s+sd}{\PYZhy{} I faced missing library errors while running Python code.}
\PY{l+s+sd}{\PYZhy{} I encountered errors while creating regression models and comparison tables.}
\PY{l+s+sd}{\PYZhy{} I carefully read the error messages, corrected the code, and installed the required libraries}
\PY{l+s+sd}{\PYZhy{} After debugging the code step by step, I was able to complete the analysis successfully}


\PY{l+s+sd}{4. How would this analysis help in a real business context?}

\PY{l+s+sd}{\PYZhy{} Real estate companies can use this analysis to estimate house prices.}
\PY{l+s+sd}{\PYZhy{} Buyers and sellers can make better decisions based on predicted house prices.}
\PY{l+s+sd}{\PYZhy{} Businesses can identify the most important factors affecting property prices.}

\PY{l+s+sd}{5. What would you do differently with more time or data?}

\PY{l+s+sd}{\PYZhy{} I would include additional features such as nearby schools, hospitals, transportation, and crime rates.}
\PY{l+s+sd}{\PYZhy{} I would perform more feature engineering to improve prediction accuracy.}


\PY{l+s+sd}{6. How do these statistical concepts connect to machine learning?}

\PY{l+s+sd}{\PYZhy{} Statistics helps us understand the dataset before building a machine learning model.}
\PY{l+s+sd}{\PYZhy{} Correlation helps in selecting the most important features.}
\PY{l+s+sd}{\PYZhy{} Mean, median, and standard deviation help summarize the data.}
\PY{l+s+sd}{\PYZhy{} Regression is one of the basic machine learning algorithms used for prediction.}
\PY{l+s+sd}{\PYZhy{} Evaluation metrics such as R² Score, MAE, and RMSE help measure model performance}

\PY{l+s+sd}{ 7. What questions do you still have?}

\PY{l+s+sd}{\PYZhy{} How do advanced machine learning models improve prediction accuracy?}
\PY{l+s+sd}{\PYZhy{} What are the best techniques for handling large real\PYZhy{}world datasets?\PYZsq{}\PYZsq{}\PYZsq{}}
\end{Verbatim}
\end{tcolorbox}

            \begin{tcolorbox}[breakable, size=fbox, boxrule=.5pt, pad at break*=1mm, opacityfill=0]
\prompt{Out}{outcolor}{123}{\boxspacing}
\begin{Verbatim}[commandchars=\\\{\}]
'\textbackslash{}n\textbackslash{}n1. What did you learn from this assignment?\textbackslash{}n\textbackslash{}n- I learned how to import
and inspect a dataset using Python.\textbackslash{}n- I learned how to clean data by handling
missing values and checking for duplicates.\textbackslash{}n- I learned how to perform
Exploratory Data Analysis (EDA) using summary statistics and visualizations.\textbackslash{}n-
I learned how to build Simple Linear Regression and Multiple Linear Regression
models.\textbackslash{}n- I learned how to evaluate model performance using R² Score, MAE, and
RMSE.\textbackslash{}n\textbackslash{}n\textbackslash{}n2. What was the most surprising finding in your data?\textbackslash{}n\textbackslash{}n- The most
surprising finding was that Median Income had the strongest relationship with
Median House Value.\textbackslash{}n- I expected features like Total Rooms or Housing Median
Age to have a stronger impact, but income was the best predictor.\textbackslash{}n- I also
found that geographical features such as Latitude and Longitude significantly
influence house prices.\textbackslash{}n\textbackslash{}n\textbackslash{}n 3. What challenges did you face and how did you
overcome them?\textbackslash{}n\textbackslash{}n- I faced missing library errors while running Python code.\textbackslash{}n-
I encountered errors while creating regression models and comparison tables.\textbackslash{}n-
I carefully read the error messages, corrected the code, and installed the
required libraries\textbackslash{}n- After debugging the code step by step, I was able to
complete the analysis successfully\textbackslash{}n\textbackslash{}n\textbackslash{}n4. How would this analysis help in a
real business context?\textbackslash{}n\textbackslash{}n- Real estate companies can use this analysis to
estimate house prices.\textbackslash{}n- Buyers and sellers can make better decisions based on
predicted house prices.\textbackslash{}n- Businesses can identify the most important factors
affecting property prices.\textbackslash{}n\textbackslash{}n5. What would you do differently with more time or
data?\textbackslash{}n\textbackslash{}n- I would include additional features such as nearby schools,
hospitals, transportation, and crime rates.\textbackslash{}n- I would perform more feature
engineering to improve prediction accuracy.\textbackslash{}n\textbackslash{}n\textbackslash{}n6. How do these statistical
concepts connect to machine learning?\textbackslash{}n\textbackslash{}n- Statistics helps us understand the
dataset before building a machine learning model.\textbackslash{}n- Correlation helps in
selecting the most important features.\textbackslash{}n- Mean, median, and standard deviation
help summarize the data.\textbackslash{}n- Regression is one of the basic machine learning
algorithms used for prediction.\textbackslash{}n- Evaluation metrics such as R² Score, MAE, and
RMSE help measure model performance\textbackslash{}n\textbackslash{}n 7. What questions do you still
have?\textbackslash{}n\textbackslash{}n- How do advanced machine learning models improve prediction
accuracy?\textbackslash{}n- What are the best techniques for handling large real-world
datasets?'
\end{Verbatim}
\end{tcolorbox}
        
    \begin{tcolorbox}[breakable, size=fbox, boxrule=1pt, pad at break*=1mm,colback=cellbackground, colframe=cellborder]
\prompt{In}{incolor}{124}{\boxspacing}
\begin{Verbatim}[commandchars=\\\{\}]
                 \PY{c+c1}{\PYZsh{}Compare with random foreset}
\end{Verbatim}
\end{tcolorbox}

    \begin{tcolorbox}[breakable, size=fbox, boxrule=1pt, pad at break*=1mm,colback=cellbackground, colframe=cellborder]
\prompt{In}{incolor}{125}{\boxspacing}
\begin{Verbatim}[commandchars=\\\{\}]
\PY{c+c1}{\PYZsh{}Compare with random foreset}
\PY{k+kn}{from}\PY{+w}{ }\PY{n+nn}{sklearn}\PY{n+nn}{.}\PY{n+nn}{ensemble}\PY{+w}{ }\PY{k+kn}{import} \PY{n}{RandomForestRegressor}
\PY{k+kn}{from}\PY{+w}{ }\PY{n+nn}{sklearn}\PY{n+nn}{.}\PY{n+nn}{metrics}\PY{+w}{ }\PY{k+kn}{import} \PY{n}{r2\PYZus{}score}\PY{p}{,} \PY{n}{mean\PYZus{}absolute\PYZus{}error}\PY{p}{,} \PY{n}{mean\PYZus{}squared\PYZus{}error}
\PY{k+kn}{import}\PY{+w}{ }\PY{n+nn}{numpy}\PY{+w}{ }\PY{k}{as}\PY{+w}{ }\PY{n+nn}{np}
\end{Verbatim}
\end{tcolorbox}

    \begin{tcolorbox}[breakable, size=fbox, boxrule=1pt, pad at break*=1mm,colback=cellbackground, colframe=cellborder]
\prompt{In}{incolor}{126}{\boxspacing}
\begin{Verbatim}[commandchars=\\\{\}]
\PY{k+kn}{from}\PY{+w}{ }\PY{n+nn}{sklearn}\PY{n+nn}{.}\PY{n+nn}{ensemble}\PY{+w}{ }\PY{k+kn}{import} \PY{n}{RandomForestRegressor}

\PY{n}{rf\PYZus{}model} \PY{o}{=} \PY{n}{RandomForestRegressor}\PY{p}{(}
    \PY{n}{n\PYZus{}estimators}\PY{o}{=}\PY{l+m+mi}{100}\PY{p}{,}
    \PY{n}{random\PYZus{}state}\PY{o}{=}\PY{l+m+mi}{42}
\PY{p}{)}

\PY{n}{rf\PYZus{}model}\PY{o}{.}\PY{n}{fit}\PY{p}{(}\PY{n}{X\PYZus{}train}\PY{p}{,} \PY{n}{y\PYZus{}train}\PY{p}{)}
\end{Verbatim}
\end{tcolorbox}

            \begin{tcolorbox}[breakable, size=fbox, boxrule=.5pt, pad at break*=1mm, opacityfill=0]
\prompt{Out}{outcolor}{126}{\boxspacing}
\begin{Verbatim}[commandchars=\\\{\}]
RandomForestRegressor(random\_state=42)
\end{Verbatim}
\end{tcolorbox}
        
    \begin{tcolorbox}[breakable, size=fbox, boxrule=1pt, pad at break*=1mm,colback=cellbackground, colframe=cellborder]
\prompt{In}{incolor}{128}{\boxspacing}
\begin{Verbatim}[commandchars=\\\{\}]
\PY{c+c1}{\PYZsh{}prediction}

\PY{n}{rf\PYZus{}pred} \PY{o}{=} \PY{n}{rf\PYZus{}model}\PY{o}{.}\PY{n}{predict}\PY{p}{(}\PY{n}{X\PYZus{}test}\PY{p}{)}
\end{Verbatim}
\end{tcolorbox}

    \begin{tcolorbox}[breakable, size=fbox, boxrule=1pt, pad at break*=1mm,colback=cellbackground, colframe=cellborder]
\prompt{In}{incolor}{129}{\boxspacing}
\begin{Verbatim}[commandchars=\\\{\}]
\PY{n}{rf\PYZus{}r2} \PY{o}{=} \PY{n}{r2\PYZus{}score}\PY{p}{(}\PY{n}{y\PYZus{}test}\PY{p}{,} \PY{n}{rf\PYZus{}pred}\PY{p}{)}

\PY{n}{rf\PYZus{}mae} \PY{o}{=} \PY{n}{mean\PYZus{}absolute\PYZus{}error}\PY{p}{(}\PY{n}{y\PYZus{}test}\PY{p}{,} \PY{n}{rf\PYZus{}pred}\PY{p}{)}

\PY{n}{rf\PYZus{}rmse} \PY{o}{=} \PY{n}{np}\PY{o}{.}\PY{n}{sqrt}\PY{p}{(}\PY{n}{mean\PYZus{}squared\PYZus{}error}\PY{p}{(}\PY{n}{y\PYZus{}test}\PY{p}{,} \PY{n}{rf\PYZus{}pred}\PY{p}{)}\PY{p}{)}

\PY{n+nb}{print}\PY{p}{(}\PY{l+s+s2}{\PYZdq{}}\PY{l+s+s2}{Random Forest Results}\PY{l+s+s2}{\PYZdq{}}\PY{p}{)}
\PY{n+nb}{print}\PY{p}{(}\PY{l+s+s2}{\PYZdq{}}\PY{l+s+s2}{R² Score:}\PY{l+s+s2}{\PYZdq{}}\PY{p}{,} \PY{n}{rf\PYZus{}r2}\PY{p}{)}
\PY{n+nb}{print}\PY{p}{(}\PY{l+s+s2}{\PYZdq{}}\PY{l+s+s2}{MAE:}\PY{l+s+s2}{\PYZdq{}}\PY{p}{,} \PY{n}{rf\PYZus{}mae}\PY{p}{)}
\PY{n+nb}{print}\PY{p}{(}\PY{l+s+s2}{\PYZdq{}}\PY{l+s+s2}{RMSE:}\PY{l+s+s2}{\PYZdq{}}\PY{p}{,} \PY{n}{rf\PYZus{}rmse}\PY{p}{)}
\end{Verbatim}
\end{tcolorbox}

    \begin{Verbatim}[commandchars=\\\{\}]
Random Forest Results
R² Score: 0.8110495190488294
MAE: 32028.281543120156
RMSE: 49759.66304963792
    \end{Verbatim}

    \begin{tcolorbox}[breakable, size=fbox, boxrule=1pt, pad at break*=1mm,colback=cellbackground, colframe=cellborder]
\prompt{In}{incolor}{137}{\boxspacing}
\begin{Verbatim}[commandchars=\\\{\}]
\PY{c+c1}{\PYZsh{}Comparison table}
\PY{n}{comparison} \PY{o}{=} \PY{n}{pd}\PY{o}{.}\PY{n}{DataFrame}\PY{p}{(}\PY{p}{\PYZob{}}
    \PY{l+s+s2}{\PYZdq{}}\PY{l+s+s2}{Model}\PY{l+s+s2}{\PYZdq{}}\PY{p}{:}\PY{p}{[}
        \PY{l+s+s2}{\PYZdq{}}\PY{l+s+s2}{Simple Linear Regression}\PY{l+s+s2}{\PYZdq{}}\PY{p}{,}
        \PY{l+s+s2}{\PYZdq{}}\PY{l+s+s2}{Multiple Linear Regression}\PY{l+s+s2}{\PYZdq{}}\PY{p}{,}
        \PY{l+s+s2}{\PYZdq{}}\PY{l+s+s2}{Random Forest}\PY{l+s+s2}{\PYZdq{}}
    \PY{p}{]}\PY{p}{,}

    \PY{l+s+s2}{\PYZdq{}}\PY{l+s+s2}{R² Score}\PY{l+s+s2}{\PYZdq{}}\PY{p}{:}\PY{p}{[}
        \PY{n}{simple\PYZus{}r2}\PY{p}{,}
        \PY{n}{multiple\PYZus{}r2}\PY{p}{,}
        \PY{n}{rf\PYZus{}r2}
    \PY{p}{]}\PY{p}{,}

    \PY{l+s+s2}{\PYZdq{}}\PY{l+s+s2}{MAE}\PY{l+s+s2}{\PYZdq{}}\PY{p}{:}\PY{p}{[}
        \PY{n}{simple\PYZus{}mae}\PY{p}{,}
        \PY{n}{multiple\PYZus{}mae}\PY{p}{,}
        \PY{n}{rf\PYZus{}mae}
    \PY{p}{]}\PY{p}{,}

    \PY{l+s+s2}{\PYZdq{}}\PY{l+s+s2}{RMSE}\PY{l+s+s2}{\PYZdq{}}\PY{p}{:}\PY{p}{[}
        \PY{n}{simple\PYZus{}rmse}\PY{p}{,}
        \PY{n}{multiple\PYZus{}rmse}\PY{p}{,}
        \PY{n}{rf\PYZus{}rmse}
    \PY{p}{]}
\PY{p}{\PYZcb{}}\PY{p}{)}

\PY{n}{comparison}
\end{Verbatim}
\end{tcolorbox}

            \begin{tcolorbox}[breakable, size=fbox, boxrule=.5pt, pad at break*=1mm, opacityfill=0]
\prompt{Out}{outcolor}{137}{\boxspacing}
\begin{Verbatim}[commandchars=\\\{\}]
                        Model  R² Score           MAE          RMSE
0    Simple Linear Regression  0.458859  62990.865301  84209.012414
1  Multiple Linear Regression  0.584449  54230.744477  73793.096814
2               Random Forest  0.811050  32028.281543  49759.663050
\end{Verbatim}
\end{tcolorbox}
        
    \begin{tcolorbox}[breakable, size=fbox, boxrule=1pt, pad at break*=1mm,colback=cellbackground, colframe=cellborder]
\prompt{In}{incolor}{138}{\boxspacing}
\begin{Verbatim}[commandchars=\\\{\}]
\PY{l+s+sd}{\PYZsq{}\PYZsq{}\PYZsq{}The Random Forest model generally performs better than Linear Regression because it can capture complex and non\PYZhy{}linear relationships between features and house prices. Unlike Linear Regression, Random Forest builds multiple decision trees and combines their predictions, reducing overfitting and improving accuracy. This often results in a higher R² score and lower MAE and RMSE values.\PYZsq{}\PYZsq{}\PYZsq{}}
\end{Verbatim}
\end{tcolorbox}

            \begin{tcolorbox}[breakable, size=fbox, boxrule=.5pt, pad at break*=1mm, opacityfill=0]
\prompt{Out}{outcolor}{138}{\boxspacing}
\begin{Verbatim}[commandchars=\\\{\}]
'The Random Forest model generally performs better than Linear Regression
because it can capture complex and non-linear relationships between features and
house prices. Unlike Linear Regression, Random Forest builds multiple decision
trees and combines their predictions, reducing overfitting and improving
accuracy. This often results in a higher R² score and lower MAE and RMSE
values.'
\end{Verbatim}
\end{tcolorbox}
        
    \begin{tcolorbox}[breakable, size=fbox, boxrule=1pt, pad at break*=1mm,colback=cellbackground, colframe=cellborder]
\prompt{In}{incolor}{139}{\boxspacing}
\begin{Verbatim}[commandchars=\\\{\}]
                   \PY{c+c1}{\PYZsh{} Feature Importance}
\end{Verbatim}
\end{tcolorbox}

    \begin{tcolorbox}[breakable, size=fbox, boxrule=1pt, pad at break*=1mm,colback=cellbackground, colframe=cellborder]
\prompt{In}{incolor}{140}{\boxspacing}
\begin{Verbatim}[commandchars=\\\{\}]
\PY{c+c1}{\PYZsh{} Feature Importance}
\PY{n}{importance} \PY{o}{=} \PY{n}{pd}\PY{o}{.}\PY{n}{DataFrame}\PY{p}{(}\PY{p}{\PYZob{}}
    \PY{l+s+s2}{\PYZdq{}}\PY{l+s+s2}{Feature}\PY{l+s+s2}{\PYZdq{}}\PY{p}{:} \PY{n}{X}\PY{o}{.}\PY{n}{columns}\PY{p}{,}
    \PY{l+s+s2}{\PYZdq{}}\PY{l+s+s2}{Importance}\PY{l+s+s2}{\PYZdq{}}\PY{p}{:} \PY{n}{rf\PYZus{}model}\PY{o}{.}\PY{n}{feature\PYZus{}importances\PYZus{}}
\PY{p}{\PYZcb{}}\PY{p}{)}

\PY{n}{importance} \PY{o}{=} \PY{n}{importance}\PY{o}{.}\PY{n}{sort\PYZus{}values}\PY{p}{(}
    \PY{n}{by}\PY{o}{=}\PY{l+s+s2}{\PYZdq{}}\PY{l+s+s2}{Importance}\PY{l+s+s2}{\PYZdq{}}\PY{p}{,}
    \PY{n}{ascending}\PY{o}{=}\PY{k+kc}{False}
\PY{p}{)}

\PY{n+nb}{print}\PY{p}{(}\PY{n}{importance}\PY{p}{)}
\end{Verbatim}
\end{tcolorbox}

    \begin{Verbatim}[commandchars=\\\{\}]
              Feature  Importance
0       median\_income    0.530068
4           longitude    0.176120
3            latitude    0.168436
1  housing\_median\_age    0.072429
2         total\_rooms    0.052947
    \end{Verbatim}

    \begin{tcolorbox}[breakable, size=fbox, boxrule=1pt, pad at break*=1mm,colback=cellbackground, colframe=cellborder]
\prompt{In}{incolor}{141}{\boxspacing}
\begin{Verbatim}[commandchars=\\\{\}]
\PY{c+c1}{\PYZsh{}barchat}
\PY{k+kn}{import}\PY{+w}{ }\PY{n+nn}{matplotlib}\PY{n+nn}{.}\PY{n+nn}{pyplot}\PY{+w}{ }\PY{k}{as}\PY{+w}{ }\PY{n+nn}{plt}
\PY{k+kn}{import}\PY{+w}{ }\PY{n+nn}{seaborn}\PY{+w}{ }\PY{k}{as}\PY{+w}{ }\PY{n+nn}{sns}

\PY{n}{plt}\PY{o}{.}\PY{n}{figure}\PY{p}{(}\PY{n}{figsize}\PY{o}{=}\PY{p}{(}\PY{l+m+mi}{8}\PY{p}{,}\PY{l+m+mi}{5}\PY{p}{)}\PY{p}{)}

\PY{n}{sns}\PY{o}{.}\PY{n}{barplot}\PY{p}{(}
    \PY{n}{data}\PY{o}{=}\PY{n}{importance}\PY{p}{,}
    \PY{n}{x}\PY{o}{=}\PY{l+s+s2}{\PYZdq{}}\PY{l+s+s2}{Importance}\PY{l+s+s2}{\PYZdq{}}\PY{p}{,}
    \PY{n}{y}\PY{o}{=}\PY{l+s+s2}{\PYZdq{}}\PY{l+s+s2}{Feature}\PY{l+s+s2}{\PYZdq{}}\PY{p}{,}
    \PY{n}{color}\PY{o}{=}\PY{l+s+s2}{\PYZdq{}}\PY{l+s+s2}{steelblue}\PY{l+s+s2}{\PYZdq{}}
\PY{p}{)}

\PY{n}{plt}\PY{o}{.}\PY{n}{title}\PY{p}{(}\PY{l+s+s2}{\PYZdq{}}\PY{l+s+s2}{Feature Importance using Random Forest}\PY{l+s+s2}{\PYZdq{}}\PY{p}{)}

\PY{n}{plt}\PY{o}{.}\PY{n}{xlabel}\PY{p}{(}\PY{l+s+s2}{\PYZdq{}}\PY{l+s+s2}{Importance Score}\PY{l+s+s2}{\PYZdq{}}\PY{p}{)}

\PY{n}{plt}\PY{o}{.}\PY{n}{ylabel}\PY{p}{(}\PY{l+s+s2}{\PYZdq{}}\PY{l+s+s2}{Feature}\PY{l+s+s2}{\PYZdq{}}\PY{p}{)}

\PY{n}{plt}\PY{o}{.}\PY{n}{grid}\PY{p}{(}\PY{k+kc}{True}\PY{p}{,} \PY{n}{linestyle}\PY{o}{=}\PY{l+s+s2}{\PYZdq{}}\PY{l+s+s2}{\PYZhy{}\PYZhy{}}\PY{l+s+s2}{\PYZdq{}}\PY{p}{,} \PY{n}{alpha}\PY{o}{=}\PY{l+m+mf}{0.5}\PY{p}{)}

\PY{n}{plt}\PY{o}{.}\PY{n}{show}\PY{p}{(}\PY{p}{)}
\end{Verbatim}
\end{tcolorbox}

    \begin{center}
    \adjustimage{max size={0.9\linewidth}{0.9\paperheight}}{output_134_0.png}
    \end{center}
    { \hspace*{\fill} \\}
    
    \begin{tcolorbox}[breakable, size=fbox, boxrule=1pt, pad at break*=1mm,colback=cellbackground, colframe=cellborder]
\prompt{In}{incolor}{142}{\boxspacing}
\begin{Verbatim}[commandchars=\\\{\}]
\PY{l+s+sd}{\PYZsq{}\PYZsq{}\PYZsq{}The feature importance analysis shows the contribution of each feature to the Random Forest model. Features with higher importance scores have a greater impact on predicting house prices. In this dataset, median\PYZus{}income is expected to be the most influential predictor, followed by geographical features such as latitude and longitude.\PYZsq{}\PYZsq{}\PYZsq{}}
\end{Verbatim}
\end{tcolorbox}

            \begin{tcolorbox}[breakable, size=fbox, boxrule=.5pt, pad at break*=1mm, opacityfill=0]
\prompt{Out}{outcolor}{142}{\boxspacing}
\begin{Verbatim}[commandchars=\\\{\}]
'The feature importance analysis shows the contribution of each feature to the
Random Forest model. Features with higher importance scores have a greater
impact on predicting house prices. In this dataset, median\_income is expected to
be the most influential predictor, followed by geographical features such as
latitude and longitude.'
\end{Verbatim}
\end{tcolorbox}
        
    \begin{tcolorbox}[breakable, size=fbox, boxrule=1pt, pad at break*=1mm,colback=cellbackground, colframe=cellborder]
\prompt{In}{incolor}{143}{\boxspacing}
\begin{Verbatim}[commandchars=\\\{\}]
\PY{k+kn}{import}\PY{+w}{ }\PY{n+nn}{joblib}

\PY{n}{joblib}\PY{o}{.}\PY{n}{dump}\PY{p}{(}\PY{n}{rf\PYZus{}model}\PY{p}{,} \PY{l+s+s2}{\PYZdq{}}\PY{l+s+s2}{house\PYZus{}model.pkl}\PY{l+s+s2}{\PYZdq{}}\PY{p}{)}

\PY{n+nb}{print}\PY{p}{(}\PY{l+s+s2}{\PYZdq{}}\PY{l+s+s2}{Model saved successfully!}\PY{l+s+s2}{\PYZdq{}}\PY{p}{)}
\end{Verbatim}
\end{tcolorbox}

    \begin{Verbatim}[commandchars=\\\{\}]
Model saved successfully!
    \end{Verbatim}

    \begin{tcolorbox}[breakable, size=fbox, boxrule=1pt, pad at break*=1mm,colback=cellbackground, colframe=cellborder]
\prompt{In}{incolor}{144}{\boxspacing}
\begin{Verbatim}[commandchars=\\\{\}]
\PY{n}{pip} \PY{n}{install} \PY{n}{plotly}
\end{Verbatim}
\end{tcolorbox}

    \begin{Verbatim}[commandchars=\\\{\}]
Requirement already satisfied: plotly in
.\textbackslash{}AppData\textbackslash{}Local\textbackslash{}Programs\textbackslash{}Python\textbackslash{}Python314\textbackslash{}Lib\textbackslash{}site-packages (6.9.0)
Requirement already satisfied: narwhals>=1.15.1 in
.\textbackslash{}AppData\textbackslash{}Local\textbackslash{}Programs\textbackslash{}Python\textbackslash{}Python314\textbackslash{}Lib\textbackslash{}site-packages (from plotly)
(2.24.0)
Requirement already satisfied: packaging in
.\textbackslash{}AppData\textbackslash{}Local\textbackslash{}Programs\textbackslash{}Python\textbackslash{}Python314\textbackslash{}Lib\textbackslash{}site-packages (from plotly) (26.2)
Note: you may need to restart the kernel to use updated packages.
    \end{Verbatim}

    \begin{tcolorbox}[breakable, size=fbox, boxrule=1pt, pad at break*=1mm,colback=cellbackground, colframe=cellborder]
\prompt{In}{incolor}{145}{\boxspacing}
\begin{Verbatim}[commandchars=\\\{\}]
\PY{l+s+sd}{\PYZsq{}\PYZsq{}\PYZsq{}import plotly.express as px}

\PY{l+s+sd}{fig = px.histogram(}
\PY{l+s+sd}{    df,}
\PY{l+s+sd}{    x=\PYZdq{}median\PYZus{}house\PYZus{}value\PYZdq{},}
\PY{l+s+sd}{    title=\PYZdq{}Distribution of House Prices\PYZdq{}}
\PY{l+s+sd}{)}

\PY{l+s+sd}{fig.show()\PYZsq{}\PYZsq{}\PYZsq{}}
\end{Verbatim}
\end{tcolorbox}

            \begin{tcolorbox}[breakable, size=fbox, boxrule=.5pt, pad at break*=1mm, opacityfill=0]
\prompt{Out}{outcolor}{145}{\boxspacing}
\begin{Verbatim}[commandchars=\\\{\}]
'import plotly.express as px\textbackslash{}n\textbackslash{}nfig = px.histogram(\textbackslash{}n    df,\textbackslash{}n
x="median\_house\_value",\textbackslash{}n    title="Distribution of House
Prices"\textbackslash{}n)\textbackslash{}n\textbackslash{}nfig.show()'
\end{Verbatim}
\end{tcolorbox}
        
    \begin{tcolorbox}[breakable, size=fbox, boxrule=1pt, pad at break*=1mm,colback=cellbackground, colframe=cellborder]
\prompt{In}{incolor}{146}{\boxspacing}
\begin{Verbatim}[commandchars=\\\{\}]
\PY{o}{!}pip\PY{+w}{ }install\PY{+w}{ }statsmodels
\end{Verbatim}
\end{tcolorbox}

    \begin{Verbatim}[commandchars=\\\{\}]
Requirement already satisfied: statsmodels in
.\textbackslash{}AppData\textbackslash{}Local\textbackslash{}Programs\textbackslash{}Python\textbackslash{}Python314\textbackslash{}Lib\textbackslash{}site-packages (0.14.6)
Requirement already satisfied: numpy<3,>=1.22.3 in
.\textbackslash{}AppData\textbackslash{}Local\textbackslash{}Programs\textbackslash{}Python\textbackslash{}Python314\textbackslash{}Lib\textbackslash{}site-packages (from statsmodels)
(2.5.1)
Requirement already satisfied: scipy!=1.9.2,>=1.8 in
.\textbackslash{}AppData\textbackslash{}Local\textbackslash{}Programs\textbackslash{}Python\textbackslash{}Python314\textbackslash{}Lib\textbackslash{}site-packages (from statsmodels)
(1.18.0)
Requirement already satisfied: pandas!=2.1.0,>=1.4 in
.\textbackslash{}AppData\textbackslash{}Local\textbackslash{}Programs\textbackslash{}Python\textbackslash{}Python314\textbackslash{}Lib\textbackslash{}site-packages (from statsmodels)
(3.0.3)
Requirement already satisfied: patsy>=0.5.6 in
.\textbackslash{}AppData\textbackslash{}Local\textbackslash{}Programs\textbackslash{}Python\textbackslash{}Python314\textbackslash{}Lib\textbackslash{}site-packages (from statsmodels)
(1.0.2)
Requirement already satisfied: packaging>=21.3 in
.\textbackslash{}AppData\textbackslash{}Local\textbackslash{}Programs\textbackslash{}Python\textbackslash{}Python314\textbackslash{}Lib\textbackslash{}site-packages (from statsmodels)
(26.2)
Requirement already satisfied: python-dateutil>=2.8.2 in
.\textbackslash{}AppData\textbackslash{}Local\textbackslash{}Programs\textbackslash{}Python\textbackslash{}Python314\textbackslash{}Lib\textbackslash{}site-packages (from
pandas!=2.1.0,>=1.4->statsmodels) (2.9.0.post0)
Requirement already satisfied: tzdata in
.\textbackslash{}AppData\textbackslash{}Local\textbackslash{}Programs\textbackslash{}Python\textbackslash{}Python314\textbackslash{}Lib\textbackslash{}site-packages (from
pandas!=2.1.0,>=1.4->statsmodels) (2026.3)
Requirement already satisfied: six>=1.5 in
.\textbackslash{}AppData\textbackslash{}Local\textbackslash{}Programs\textbackslash{}Python\textbackslash{}Python314\textbackslash{}Lib\textbackslash{}site-packages (from python-
dateutil>=2.8.2->pandas!=2.1.0,>=1.4->statsmodels) (1.17.0)
    \end{Verbatim}

    \begin{tcolorbox}[breakable, size=fbox, boxrule=1pt, pad at break*=1mm,colback=cellbackground, colframe=cellborder]
\prompt{In}{incolor}{147}{\boxspacing}
\begin{Verbatim}[commandchars=\\\{\}]
\PY{k+kn}{import}\PY{+w}{ }\PY{n+nn}{plotly}\PY{n+nn}{.}\PY{n+nn}{express}\PY{+w}{ }\PY{k}{as}\PY{+w}{ }\PY{n+nn}{px}

\PY{n}{fig} \PY{o}{=} \PY{n}{px}\PY{o}{.}\PY{n}{scatter}\PY{p}{(}
    \PY{n}{df}\PY{p}{,}
    \PY{n}{x}\PY{o}{=}\PY{l+s+s2}{\PYZdq{}}\PY{l+s+s2}{median\PYZus{}income}\PY{l+s+s2}{\PYZdq{}}\PY{p}{,}
    \PY{n}{y}\PY{o}{=}\PY{l+s+s2}{\PYZdq{}}\PY{l+s+s2}{median\PYZus{}house\PYZus{}value}\PY{l+s+s2}{\PYZdq{}}\PY{p}{,}
    \PY{n}{trendline}\PY{o}{=}\PY{l+s+s2}{\PYZdq{}}\PY{l+s+s2}{ols}\PY{l+s+s2}{\PYZdq{}}\PY{p}{,}
    \PY{n}{title}\PY{o}{=}\PY{l+s+s2}{\PYZdq{}}\PY{l+s+s2}{Median Income vs Median House Value}\PY{l+s+s2}{\PYZdq{}}
\PY{p}{)}

\PY{n}{fig}\PY{o}{.}\PY{n}{show}\PY{p}{(}\PY{p}{)}
\end{Verbatim}
\end{tcolorbox}

    \begin{center}
    \adjustimage{max size={0.9\linewidth}{0.9\paperheight}}{output_140_0.png}
    \end{center}
    { \hspace*{\fill} \\}
    
    \begin{tcolorbox}[breakable, size=fbox, boxrule=1pt, pad at break*=1mm,colback=cellbackground, colframe=cellborder]
\prompt{In}{incolor}{148}{\boxspacing}
\begin{Verbatim}[commandchars=\\\{\}]
\PY{k+kn}{import}\PY{+w}{ }\PY{n+nn}{plotly}
\PY{n+nb}{print}\PY{p}{(}\PY{n}{plotly}\PY{o}{.}\PY{n}{\PYZus{}\PYZus{}version\PYZus{}\PYZus{}}\PY{p}{)}
\end{Verbatim}
\end{tcolorbox}

    \begin{Verbatim}[commandchars=\\\{\}]
6.9.0
    \end{Verbatim}

    \begin{tcolorbox}[breakable, size=fbox, boxrule=1pt, pad at break*=1mm,colback=cellbackground, colframe=cellborder]
\prompt{In}{incolor}{149}{\boxspacing}
\begin{Verbatim}[commandchars=\\\{\}]
\PY{k+kn}{import}\PY{+w}{ }\PY{n+nn}{plotly}\PY{n+nn}{.}\PY{n+nn}{express}\PY{+w}{ }\PY{k}{as}\PY{+w}{ }\PY{n+nn}{px}
\end{Verbatim}
\end{tcolorbox}

    \begin{tcolorbox}[breakable, size=fbox, boxrule=1pt, pad at break*=1mm,colback=cellbackground, colframe=cellborder]
\prompt{In}{incolor}{150}{\boxspacing}
\begin{Verbatim}[commandchars=\\\{\}]
\PY{k+kn}{import}\PY{+w}{ }\PY{n+nn}{plotly}\PY{n+nn}{.}\PY{n+nn}{express}\PY{+w}{ }\PY{k}{as}\PY{+w}{ }\PY{n+nn}{px}

\PY{n}{fig} \PY{o}{=} \PY{n}{px}\PY{o}{.}\PY{n}{scatter}\PY{p}{(}
    \PY{n}{x}\PY{o}{=}\PY{p}{[}\PY{l+m+mi}{1}\PY{p}{,} \PY{l+m+mi}{2}\PY{p}{,} \PY{l+m+mi}{3}\PY{p}{,} \PY{l+m+mi}{4}\PY{p}{,} \PY{l+m+mi}{5}\PY{p}{]}\PY{p}{,}
    \PY{n}{y}\PY{o}{=}\PY{p}{[}\PY{l+m+mi}{2}\PY{p}{,} \PY{l+m+mi}{4}\PY{p}{,} \PY{l+m+mi}{6}\PY{p}{,} \PY{l+m+mi}{8}\PY{p}{,} \PY{l+m+mi}{10}\PY{p}{]}\PY{p}{,}
    \PY{n}{title}\PY{o}{=}\PY{l+s+s2}{\PYZdq{}}\PY{l+s+s2}{Test Scatter Plot}\PY{l+s+s2}{\PYZdq{}}
\PY{p}{)}

\PY{n}{fig}\PY{o}{.}\PY{n}{show}\PY{p}{(}\PY{p}{)}
\end{Verbatim}
\end{tcolorbox}

    \begin{center}
    \adjustimage{max size={0.9\linewidth}{0.9\paperheight}}{output_143_0.png}
    \end{center}
    { \hspace*{\fill} \\}
    
    \begin{tcolorbox}[breakable, size=fbox, boxrule=1pt, pad at break*=1mm,colback=cellbackground, colframe=cellborder]
\prompt{In}{incolor}{151}{\boxspacing}
\begin{Verbatim}[commandchars=\\\{\}]
\PY{k+kn}{import}\PY{+w}{ }\PY{n+nn}{plotly}\PY{n+nn}{.}\PY{n+nn}{io}\PY{+w}{ }\PY{k}{as}\PY{+w}{ }\PY{n+nn}{pio}
\PY{n}{pio}\PY{o}{.}\PY{n}{renderers}\PY{o}{.}\PY{n}{default} \PY{o}{=} \PY{l+s+s2}{\PYZdq{}}\PY{l+s+s2}{browser}\PY{l+s+s2}{\PYZdq{}}
\end{Verbatim}
\end{tcolorbox}

    \begin{tcolorbox}[breakable, size=fbox, boxrule=1pt, pad at break*=1mm,colback=cellbackground, colframe=cellborder]
\prompt{In}{incolor}{152}{\boxspacing}
\begin{Verbatim}[commandchars=\\\{\}]
\PY{k+kn}{import}\PY{+w}{ }\PY{n+nn}{plotly}\PY{n+nn}{.}\PY{n+nn}{express}\PY{+w}{ }\PY{k}{as}\PY{+w}{ }\PY{n+nn}{px}

\PY{n}{fig} \PY{o}{=} \PY{n}{px}\PY{o}{.}\PY{n}{scatter}\PY{p}{(}
    \PY{n}{df}\PY{p}{,}
    \PY{n}{x}\PY{o}{=}\PY{l+s+s2}{\PYZdq{}}\PY{l+s+s2}{median\PYZus{}income}\PY{l+s+s2}{\PYZdq{}}\PY{p}{,}
    \PY{n}{y}\PY{o}{=}\PY{l+s+s2}{\PYZdq{}}\PY{l+s+s2}{median\PYZus{}house\PYZus{}value}\PY{l+s+s2}{\PYZdq{}}\PY{p}{,}
    \PY{n}{title}\PY{o}{=}\PY{l+s+s2}{\PYZdq{}}\PY{l+s+s2}{Median Income vs Median House Value}\PY{l+s+s2}{\PYZdq{}}
\PY{p}{)}

\PY{n}{fig}\PY{o}{.}\PY{n}{show}\PY{p}{(}\PY{p}{)}
\end{Verbatim}
\end{tcolorbox}

    \begin{tcolorbox}[breakable, size=fbox, boxrule=1pt, pad at break*=1mm,colback=cellbackground, colframe=cellborder]
\prompt{In}{incolor}{153}{\boxspacing}
\begin{Verbatim}[commandchars=\\\{\}]
\PY{k+kn}{import}\PY{+w}{ }\PY{n+nn}{plotly}\PY{n+nn}{.}\PY{n+nn}{express}\PY{+w}{ }\PY{k}{as}\PY{+w}{ }\PY{n+nn}{px}

\PY{n}{fig} \PY{o}{=} \PY{n}{px}\PY{o}{.}\PY{n}{histogram}\PY{p}{(}
    \PY{n}{df}\PY{p}{,}
    \PY{n}{x}\PY{o}{=}\PY{l+s+s2}{\PYZdq{}}\PY{l+s+s2}{median\PYZus{}house\PYZus{}value}\PY{l+s+s2}{\PYZdq{}}\PY{p}{,}
    \PY{n}{title}\PY{o}{=}\PY{l+s+s2}{\PYZdq{}}\PY{l+s+s2}{Distribution of House Prices}\PY{l+s+s2}{\PYZdq{}}
\PY{p}{)}

\PY{n}{fig}\PY{o}{.}\PY{n}{show}\PY{p}{(}\PY{p}{)}
\end{Verbatim}
\end{tcolorbox}

    \begin{tcolorbox}[breakable, size=fbox, boxrule=1pt, pad at break*=1mm,colback=cellbackground, colframe=cellborder]
\prompt{In}{incolor}{154}{\boxspacing}
\begin{Verbatim}[commandchars=\\\{\}]
\PY{n}{fig} \PY{o}{=} \PY{n}{px}\PY{o}{.}\PY{n}{box}\PY{p}{(}
    \PY{n}{df}\PY{p}{,}
    \PY{n}{y}\PY{o}{=}\PY{l+s+s2}{\PYZdq{}}\PY{l+s+s2}{median\PYZus{}income}\PY{l+s+s2}{\PYZdq{}}\PY{p}{,}
    \PY{n}{title}\PY{o}{=}\PY{l+s+s2}{\PYZdq{}}\PY{l+s+s2}{Median Income Box Plot}\PY{l+s+s2}{\PYZdq{}}
\PY{p}{)}

\PY{n}{fig}\PY{o}{.}\PY{n}{show}\PY{p}{(}\PY{p}{)}
\end{Verbatim}
\end{tcolorbox}

    \begin{tcolorbox}[breakable, size=fbox, boxrule=1pt, pad at break*=1mm,colback=cellbackground, colframe=cellborder]
\prompt{In}{incolor}{155}{\boxspacing}
\begin{Verbatim}[commandchars=\\\{\}]
\PY{o}{!}pip\PY{+w}{ }install\PY{+w}{ }streamlit
\end{Verbatim}
\end{tcolorbox}

    \begin{Verbatim}[commandchars=\\\{\}]
Requirement already satisfied: streamlit in
.\textbackslash{}AppData\textbackslash{}Local\textbackslash{}Programs\textbackslash{}Python\textbackslash{}Python314\textbackslash{}Lib\textbackslash{}site-packages (1.59.2)
Requirement already satisfied: altair!=5.4.0,!=5.4.1,<7,>=4.0 in
.\textbackslash{}AppData\textbackslash{}Local\textbackslash{}Programs\textbackslash{}Python\textbackslash{}Python314\textbackslash{}Lib\textbackslash{}site-packages (from streamlit)
(6.2.2)
Requirement already satisfied: blinker<2,>=1.5.0 in
.\textbackslash{}AppData\textbackslash{}Local\textbackslash{}Programs\textbackslash{}Python\textbackslash{}Python314\textbackslash{}Lib\textbackslash{}site-packages (from streamlit)
(1.9.0)
Requirement already satisfied: cachetools<8,>=5.5 in
.\textbackslash{}AppData\textbackslash{}Local\textbackslash{}Programs\textbackslash{}Python\textbackslash{}Python314\textbackslash{}Lib\textbackslash{}site-packages (from streamlit)
(7.1.4)
Requirement already satisfied: click<9,>=7.0 in
.\textbackslash{}AppData\textbackslash{}Local\textbackslash{}Programs\textbackslash{}Python\textbackslash{}Python314\textbackslash{}Lib\textbackslash{}site-packages (from streamlit)
(8.4.2)
Requirement already satisfied: gitpython!=3.1.19,<4,>=3.0.7 in
.\textbackslash{}AppData\textbackslash{}Local\textbackslash{}Programs\textbackslash{}Python\textbackslash{}Python314\textbackslash{}Lib\textbackslash{}site-packages (from streamlit)
(3.1.52)
Requirement already satisfied: numpy<3,>=1.23 in
.\textbackslash{}AppData\textbackslash{}Local\textbackslash{}Programs\textbackslash{}Python\textbackslash{}Python314\textbackslash{}Lib\textbackslash{}site-packages (from streamlit)
(2.5.1)
Requirement already satisfied: packaging>=20 in
.\textbackslash{}AppData\textbackslash{}Local\textbackslash{}Programs\textbackslash{}Python\textbackslash{}Python314\textbackslash{}Lib\textbackslash{}site-packages (from streamlit)
(26.2)
Requirement already satisfied: pandas<4,>=1.4.0 in
.\textbackslash{}AppData\textbackslash{}Local\textbackslash{}Programs\textbackslash{}Python\textbackslash{}Python314\textbackslash{}Lib\textbackslash{}site-packages (from streamlit)
(3.0.3)
Requirement already satisfied: pillow<13,>=7.1.0 in
.\textbackslash{}AppData\textbackslash{}Local\textbackslash{}Programs\textbackslash{}Python\textbackslash{}Python314\textbackslash{}Lib\textbackslash{}site-packages (from streamlit)
(12.3.0)
Requirement already satisfied: pydeck<1,>=0.8.0b4 in
.\textbackslash{}AppData\textbackslash{}Local\textbackslash{}Programs\textbackslash{}Python\textbackslash{}Python314\textbackslash{}Lib\textbackslash{}site-packages (from streamlit)
(0.9.3)
Requirement already satisfied: protobuf<8,>=3.20 in
.\textbackslash{}AppData\textbackslash{}Local\textbackslash{}Programs\textbackslash{}Python\textbackslash{}Python314\textbackslash{}Lib\textbackslash{}site-packages (from streamlit)
(7.35.1)
Requirement already satisfied: pyarrow<25,>=7.0 in
.\textbackslash{}AppData\textbackslash{}Local\textbackslash{}Programs\textbackslash{}Python\textbackslash{}Python314\textbackslash{}Lib\textbackslash{}site-packages (from streamlit)
(24.0.0)
Requirement already satisfied: requests<3,>=2.27 in
.\textbackslash{}AppData\textbackslash{}Local\textbackslash{}Programs\textbackslash{}Python\textbackslash{}Python314\textbackslash{}Lib\textbackslash{}site-packages (from streamlit)
(2.34.2)
Requirement already satisfied: tenacity<10,>=8.1.0 in
.\textbackslash{}AppData\textbackslash{}Local\textbackslash{}Programs\textbackslash{}Python\textbackslash{}Python314\textbackslash{}Lib\textbackslash{}site-packages (from streamlit)
(9.1.4)
Requirement already satisfied: toml<2,>=0.10.1 in
.\textbackslash{}AppData\textbackslash{}Local\textbackslash{}Programs\textbackslash{}Python\textbackslash{}Python314\textbackslash{}Lib\textbackslash{}site-packages (from streamlit)
(0.10.2)
Requirement already satisfied: typing-extensions<5,>=4.10.0 in
.\textbackslash{}AppData\textbackslash{}Local\textbackslash{}Programs\textbackslash{}Python\textbackslash{}Python314\textbackslash{}Lib\textbackslash{}site-packages (from streamlit)
(4.16.0)
Requirement already satisfied: starlette>=0.40.0 in
.\textbackslash{}AppData\textbackslash{}Local\textbackslash{}Programs\textbackslash{}Python\textbackslash{}Python314\textbackslash{}Lib\textbackslash{}site-packages (from streamlit)
(1.3.1)
Requirement already satisfied: uvicorn>=0.30.0 in
.\textbackslash{}AppData\textbackslash{}Local\textbackslash{}Programs\textbackslash{}Python\textbackslash{}Python314\textbackslash{}Lib\textbackslash{}site-packages (from streamlit)
(0.51.0)
Requirement already satisfied: httptools>=0.6.3 in
.\textbackslash{}AppData\textbackslash{}Local\textbackslash{}Programs\textbackslash{}Python\textbackslash{}Python314\textbackslash{}Lib\textbackslash{}site-packages (from streamlit)
(0.8.0)
Requirement already satisfied: anyio>=4.0.0 in
.\textbackslash{}AppData\textbackslash{}Local\textbackslash{}Programs\textbackslash{}Python\textbackslash{}Python314\textbackslash{}Lib\textbackslash{}site-packages (from streamlit)
(4.14.2)
Requirement already satisfied: python-multipart>=0.0.10 in
.\textbackslash{}AppData\textbackslash{}Local\textbackslash{}Programs\textbackslash{}Python\textbackslash{}Python314\textbackslash{}Lib\textbackslash{}site-packages (from streamlit)
(0.0.32)
Requirement already satisfied: websockets>=12.0.0 in
.\textbackslash{}AppData\textbackslash{}Local\textbackslash{}Programs\textbackslash{}Python\textbackslash{}Python314\textbackslash{}Lib\textbackslash{}site-packages (from streamlit)
(16.1.1)
Requirement already satisfied: itsdangerous>=2.1.2 in
.\textbackslash{}AppData\textbackslash{}Local\textbackslash{}Programs\textbackslash{}Python\textbackslash{}Python314\textbackslash{}Lib\textbackslash{}site-packages (from streamlit)
(2.2.0)
Requirement already satisfied: watchdog<7,>=2.1.5 in
.\textbackslash{}AppData\textbackslash{}Local\textbackslash{}Programs\textbackslash{}Python\textbackslash{}Python314\textbackslash{}Lib\textbackslash{}site-packages (from streamlit)
(6.0.0)
Requirement already satisfied: jinja2 in
.\textbackslash{}AppData\textbackslash{}Local\textbackslash{}Programs\textbackslash{}Python\textbackslash{}Python314\textbackslash{}Lib\textbackslash{}site-packages (from
altair!=5.4.0,!=5.4.1,<7,>=4.0->streamlit) (3.1.6)
Requirement already satisfied: jsonschema>=3.0 in
.\textbackslash{}AppData\textbackslash{}Local\textbackslash{}Programs\textbackslash{}Python\textbackslash{}Python314\textbackslash{}Lib\textbackslash{}site-packages (from
altair!=5.4.0,!=5.4.1,<7,>=4.0->streamlit) (4.26.0)
Requirement already satisfied: narwhals>=2.4.0 in
.\textbackslash{}AppData\textbackslash{}Local\textbackslash{}Programs\textbackslash{}Python\textbackslash{}Python314\textbackslash{}Lib\textbackslash{}site-packages (from
altair!=5.4.0,!=5.4.1,<7,>=4.0->streamlit) (2.24.0)
Requirement already satisfied: colorama in
.\textbackslash{}AppData\textbackslash{}Local\textbackslash{}Programs\textbackslash{}Python\textbackslash{}Python314\textbackslash{}Lib\textbackslash{}site-packages (from
click<9,>=7.0->streamlit) (0.4.6)
Requirement already satisfied: gitdb<5,>=4.0.1 in
.\textbackslash{}AppData\textbackslash{}Local\textbackslash{}Programs\textbackslash{}Python\textbackslash{}Python314\textbackslash{}Lib\textbackslash{}site-packages (from
gitpython!=3.1.19,<4,>=3.0.7->streamlit) (4.0.12)
Requirement already satisfied: smmap<6,>=3.0.1 in
.\textbackslash{}AppData\textbackslash{}Local\textbackslash{}Programs\textbackslash{}Python\textbackslash{}Python314\textbackslash{}Lib\textbackslash{}site-packages (from
gitdb<5,>=4.0.1->gitpython!=3.1.19,<4,>=3.0.7->streamlit) (5.0.3)
Requirement already satisfied: python-dateutil>=2.8.2 in
.\textbackslash{}AppData\textbackslash{}Local\textbackslash{}Programs\textbackslash{}Python\textbackslash{}Python314\textbackslash{}Lib\textbackslash{}site-packages (from
pandas<4,>=1.4.0->streamlit) (2.9.0.post0)
Requirement already satisfied: tzdata in
.\textbackslash{}AppData\textbackslash{}Local\textbackslash{}Programs\textbackslash{}Python\textbackslash{}Python314\textbackslash{}Lib\textbackslash{}site-packages (from
pandas<4,>=1.4.0->streamlit) (2026.3)
Requirement already satisfied: charset\_normalizer<4,>=2 in
.\textbackslash{}AppData\textbackslash{}Local\textbackslash{}Programs\textbackslash{}Python\textbackslash{}Python314\textbackslash{}Lib\textbackslash{}site-packages (from
requests<3,>=2.27->streamlit) (3.4.9)
Requirement already satisfied: idna<4,>=2.5 in
.\textbackslash{}AppData\textbackslash{}Local\textbackslash{}Programs\textbackslash{}Python\textbackslash{}Python314\textbackslash{}Lib\textbackslash{}site-packages (from
requests<3,>=2.27->streamlit) (3.18)
Requirement already satisfied: urllib3<3,>=1.26 in
.\textbackslash{}AppData\textbackslash{}Local\textbackslash{}Programs\textbackslash{}Python\textbackslash{}Python314\textbackslash{}Lib\textbackslash{}site-packages (from
requests<3,>=2.27->streamlit) (2.7.0)
Requirement already satisfied: certifi>=2023.5.7 in
.\textbackslash{}AppData\textbackslash{}Local\textbackslash{}Programs\textbackslash{}Python\textbackslash{}Python314\textbackslash{}Lib\textbackslash{}site-packages (from
requests<3,>=2.27->streamlit) (2026.6.17)
Requirement already satisfied: MarkupSafe>=2.0 in
.\textbackslash{}AppData\textbackslash{}Local\textbackslash{}Programs\textbackslash{}Python\textbackslash{}Python314\textbackslash{}Lib\textbackslash{}site-packages (from
jinja2->altair!=5.4.0,!=5.4.1,<7,>=4.0->streamlit) (3.0.3)
Requirement already satisfied: attrs>=22.2.0 in
.\textbackslash{}AppData\textbackslash{}Local\textbackslash{}Programs\textbackslash{}Python\textbackslash{}Python314\textbackslash{}Lib\textbackslash{}site-packages (from
jsonschema>=3.0->altair!=5.4.0,!=5.4.1,<7,>=4.0->streamlit) (26.1.0)
Requirement already satisfied: jsonschema-specifications>=2023.03.6 in
.\textbackslash{}AppData\textbackslash{}Local\textbackslash{}Programs\textbackslash{}Python\textbackslash{}Python314\textbackslash{}Lib\textbackslash{}site-packages (from
jsonschema>=3.0->altair!=5.4.0,!=5.4.1,<7,>=4.0->streamlit) (2025.9.1)
Requirement already satisfied: referencing>=0.28.4 in
.\textbackslash{}AppData\textbackslash{}Local\textbackslash{}Programs\textbackslash{}Python\textbackslash{}Python314\textbackslash{}Lib\textbackslash{}site-packages (from
jsonschema>=3.0->altair!=5.4.0,!=5.4.1,<7,>=4.0->streamlit) (0.37.0)
Requirement already satisfied: rpds-py>=0.25.0 in
.\textbackslash{}AppData\textbackslash{}Local\textbackslash{}Programs\textbackslash{}Python\textbackslash{}Python314\textbackslash{}Lib\textbackslash{}site-packages (from
jsonschema>=3.0->altair!=5.4.0,!=5.4.1,<7,>=4.0->streamlit) (2026.6.3)
Requirement already satisfied: six>=1.5 in
.\textbackslash{}AppData\textbackslash{}Local\textbackslash{}Programs\textbackslash{}Python\textbackslash{}Python314\textbackslash{}Lib\textbackslash{}site-packages (from python-
dateutil>=2.8.2->pandas<4,>=1.4.0->streamlit) (1.17.0)
Requirement already satisfied: h11>=0.8 in
.\textbackslash{}AppData\textbackslash{}Local\textbackslash{}Programs\textbackslash{}Python\textbackslash{}Python314\textbackslash{}Lib\textbackslash{}site-packages (from
uvicorn>=0.30.0->streamlit) (0.16.0)
    \end{Verbatim}

    \begin{tcolorbox}[breakable, size=fbox, boxrule=1pt, pad at break*=1mm,colback=cellbackground, colframe=cellborder]
\prompt{In}{incolor}{156}{\boxspacing}
\begin{Verbatim}[commandchars=\\\{\}]
\PY{n}{model}\PY{o}{.}\PY{n}{fit}\PY{p}{(}\PY{n}{X\PYZus{}train}\PY{p}{,} \PY{n}{y\PYZus{}train}\PY{p}{)}
\end{Verbatim}
\end{tcolorbox}

            \begin{tcolorbox}[breakable, size=fbox, boxrule=.5pt, pad at break*=1mm, opacityfill=0]
\prompt{Out}{outcolor}{156}{\boxspacing}
\begin{Verbatim}[commandchars=\\\{\}]
LinearRegression()
\end{Verbatim}
\end{tcolorbox}
        
    \begin{tcolorbox}[breakable, size=fbox, boxrule=1pt, pad at break*=1mm,colback=cellbackground, colframe=cellborder]
\prompt{In}{incolor}{157}{\boxspacing}
\begin{Verbatim}[commandchars=\\\{\}]
\PY{k+kn}{import}\PY{+w}{ }\PY{n+nn}{joblib}

\PY{n}{joblib}\PY{o}{.}\PY{n}{dump}\PY{p}{(}\PY{n}{model}\PY{p}{,} \PY{l+s+s2}{\PYZdq{}}\PY{l+s+s2}{house\PYZus{}model.pkl}\PY{l+s+s2}{\PYZdq{}}\PY{p}{)}

\PY{n+nb}{print}\PY{p}{(}\PY{l+s+s2}{\PYZdq{}}\PY{l+s+s2}{Model saved successfully!}\PY{l+s+s2}{\PYZdq{}}\PY{p}{)}
\end{Verbatim}
\end{tcolorbox}

    \begin{Verbatim}[commandchars=\\\{\}]
Model saved successfully!
    \end{Verbatim}

    \begin{tcolorbox}[breakable, size=fbox, boxrule=1pt, pad at break*=1mm,colback=cellbackground, colframe=cellborder]
\prompt{In}{incolor}{158}{\boxspacing}
\begin{Verbatim}[commandchars=\\\{\}]
\PY{k+kn}{import}\PY{+w}{ }\PY{n+nn}{os}

\PY{n+nb}{print}\PY{p}{(}\PY{n}{os}\PY{o}{.}\PY{n}{getcwd}\PY{p}{(}\PY{p}{)}\PY{p}{)}
\PY{n+nb}{print}\PY{p}{(}\PY{n}{os}\PY{o}{.}\PY{n}{listdir}\PY{p}{(}\PY{p}{)}\PY{p}{)}
\end{Verbatim}
\end{tcolorbox}

    \begin{Verbatim}[commandchars=\\\{\}]
C:\textbackslash{}Users\textbackslash{}deekshitha.kotian
['.bash\_history', '.gitconfig', '.idlerc', '.ipynb\_checkpoints', '.ipython',
'.jupyter', '.lesshst', '.python\_history', '.streamlit', '.vscode', '.vscode-
shared', 'app.py', 'AppData', 'Application Data', 'Contacts', 'Cookies',
'Documents', 'Downloads', 'Favorites', 'house\_model.pkl', 'housing.csv',
'Links', 'Local Settings', 'Music', 'My Documents', 'NetHood', 'NTUSER.DAT',
'ntuser.dat.LOG1', 'ntuser.dat.LOG2', 'NTUSER.DAT\{c4d96ec7-771b-11f1-85ca-
ac91a10ae34c\}.TM.blf', 'NTUSER.DAT\{c4d96ec7-771b-11f1-85ca-
ac91a10ae34c\}.TMContainer00000000000000000001.regtrans-ms',
'NTUSER.DAT\{c4d96ec7-771b-11f1-85ca-
ac91a10ae34c\}.TMContainer00000000000000000002.regtrans-ms', 'ntuser.ini',
'OneDrive', 'OneDrive - VCTI', 'PrintHood', 'Python-Training', 'Recent', 'Saved
Games', 'Searches', 'SendTo', 'source', 'Start Menu', 'Templates',
'Untitled-1.py', 'Untitled.ipynb', 'Untitled1.ipynb', 'Untitled1.PDF', 'Videos']
    \end{Verbatim}

    \begin{tcolorbox}[breakable, size=fbox, boxrule=1pt, pad at break*=1mm,colback=cellbackground, colframe=cellborder]
\prompt{In}{incolor}{110}{\boxspacing}
\begin{Verbatim}[commandchars=\\\{\}]
\PY{k+kn}{import}\PY{+w}{ }\PY{n+nn}{joblib}

\PY{n}{model} \PY{o}{=} \PY{n}{joblib}\PY{o}{.}\PY{n}{load}\PY{p}{(}\PY{l+s+s2}{\PYZdq{}}\PY{l+s+s2}{house\PYZus{}model.pkl}\PY{l+s+s2}{\PYZdq{}}\PY{p}{)}
\end{Verbatim}
\end{tcolorbox}

    \begin{tcolorbox}[breakable, size=fbox, boxrule=1pt, pad at break*=1mm,colback=cellbackground, colframe=cellborder]
\prompt{In}{incolor}{ }{\boxspacing}
\begin{Verbatim}[commandchars=\\\{\}]

\end{Verbatim}
\end{tcolorbox}


    % Add a bibliography block to the postdoc
    
    
    
\end{document}
